\section{Sufficient statistics for the initial state}\label{m:sec:suff}

Chapter~\ref{ch:gbc} shows that Bayes' rule can be learned as a map from data to posterior using simulated pairs, and Chapter~\ref{ch:manifold} applied the idea to the initial condition of an evolution equation (Algorithm~\ref{m:alg:gbc}). The method can be no better than the summary of the data it conditions on: a statistic that discards information about the initial condition limits every generator trained on it, whatever the simulation budget, while the Bernstein von Mises theorem of Nickl~\cite{nicklbvm2024} (Theorem~\ref{m:thm:nickl2}) describes the posterior given all of the data. This chapter asks what to condition on and what is lost by conditioning on something else. For a linear equation the answer is exact: the data run backward through the adjoint equation form a sufficient statistic for the initial state, and with a Gaussian prior the map from that statistic to the posterior is affine. For a nonlinear equation no finite dimensional statistic is exactly sufficient, and the reduction is exact only to the order the Bernstein von Mises theorem controls, which is the order the transfer theorem of Chapter~\ref{ch:generators} needs. We construct a one step efficient score summary that agrees with the centering of Nickl's theorem to $o(N^{-1/2})$ under the conditions of Proposition~\ref{m:prop:suff}, and Section~\ref{m:sec:abcpls} places it among the summary statistics of approximate Bayesian computation and shows that, in the linear case, the construction is partial least squares.

The posterior mean is exact for the first moment at every sample size and, under the uniform Bernstein von Mises statement of Assumption~\ref{m:ass:U}, sufficient for the whole posterior at the scale $N^{-1/2}$ (Section~\ref{m:sec:jw}). A summary that is not sufficient leaves a gap in Wasserstein distance, which Section~\ref{m:sec:suffbounds} brackets between two quantities computable from the model. Under dissipation the effective dimension of the adjoint statistic falls with the viscosity, and a generator required to be uniformly accurate on the raw data needs a super exponential simulation budget, whereas on the statistic the budget is linear in the effective dimension. The results for linear equations are unconditional; in the nonlinear case everything that concerns the posterior given the summary rests on Assumption~\ref{m:ass:U}, left as Open Question~\ref{m:op:unif}, and the budget linear in the effective dimension rests also on Assumption~\ref{m:ass:u2p}; every statement that uses either assumption says so.

\subsection{Linear dissipative equations: the adjoint evolution is exactly sufficient}

Let the evolution be linear, $u_\theta(t)=S_t\theta$ with $S_t$ a strongly continuous semigroup on $L^2_0(\Omega)$ diagonal in the basis $(e_j)$ with eigenvalues $s_j(t)$; for the heat equation $s_j(t)=e^{-\nu\lambda_jt}$, for Stokes the same on divergence free fields. Write $\theta=\sum_j\theta_je_j$ and truncate to $D$ modes, so that $\theta\in\R^D$ and the design matrix is
\begin{equation}\label{m:eq:design}
\Phi_{ij}=s_j(t_i)\,e_j(x_i),\qquad i\le N,\ j\le D,\qquad Y=\Phi\theta+\varepsilon .
\end{equation}
With a Gaussian prior $\theta\sim N(0,\Sigma_0)$ this is a Gaussian linear model, and its posterior is $N(m_N,V_N)$ with
\begin{equation}\label{m:eq:linpost}
V_N^{-1}=\Phi^\top\Phi+\Sigma_0^{-1},\qquad m_N=V_N\,\Phi^\top Y .
\end{equation}
By the factorization theorem, conditional on the design $(X_i)$ the statistic $\Phi^\top Y\in\R^D$ is sufficient for $\theta$, and since the design is ancillary the pair $(\Phi^\top Y,\Phi^\top\Phi)$ is sufficient unconditionally.

The statistic $\Phi^\top Y$ has a PDE meaning, which we record as a lemma because it is the construction that generalizes.

\begin{lemma}[Adjoint representation of the sufficient statistic]\label{m:lem:adjoint}
Let $S_t$ be a strongly continuous semigroup on $L^2(\Omega)$ with $S_te_j=s_j(t)e_j$, and let $r=\sum_{i\le N}Y_i\,\delta_{t_i}\otimes\delta_{x_i}$ be the empirical data measure on $[0,T]\times\Omega$. Then for every $h=\sum_jh_je_j\in\mathrm{span}\{e_j:j\le D\}$,
\[
\big\langle\Phi^\top Y,\,h\big\rangle_{\R^D}=\sum_{i=1}^NY_i\,(S_{t_i}h)(x_i)=\Big\langle h,\ \int_0^TS_t^*\,r(t)\,dt\Big\rangle ,
\]
so that $\Phi^\top Y$ is the coefficient vector of $w(0)$, where $w$ is the mild solution of the backward problem $-\partial_sw=A^*w+r$, $w(T)=0$, with $A$ the generator of $S_t$.
\end{lemma}

\begin{proof}
By definition $(\Phi^\top Y)_j=\sum_iY_is_j(t_i)e_j(x_i)$, so
\[
\langle\Phi^\top Y,h\rangle=\sum_iY_i\sum_jh_js_j(t_i)e_j(x_i)=\sum_iY_i(S_{t_i}h)(x_i).
\] Since $(S_th)(x)=\langle S_th,\delta_x\rangle=\langle h,S_t^*\delta_x\rangle$ in the sense of distributions, the sum equals $\langle h,\sum_iY_iS_{t_i}^*\delta_{x_i}\rangle=\langle h,\int_0^TS_t^*r(t)dt\rangle$. The variation of constants formula for the backward problem gives $w(0)=\int_0^TS^*_tr(t)\,dt$.
\end{proof}

For the heat equation the semigroup is self adjoint and the backward problem of Lemma~\ref{m:lem:adjoint} reads
\begin{equation}\label{m:eq:adjoint}
\begin{aligned}
&-\partial_sw=\nu\Delta w+\sum_{i=1}^NY_i\,\delta_{(t_i,x_i)},\qquad w(T)=0,\\
&\Phi^\top Y=\big(\langle w(0),e_j\rangle\big)_{j\le D},
\end{aligned}
\end{equation}
so the sufficient statistic for the initial state is the data run backward in time through the adjoint of the forward evolution and read off at $t=0$. Each observation at time $t_i$ contributes to mode $j$ with weight $e^{-\nu\lambda_jt_i}$, which is the fraction of mode $j$ that survived to the observation time; observations taken after the mode has decayed carry no information about it, which is the mechanism behind the dissipation barrier of Section~\ref{m:sec:heat}. The Gram matrix $\Phi^\top\Phi$ concentrates, by the law of large numbers over the uniform design, around $N$ times the information operator of Section~\ref{m:sec:papers}:
\begin{equation}\label{m:eq:gram}
\begin{aligned}
\frac1N(\Phi^\top\Phi)_{jj'}&\to\frac1T\int_0^Ts_j(t)s_{j'}(t)\,dt\,\langle e_j,e_{j'}\rangle=\delta_{jj'}\,\iota_j,\\
\iota_j&:=\frac1T\int_0^Ts_j(t)^2dt ,
\end{aligned}
\end{equation}
and for the heat equation $\iota_j=(1-e^{-2\nu\lambda_jT})/(2\nu\lambda_jT)$, which is \eqref{m:eq:modeinfo} divided by $T$. Off diagonal terms are $O_P(N^{-1/2})$ by orthonormality of the $e_j$.

\begin{proposition}[Exact sufficiency, linear case]\label{m:prop:linsuff}
In the linear model \eqref{m:eq:design} with Gaussian prior, the adjoint statistic $S_N(Z)=\Phi^\top Y=w(0)$ from \eqref{m:eq:adjoint}, together with the design, is sufficient for $\theta$; the posterior given $S_N$ is exactly the posterior given $Z^{(N)}$; and the conditional quantile map is exactly affine, with $z_\tau$ the $\tau$ quantile of the standard normal law,
\begin{equation}\label{m:eq:linquant}
Q_j(\tau\mid S_N)=\big(V_NS_N\big)_j+z_\tau\,(V_N)_{jj}^{1/2},\qquad V_N=(\Phi^\top\Phi+\Sigma_0^{-1})^{-1}.
\end{equation}
The slope of that map is $V_N$, not the identity. It is known in closed form from the design alone, so the generator has no slope to learn: writing $\tilde S_N=V_NS_N$ for the preconditioned summary, which is the posterior mean itself, the map $\tilde S_N\mapsto Q_j(\tau\mid S_N)$ has slope exactly one and, apart from the data independent scale term $z_\tau(V_N)_{jj}^{1/2}$, intercept zero. Consequently GBC with summary $\tilde S_N$ and an identity plus scale generator is exact for every $N$, up to the estimation of the $D\times D$ scale from $M$ simulations. In the linear model that scale is $V_N$ itself and is known in closed form, so even the estimation is optional; Remark~\ref{m:rem:emp} discusses the empirical quantile head, which is the one to use at moderate $N$.
\end{proposition}

\begin{proof}
Write the linear model as $Y=\Phi\theta+\varepsilon$ with $\varepsilon\sim N(0,\Id_N)$. Its density factorizes as
\[
\exp\big(-\tfrac12\norm{Y}^2\big)\cdot\exp\big(\theta^\top\Phi^\top Y-\tfrac12\theta^\top\Phi^\top\Phi\,\theta\big),
\]
in which $\theta$ enters only through $\Phi^\top Y$ and the design, so $S_N=\Phi^\top Y$ is sufficient by the factorization theorem, and Lemma~\ref{m:lem:adjoint} identifies it with $w(0)$. Conjugacy with the prior $N(0,\Sigma_0)$ gives $\Law(\theta\mid Z^{(N)})=N(V_N\Phi^\top Y,V_N)$ with $V_N=(\Phi^\top\Phi+\Sigma_0^{-1})^{-1}$, which depends on the data only through $S_N$; it is therefore also $\Law(\theta\mid S_N)$. Its $j$th marginal is $N\big((V_NS_N)_j,(V_N)_{jj}\big)$, whose $\tau$ quantile is \eqref{m:eq:linquant}. As a map of $S_N$ the slope is $V_N$; in the preconditioned parametrization $\tilde S_N=V_NS_N$, \eqref{m:eq:linquant} is $\tilde S_{N,j}+z_\tau(V_N)_{jj}^{1/2}$, an identity skip of slope one and a scale that does not depend on the data, so estimating the scale from $M$ simulations is the only error GBC makes.
\end{proof}

\subsection{Nonlinear equations: the one step efficient score statistic}

For nonlinear $\Gcal$ there is no finite dimensional sufficient statistic. But the Bernstein von Mises theorem says that the posterior is asymptotically $N(\hat\theta_N,T(\I^*_{\theta_0}\I_{\theta_0})^{-1}/N)$ with a data independent covariance, so the centering $\hat\theta_N$ carries the information at the $\sqrt N$ scale (Theorem~\ref{m:thm:jwbvm} makes this precise under Assumption~\ref{m:ass:U}). The centering \eqref{m:eq:thetahat} is not computable, since it involves $\theta_0$ through $\I_{\theta_0}$ and through $\varepsilon_i=Y_i-u_{\theta_0}(X_i)$. Replacing $\theta_0$ by a preliminary estimate is the one step device of semiparametric statistics~\cite{vandervaart1998}.

\begin{definition}[One step efficient score summary]\label{m:def:summary}
Split $Z^{(N)}$ into halves $Z^{[1]},Z^{[2]}$ of sizes differing by at most one, with index sets $\mathcal I_1,\mathcal I_2$. For $b=1,2$ let $\bar\theta^{[b]}_N$ be computed from $Z^{[b]}$ alone, take values in the parameter space and satisfy $\norm{\bar\theta^{[b]}_N-\theta_0}_{H^\xi}=O_P(\tilde\delta_N(\xi))$ for $0\le\xi\le\bar\gamma$, where $\bar\gamma$ is the Sobolev index of the localization ball in the proofs of Nickl~\cite{nicklbvm2024}, chosen as in Section~\ref{m:sec:main}, and $\tilde\delta_N(\xi)=\delta_N^{(\bar\gamma-\xi)/(\bar\gamma+1)}$, so that at $\xi=\bar\gamma$ the estimate is bounded in $H^{\bar\gamma}$ in probability, as Condition 2 of Nickl~\cite{nicklbvm2024} requires of the points at which it is applied; for instance the $t=0$ trace of a nonparametric space time smoother, or the mean of a cheap first pass GBC. Write $\bar e_j^{\,\vartheta}=T(\I^*_\vartheta\I_\vartheta)^{-1}e_j$ and, for $j\le J_N\asymp N\delta_N^2$,
\begin{equation}\label{m:eq:summary}
\begin{aligned}
S_{N,j}(Z)&=\frac12\sum_{b=1,2}\Big[\langle e_j,\bar\theta^{[b]}_N\rangle\\
&\quad+\frac{1}{|\mathcal I_{3-b}|}\sum_{i\in\mathcal I_{3-b}}\Big\langle\I_{\bar\theta^{[b]}_N}\,p_N\big(\bar e_j^{\,\bar\theta^{[b]}_N}\big)(X_i),\;Y_i-u_{\bar\theta^{[b]}_N}(X_i)\Big\rangle_W\Big],
\end{aligned}
\end{equation}
where $S_N(Z)=(S_{N,j})_{j\le J_N}$. The preliminary estimate is formed on one half and the one step correction taken over the complementary half, and the two are averaged: this is cross fitting in the sense of Chernozhukov and coauthors~\cite{chernozhukov2018}.
\end{definition}

Cross fitting rather than a single split, because two things must hold at once. Each summand in \eqref{m:eq:summary} must be independent of the estimate that weights it, or the errors $\varepsilon_i$ it weights also drive $\bar\theta_N$, the average is not conditionally Gaussian given the design, and the step on the stochastic term in the proof of Proposition~\ref{m:prop:suff} fails. And the statistic must carry the information of all $N$ observations, or it is not $\sqrt N$ equivalent to the centering \eqref{m:eq:thetahat}. A single split gives the first and not the second, correcting on one half only and so doubling the efficient conditional variance; averaging over all $N$ without splitting gives the second and not the first. The alternative is an entropy condition on the class of admissible $\bar\theta_N$, which we do not pursue. In the linear model the question does not arise for the regularized form \eqref{m:eq:regsummary} of Remark~\ref{m:rem:reg}, which is what we compute, since Lemma~\ref{m:lem:exact} makes it exact and independent of $\bar\theta_N$ altogether; that is why the choice is settled here by argument and not by the experiments of Section~\ref{m:sec:experiments}.

The second term is again an adjoint evolution driven by data. Fix a fold $b$, write $\bar\theta_N$ for $\bar\theta^{[b]}_N$, let $\mathcal I=\mathcal I_{3-b}$ be the complementary index set and $r_i=Y_i-u_{\bar\theta_N}(X_i)$ the residuals on it; then $\frac1{|\mathcal I|}\sum_{i\in\mathcal I}\langle\I_{\bar\theta_N}h(X_i),r_i\rangle=\langle h,\frac1{|\mathcal I|}\I^*_{\bar\theta_N}r\rangle$, so that fold's contribution to $S_N$ is the $J_N$ dimensional projection of
\begin{equation}\label{m:eq:onestep}
\hat S^{[b]}_N=\bar\theta_N+T(\I^*_{\bar\theta_N}\I_{\bar\theta_N})^{-1}p_N\,\frac1{|\mathcal I|}\I^*_{\bar\theta_N}\Big[\sum_{i\in\mathcal I}r_i\,\delta_{X_i}\Big],
\end{equation}
which is one Newton (Gauss Newton, or Fisher scoring) step from $\bar\theta_N$ on the least squares criterion, with the Fisher information as Hessian and the gradient projected onto the first $J_N$ modes, and we write $\hat S_N=\hat S^{J_N}_N=\sum_{j\le J_N}S_{N,j}e_j$ for the projection of $\frac12(\hat S^{[1]}_N+\hat S^{[2]}_N)$ onto those modes. The computation is: one nonlinear forward solve to get $u_{\bar\theta_N}$ and the residuals; one adjoint solve of the linearized equation driven by the residuals as point sources to get $\I^*_{\bar\theta_N}r$; and the application of $(\I^*\I)^{-1}$, which by \eqref{m:eq:fredholm} is a Laplacian type operator plus a compact correction and can be applied by a few conjugate gradient iterations each costing one forward and one adjoint linearized solve. This is the adjoint machinery of variational data assimilation (4DVar), introduced for meteorology by Le Dimet and Talagrand~\cite{ledimettalagrand1986} and treated mathematically by Law, Stuart and Zygalakis~\cite{law2015}, used here not to find the MAP estimate but to construct the summary statistic. Everywhere below, where the fold index is immaterial we write $\bar\theta_N$ for either estimate: the two satisfy the same rate, and the localization set of Section~\ref{m:sec:loc} and the plug in covariance may be built from either.

\begin{remark}[The regularized summary, and why it is the one to compute]\label{m:rem:reg}
Equation~\eqref{m:eq:onestep} inverts the population information $\I^*\I$ and carries no prior term, because both choices are asymptotically immaterial for a fixed prior and a fixed number of modes: $\Sigma_0^{-1}$ is then $O(1)$ against an information of order $N$, and the realized Gram matrix divided by $N$ concentrates on $\diag(\iota_j)$ at rate $N^{-1/2}$, by Chebyshev's inequality applied to its entries. Neither is immaterial at the sample sizes of Section~\ref{m:sec:heat}. Writing $F_N$ for the realized information $\I^*_{\bar\theta_N}\I_{\bar\theta_N}$ evaluated on the design actually drawn and summed over the $N$ observations, the regularized form
\begin{equation}\label{m:eq:regsummary}
\hat S^{\rm reg}_N=\bar\theta_N+\big(F_N+\Sigma_0^{-1}\big)^{-1}\Big(\I^*_{\bar\theta_N}\Big[\sum_ir_i\,\delta_{X_i}\Big]-\Sigma_0^{-1}\bar\theta_N\Big)
\end{equation}
costs the same one forward and one adjoint solve, and in the linear Gaussian case is exactly the posterior mean for every $N$ and independent of $\bar\theta_N$, by Lemma~\ref{m:lem:exact}. Writing $\hat S^{\rm real}_N=\bar\theta_N+F_N^{-1}\I^*_{\bar\theta_N}\big[\sum_ir_i\,\delta_{X_i}\big]$ for the one step summary with the realized information in place of the population one, \eqref{m:eq:regsummary} rearranges to $\hat S^{\rm reg}_N=(F_N+\Sigma_0^{-1})^{-1}F_N\hat S^{\rm real}_N$, that is, $\hat S^{\rm reg}_N-\hat S^{\rm real}_N=-(F_N+\Sigma_0^{-1})^{-1}\Sigma_0^{-1}\hat S^{\rm real}_N$: the regularized summary is the realized one shrunk toward the prior mean. It is therefore $\sqrt N$ equivalent to $\hat S_N$ in $H^{-k-2}$ whenever this shrinkage and the change of information $\hat S^{\rm real}_N-\hat S_N$ are $o_P(N^{-1/2})$ in that norm, and every statement below that uses the summary only through its distance to the centering \eqref{m:eq:thetahat} in $H^{-k-2}$ then applies to it; what concerns the posterior given the summary (Assumption~\ref{m:ass:U}) or the summary in $L^2$ and $H^{\bar\gamma}$ (Theorem~\ref{m:thm:qnnA}) has to be checked for \eqref{m:eq:regsummary} itself. For a fixed prior and a fixed number of modes both are. Let $G_N$ be the population counterpart of $F_N$, the information $T^{-1}\I^*_{\bar\theta_N}\I_{\bar\theta_N}$ times the number of residuals summed, and suppose $(F_N-G_N)/N=O_P(N^{-1/2})$ and $G_N/N$ has an invertible limit, as for the heat and Stokes equations, where $\I_{\bar\theta_N}$ does not depend on $\bar\theta_N$ and the entries of $\Phi^\top\Phi/N$ are averages of independent bounded variables. Then $\norm{(F_N+\Sigma_0^{-1})^{-1}}=O_P(N^{-1})$ while $\hat S^{\rm real}_N=O_P(1)$, and $\hat S^{\rm real}_N-\hat S^{\rm pop}_N=F_N^{-1}(G_N-F_N)(\hat S^{\rm pop}_N-\bar\theta_N)$, for $\hat S^{\rm pop}_N$ the same statistic with $G_N$ in place of $F_N$, is $O_P(N^{-1/2})$ times the one step correction, which tends to zero. For a linear forward map $\hat S^{\rm pop}_N$ and each fold of $\hat S_N$ are $\theta_0$ plus the efficient term of their own residuals plus $G^{-1}(G-F)(\bar\theta-\theta_0)=o_P(N^{-1/2})$, formed with their own $F$, $G$ and $\bar\theta$, so that $\hat S^{\rm pop}_N-\hat S_N=o_P(N^{-1/2})$ as well. When the number of modes grows like $J_N$ under the rescaled prior \eqref{m:eq:prior}, the shrinkage is not small on the modes with $\kappa_j\le1$ of Proposition~\ref{m:prop:pr}, and this argument fails. Propositions~\ref{m:prop:regheat} and~\ref{m:prop:regfixed} below prove the $\sqrt N$ equivalence of \eqref{m:eq:regsummary} with the centering \eqref{m:eq:thetahat} in $H^{-k-2}$ for the heat and Stokes equations, and Proposition~\ref{m:prop:regrd} for the reaction diffusion class; the equivalence with the one step summary of Definition~\ref{m:def:summary} follows for the heat equation and for the reaction diffusion class wherever Proposition~\ref{m:prop:suff} applies. For other nonlinear forward maps neither is proved, and Remark~\ref{m:rem:regnonlin} lists what a proof along the same lines needs. The difference at finite $N$ is not small. On the heat equation at $J=20$, $N=8000$, $M=4000$, over twelve replicates, GBC built on \eqref{m:eq:regsummary} reproduces the exact posterior standard deviation to within $0.3\%$ in every regime, while GBC built on \eqref{m:eq:onestep} overstates it by a factor $3.0\times10^4$ when the prior dominates, and, when the data dominate, misplaces the posterior mean by $410$ posterior standard deviations on average over the modes while getting the spread right. The second failure is caused by $N\iota_j$ in place of $F_N$ and is untouched by the training design of Algorithm~\ref{m:alg:loc}: it is the same in every radius and with or without the rescaling. We therefore compute \eqref{m:eq:regsummary} throughout and read \eqref{m:eq:onestep} as its asymptotic form.
\end{remark}

\begin{proposition}[Equivalence with the efficient centering]\label{m:prop:suff}
Suppose that Conditions 1 and 2 of Nickl~\cite{nicklbvm2024} hold with $\gamma$ and $k$ as in Theorem~\ref{m:thm:nickl2}, let $\kappa$ satisfy $2d+\zeta<\kappa<k-d/2$ as in (30) there, and suppose that $\sqrt N\,\tilde\delta_N(0)\tilde\delta_N(\zeta)\to0$, that $\gamma<(\bar\gamma+1)(\kappa-d/2)$ unless $\I^*_{\theta_0}\I_{\theta_0}$ is diagonal in $(e_j)$ or the preliminary estimates converge at the rate $\delta_N$ in $H^{-1}$, and that the weights of \eqref{m:eq:summary} depend on the linearization point in a Lipschitz way: for every $B>0$ there is $L=L(B)$ such that for all $\vartheta$ with $\norm\vartheta_{H^{\bar\gamma}}\le B$ and all $j\le J_N$
\begin{equation}\label{m:eq:lipweight}
\big\|\I_\vartheta p_N\big(\bar e_j^{\,\vartheta}\big)-\I_{\theta_0}p_N\big(\bar e_j^{\,\theta_0}\big)\big\|_{L^2(\mathcal X)}\le L\,\norm{\vartheta-\theta_0}_{H^\zeta}\norm{e_j}_{H^{\kappa+2}} .
\end{equation}
Then the projected one step statistic $\hat S_N^{J_N}=\sum_{j\le J_N}S_{N,j}e_j$ satisfies
\[
\sqrt N\,\norm{\hat S^{J_N}_N-\hat\theta_N}_{H^{-k-2}}\to^P0 ,
\]
and consequently $\Wone_1\big(\Law(\sqrt N(\theta-\hat S_N^{J_N})\mid Z^{(N)}),\Ncal_{\theta_0}\big)\to^P0$.
\end{proposition}

\begin{proof}
In this proof, as in Nickl~\cite{nicklbvm2024}, $L^2(\mathcal X)$ carries the law of the design, the uniform probability on $\mathcal X$; this divides squared norms by $T$ and turns $T(\I^*\I)^{-1}$ into the inverse of the information operator of this normalization, so that $\bar e_j$ is unchanged and (51) there and the conditional moments below hold as written. Fix a fold $b$ and argue on it, writing $\bar\theta_N$ for $\bar\theta^{[b]}_N$, $\mathcal I=\mathcal I_{3-b}$ and $n=|\mathcal I|$; the average over the two folds in \eqref{m:eq:summary} then obeys the same bounds by the triangle inequality. Since $\bar\theta_N$ is bounded in $H^{\bar\gamma}$ in probability (Definition~\ref{m:def:summary}), it suffices to argue on the events $\{\norm{\bar\theta_N}_{H^{\bar\gamma}}\le B\}$, whose probability tends to one as $B\to\infty$ uniformly in $N$, and on which Conditions 2C, 2D and \eqref{m:eq:lipweight} apply at $\bar\theta_N$. For $j\le J_N$ let $w_j=\I_{\theta_0}p_N(\bar e_j)$ be the weight of \eqref{m:eq:thetahat}, $w^b_j=\I_{\bar\theta_N}p_N(\bar e_j^{\,\bar\theta_N})$ the weight of \eqref{m:eq:summary}, and $C_j=\norm{e_j}_{H^{\kappa+2}}=(1+\lambda_j)^{(\kappa+2)/2}$. By (52) of Nickl~\cite{nicklbvm2024} and the choice of $\kappa$, $\norm{\bar e_j}_{H^\kappa}\lesssim C_j$, so $\norm{w_j}_{L^2(\mathcal X)}\lesssim C_j$ by continuity of $\I_{\theta_0}$ on $L^2$ (Condition 2C), and $\norm{w^b_j}_{L^2(\mathcal X)}\lesssim C_j$ by \eqref{m:eq:lipweight}. The constants below do not depend on $j$, and the only random ones are $\norm{\bar\theta_N-\theta_0}_{H^\zeta}/\tilde\delta_N(\zeta)$ and $\norm{\bar\theta_N-\theta_0}_{L^2}/\tilde\delta_N(0)$, which are $O_P(1)$. With $r_i=\varepsilon_i+(u_{\theta_0}-u_{\bar\theta_N})(X_i)$ the fold contributes $\langle e_j,\bar\theta_N\rangle+n^{-1}\sum_{i\in\mathcal I}\langle w^b_j(X_i),r_i\rangle_W$ to $S_{N,j}$, and we compare it with $\hat\Psi^b_{N,j}=\langle e_j,\theta_0\rangle+n^{-1}\sum_{i\in\mathcal I}\langle w_j(X_i),\varepsilon_i\rangle_W$. The average of $\hat\Psi^1_{N,j}$ and $\hat\Psi^2_{N,j}$ differs from $\hat\Psi_{N,j}$ of \eqref{m:eq:thetahat} only in that each summand carries the weight $1/(2n)$ in place of $1/N$; the halves differ in size by at most one, so the weights differ by $O(N^{-2})$ and the difference has variance $O(N^{-3}C_j^2)$, which is negligible in Step 4. The difference between the contribution of the fold and $\hat\Psi^b_{N,j}$ splits into a bias term and a stochastic term, and the coordinates $j>J_N$, which the summary omits, give a truncation term.

\emph{Step 1: the bias term.} It is $\langle e_j,\bar\theta_N-\theta_0\rangle+n^{-1}\sum_{i\in\mathcal I}\langle w^b_j(X_i),(u_{\theta_0}-u_{\bar\theta_N})(X_i)\rangle_W$. The pairs $(X_i,\varepsilon_i)$, $i\in\mathcal I$, are independent of $\bar\theta_N$, so the average is its conditional mean $\langle w^b_j,u_{\theta_0}-u_{\bar\theta_N}\rangle_{L^2(\mathcal X)}$ plus a fluctuation $F_j$ with conditional mean zero and conditional second moment at most $n^{-1}\norm{u_{\theta_0}-u_{\bar\theta_N}}^2_\infty\norm{w^b_j}^2_{L^2(\mathcal X)}\lesssim N^{-1}\tilde\delta_N(\zeta)^2C_j^2$, by Condition 2D. Condition 2C linearizes $u_{\theta_0}-u_{\bar\theta_N}=-\I_{\theta_0}[\bar\theta_N-\theta_0]-R$ with $\norm R_{L^2(\mathcal X)}\lesssim\tilde\delta_N(0)^2$, and the identity $\langle\I_{\theta_0}\bar e_j,\I_{\theta_0}h\rangle_{L^2(\mathcal X)}=\langle e_j,h\rangle$, equation (51) of Nickl~\cite{nicklbvm2024}, cancels $\langle e_j,\bar\theta_N-\theta_0\rangle$. Three remainders survive: the quadratic term $|\langle w^b_j,R\rangle|\lesssim\tilde\delta_N(0)^2C_j$; the change of weight $|\langle w^b_j-w_j,\I_{\theta_0}[\bar\theta_N-\theta_0]\rangle|\lesssim\tilde\delta_N(\zeta)\tilde\delta_N(0)C_j$, by \eqref{m:eq:lipweight}; and the projection error $|\langle\I_{\theta_0}(\bar e_j-p_N\bar e_j),\I_{\theta_0}[\bar\theta_N-\theta_0]\rangle|\lesssim\lambda_{J_N}^{-\kappa/2}\norm{\bar e_j}_{H^\kappa}\tilde\delta_N(0)\lesssim J_N^{-\kappa/d}\tilde\delta_N(0)C_j$.

\emph{Step 2: the stochastic term.} It is $B_j=n^{-1}\sum_{i\in\mathcal I}\langle(w^b_j-w_j)(X_i),\varepsilon_i\rangle_W$. By Definition~\ref{m:def:summary} the estimate $\bar\theta_N$ is computed from $Z^{[b]}$, which is disjoint from the half $\mathcal I$ being averaged, so given $\bar\theta_N$ the summands are independent and centered, and $E[B_j^2\mid\bar\theta_N]=n^{-1}\norm{w^b_j-w_j}^2_{L^2(\mathcal X)}\lesssim N^{-1}\tilde\delta_N(\zeta)^2C_j^2$ by \eqref{m:eq:lipweight}. Without the cross fitting this step is unavailable, since the weights would depend on the same errors being averaged.

\emph{Step 3: the truncation term.} Equation (69) of Nickl~\cite{nicklbvm2024}, summed over $j>J_N$ only, gives $N\,E\norm{\sum_{j>J_N}(\hat\Psi_{N,j}-\langle e_j,\theta_0\rangle)e_j}^2_{H^{-k-2}}\lesssim\sum_{j>J_N}\lambda_j^{-k}\to0$, and the coordinates of $\theta_0$ satisfy $N\norm{\sum_{j>J_N}\langle e_j,\theta_0\rangle e_j}^2_{H^{-k-2}}\le N\lambda_{J_N}^{-k-2-\gamma}\norm{\theta_0}^2_{H^\gamma}\to0$, since $\lambda_{J_N}\asymp J_N^{2/d}\asymp N^{2/(2\gamma+d)}$ and $k+2>d/2$.

\emph{Step 4: assembly.} For the fold $b$ and $j\le J_N$, Steps 1 and 2 write the difference between its contribution and $\hat\Psi^b_{N,j}$ as $a_j+F_j+B_j$, with $|a_j|\lesssim\eta_NC_j$, where $\eta_N=\tilde\delta_N(0)^2+\tilde\delta_N(\zeta)\tilde\delta_N(0)+J_N^{-\kappa/d}\tilde\delta_N(0)$, and $E[F_j^2+B_j^2\mid\bar\theta_N]\lesssim N^{-1}\tilde\delta_N(\zeta)^2C_j^2$. Since $\norm{\sum_{j\le J_N}c_je_j}^2_{H^{-k-2}}\le\sum_{j\le J_N}\lambda_j^{-k-2}c_j^2$, and $\lambda_j\asymp j^{2/d}$ makes $\sum_j\lambda_j^{-k-2}C_j^2$ finite because $k-\kappa>d/2$, these bounds may be summed over $j$, with Markov's inequality given $\bar\theta_N$ for the stochastic part. Averaging over the two folds and adding the fold size term of the first paragraph,
\[
\sqrt N\,\Big\|\sum_{j\le J_N}(S_{N,j}-\hat\Psi_{N,j})e_j\Big\|_{H^{-k-2}}\lesssim\sqrt N\,\eta_N+O_P\big(\tilde\delta_N(\zeta)+N^{-1}\big).
\]
Now $\tilde\delta_N(\xi)=\delta_N^{(\bar\gamma-\xi)/(\bar\gamma+1)}$ is nondecreasing in $\xi$, because $\delta_N<1$ and the exponent decreases in $\xi$; hence $\tilde\delta_N(0)\le\tilde\delta_N(\zeta)$ and so $\tilde\delta_N(0)^2\le\tilde\delta_N(0)\tilde\delta_N(\zeta)$. The hypothesis $\sqrt N\,\tilde\delta_N(0)\tilde\delta_N(\zeta)\to0$ therefore removes the first two terms of $\sqrt N\,\eta_N$, and $\tilde\delta_N(\zeta)\to0$ removes the stochastic terms. The third term of $\eta_N$ bounds the projection error of Step 1 and may be dropped when $\I^*_{\theta_0}\I_{\theta_0}$ is diagonal in $(e_j)$, since then $p_N\bar e_j=\bar e_j$ for $j\le J_N$ and that error is zero. When $\norm{\bar\theta_N-\theta_0}_{H^{-1}}=O_P(\delta_N)$, the bound $\norm{\I_{\theta_0}h}_{L^2(\mathcal X)}\lesssim\norm h_{H^{-1}}$ for $h\in L^2_0(\Omega)$, which Remark~\ref{m:rem:infobound} derives from Condition 2, bounds that error by a multiple of $\norm{\bar e_j-p_N\bar e_j}_{H^{-1}}\norm{\bar\theta_N-\theta_0}_{H^{-1}}\lesssim\lambda_{J_N}^{-(\kappa+1)/2}C_j\norm{\bar\theta_N-\theta_0}_{H^{-1}}$, so the third term may be replaced by $\lambda_{J_N}^{-(\kappa+1)/2}\delta_N$, with $\norm{\bar\theta_N-\theta_0}_{H^{-1}}/\delta_N=O_P(1)$ as a further random constant, and $\sqrt N\lambda_{J_N}^{-(\kappa+1)/2}\delta_N\asymp N^{(d/2-\kappa-1)/(2\gamma+d)}\to0$. Otherwise $\sqrt N\,J_N^{-\kappa/d}\tilde\delta_N(0)\asymp N^{(\gamma/(\bar\gamma+1)+d/2-\kappa)/(2\gamma+d)}$, which tends to zero because $\gamma<(\bar\gamma+1)(\kappa-d/2)$; this is the bound $\sqrt N\,\tilde\delta_N(0)J_N^{-\kappa/d}=o(1)$ that the proof of Theorem~4 of Nickl~\cite{nicklbvm2024} uses after (58) there. Combining with Step 3 gives $\sqrt N\norm{\hat S^{J_N}_N-\hat\theta_N}_{H^{-k-2}}\to^P0$. The last assertion is Theorem~\ref{m:thm:nickl2} with the centering $\hat\theta_N$, equation (80) there, together with translation invariance of $\Wone_1$.
\end{proof}

Condition \eqref{m:eq:lipweight} holds with $L=0$ for the heat and Stokes equations, where $\I_\vartheta$ does not depend on $\vartheta$. For a nonlinear forward map it asks that the linearization $\vartheta\mapsto\I_\vartheta$ and the inverse information $\vartheta\mapsto(\I^*_\vartheta\I_\vartheta)^{-1}$ be Lipschitz near $\theta_0$ in the norms shown, which Condition 2 does not supply: Condition 2C linearizes at $\theta_0$ only, and Condition 2D bounds $\Gcal(\vartheta)-\Gcal(\theta_0)$, not $\I_\vartheta-\I_{\theta_0}$. We therefore impose it.

Proposition~\ref{m:prop:suff} is a statement about the posterior given all the data: it says that $\hat S^{J_N}_N$ is a computable center for the Gaussian limit of $\PiN$. It does not say that the posterior given $\hat S^{J_N}_N$ alone is close to $\PiN$, that is, that the summary is sufficient for the whole law at the scale $\sqrt N$. That is Theorem~\ref{m:thm:jwbvm}, which holds under Assumption~\ref{m:ass:U}, the uniform Bernstein von Mises statement that is left open.

Proposition~\ref{m:prop:suff} controls the summary in $H^{-k-2}$, the norm of the Bernstein von Mises theorem. Theorem~\ref{m:thm:qnnA} needs a bound on its error in $L^2$ and in $H^{\bar\gamma}$ as well, and the next proposition supplies one under two conditions. The first is a two sided bound on the linearizations at which the weights of \eqref{m:eq:summary} are computed: for every $B>0$ there are constants $0<c_B\le C_B<\infty$ such that, for every $\vartheta$ in the parameter space with $\norm\vartheta_{H^{\bar\gamma}}\le B$,
\begin{equation}\label{m:eq:infobound}
c_B\norm h_{H^{-1}}\le\norm{\I_\vartheta h}_{L^2(\mathcal X)}\le C_B\norm h_{H^{-1}}\qquad\text{for all }h\in L^2_0(\Omega),
\end{equation}
with $\Hds{-1}$ and $\Hds{}$ in place of $H^{-1}$ and $L^2_0(\Omega)$ for the Stokes and Navier Stokes equations. The second is the rate of Definition~\ref{m:def:summary} extended to $\xi=-1$, where the formula for $\tilde\delta_N(\xi)$ gives $\tilde\delta_N(-1)=\delta_N$:
\begin{equation}\label{m:eq:pilotrate}
\norm{\bar\theta^{[b]}_N-\theta_0}_{H^{-1}}=O_P(\delta_N),\qquad b=1,2 .
\end{equation}

\begin{proposition}[The summary in Sobolev norms]\label{m:prop:sumrate}
Suppose that Conditions 1 and 2 of Nickl~\cite{nicklbvm2024} hold, that the linearizations in \eqref{m:eq:summary} satisfy \eqref{m:eq:infobound}, and that the preliminary estimates of Definition~\ref{m:def:summary} take values in the parameter space and satisfy \eqref{m:eq:pilotrate}. Then the projected one step statistic $\hat S^{J_N}_N=\sum_{j\le J_N}S_{N,j}e_j$ satisfies
\begin{equation}\label{m:eq:sumrate}
\norm{\hat S^{J_N}_N-\theta_0}_{H^s}=O_P\big(N^{(s+1-\gamma)/(2\gamma+d)}\big)\qquad\text{for every }-1\le s\le\gamma .
\end{equation}
If moreover $\bar\gamma<\gamma-1$, as in the range of Section~\ref{m:sec:main}, then $\norm{\hat S^{J_N}_N-\theta_0}_{L^2}=o_P(\tilde\delta_N(0))$ and $\norm{\hat S^{J_N}_N-\theta_0}_{H^{\bar\gamma}}=o_P(1)$, which is the hypothesis of Theorem~\ref{m:thm:qnnA} on the summary. For a linear forward map whose information operator is diagonal in $(e_j)$, as for the heat and Stokes equations, the same conclusions hold without \eqref{m:eq:pilotrate}.
\end{proposition}

\begin{proof}
Fix a fold $b$, write $\bar\theta=\bar\theta^{[b]}_N$, $h=\bar\theta-\theta_0$, $\mathcal I=\mathcal I_{3-b}$, $n=|\mathcal I|$ and $J=J_N$, let $w^b_j=\I_{\bar\theta}p_N(\bar e_j^{\,\bar\theta})$ be the weight of $e_j$ in \eqref{m:eq:summary}, as in the proof of Proposition~\ref{m:prop:suff}, and let $S^b_j$ be the contribution of the fold to $S_{N,j}$, the bracket in \eqref{m:eq:summary}, so that $\hat S^{J_N}_N-\theta_0=\frac12\sum_{b=1,2}\sum_{j\le J}(S^b_j-\langle e_j,\theta_0\rangle)e_j-(\theta_0-p_N\theta_0)$. Write $\rho_N(s)=(1+\lambda_J)^{(s+1)/2}\delta_N$. Since $\lambda_J\asymp J^{2/d}$ and $J\asymp N\delta_N^2=N^{d/(2\gamma+d)}$, $\rho_N(s)\asymp N^{(s+1-\gamma)/(2\gamma+d)}$, and $J/n\lesssim\delta_N^2$ because $n\ge(N-1)/2$. For $s\le\gamma$ the truncation term satisfies $\norm{\theta_0-p_N\theta_0}_{H^s}\le(1+\lambda_J)^{-(\gamma-s)/2}\norm{\theta_0}_{H^\gamma}$, which is $\rho_N(s)$ times a multiple of $(1+\lambda_J)^{-(\gamma+1)/2}\delta_N^{-1}\asymp N^{-1/(2\gamma+d)}$. Two elementary inequalities are used throughout. For $f$ in the span of $e_1,\dots,e_J$ and $s\ge-1$, $\norm f_{H^s}\le(1+\lambda_J)^{(s+1)/2}\norm f_{H^{-1}}$, since $\lambda_j\le\lambda_J$ for $j\le J$. For $-1\le\xi\le\bar\gamma$, H\"older's inequality applied to the coefficients gives $\norm f_{H^\xi}\le\norm f_{H^{-1}}^{(\bar\gamma-\xi)/(\bar\gamma+1)}\norm f_{H^{\bar\gamma}}^{(\xi+1)/(\bar\gamma+1)}$.

\emph{Step 1: the weights.} Fix $B\ge\norm{\theta_0}_{H^{\bar\gamma}}$ and $\vartheta$ in the parameter space with $\norm\vartheta_{H^{\bar\gamma}}\le B$. By the upper bound in \eqref{m:eq:infobound}, $\I_\vartheta$ extends by continuity to $H^{-1}_0$, and the form $(f,g)\mapsto\langle\I_\vartheta f,\I_\vartheta g\rangle_{L^2(\mathcal X)}$ is bounded on $H^{-1}_0$; by the lower bound it is coercive. The dual of $H^{-1}_0$ is $H^1_0$ under the pairing of $L^2$, so the Lax Milgram lemma gives, for every $\phi\in H^1_0$, a unique $f\in H^{-1}_0$ with $\langle\I_\vartheta f,\I_\vartheta g\rangle_{L^2(\mathcal X)}=\langle\phi,g\rangle$ for all $g\in H^{-1}_0$. This $f$ is $(\I^*_\vartheta\I_\vartheta)^{-1}\phi$, the information operator being taken on the mean zero space on which the parameter lives: when $\phi$ lies in the range of that operator on $L^2_0(\Omega)$ the two coincide by uniqueness, since $L^2_0(\Omega)\subset H^{-1}_0$, and otherwise $f$ is how we read the inverse in \eqref{m:eq:summary}. The choice $g=f$ gives $c_B^2\norm f^2_{H^{-1}}\le\norm\phi_{H^1}\norm f_{H^{-1}}$. Hence, if $\norm{\bar\theta}_{H^{\bar\gamma}}\le B$, then for $\phi=\sum_{j\le J}a_je_j$ the weight $w^b_\phi=\sum_{j\le J}a_jw^b_j=\I_{\bar\theta}p_N\big(T(\I^*_{\bar\theta}\I_{\bar\theta})^{-1}\phi\big)$ satisfies
\[
\norm{w^b_\phi}_{L^2(\mathcal X)}\le C_B\norm{p_NT(\I^*_{\bar\theta}\I_{\bar\theta})^{-1}\phi}_{H^{-1}}\le TC_Bc_B^{-2}\norm\phi_{H^1} .
\]
For $v\in L^2(\mathcal X)$ let $\beta(v)=\sum_{j\le J}T^{-1}\langle w^b_j,v\rangle_{L^2(\mathcal X)}e_j$. Then $\langle\beta(v),\phi\rangle=T^{-1}\langle w^b_\phi,v\rangle_{L^2(\mathcal X)}$, and the norm of $H^{-1}$ of an element of the span of $e_1,\dots,e_J$ is the supremum of its pairings with the $\phi$ in that span with $\norm\phi_{H^1}\le1$, so
\[
\norm{\beta(v)}_{H^{-1}}\le C_Bc_B^{-2}\norm v_{L^2(\mathcal X)} .
\]

\emph{Step 2: decomposition.} Let $E_B=\{\norm{\bar\theta}_{H^{\bar\gamma}}\le B\}$, an event determined by $Z^{[b]}$ whose probability tends to one as $B\to\infty$, uniformly in $N$, by Definition~\ref{m:def:summary}. Write $g=u_{\theta_0}-u_{\bar\theta}$, so that $Y_i-u_{\bar\theta}(X_i)=\varepsilon_i+g(X_i)$, and $Z_{ij}=\langle w^b_j(X_i),\varepsilon_i+g(X_i)\rangle_W$. Given $Z^{[b]}$, the pairs $(X_i,\varepsilon_i)$, $i\in\mathcal I$, are independent and identically distributed and independent of $\bar\theta$, $X_i$ has density $1/T$ on $\mathcal X$, and $\varepsilon_i$ is centered and independent of $X_i$. Hence $E[Z_{ij}\mid Z^{[b]}]=T^{-1}\langle w^b_j,g\rangle_{L^2(\mathcal X)}$ and
\[
\begin{gathered}
\sum_{j\le J}\big(S^b_j-\langle e_j,\theta_0\rangle\big)e_j=p_Nh+\beta(g)+\mathcal E,\\
\mathcal E=\sum_{j\le J}\Big(\frac1n\sum_{i\in\mathcal I}Z_{ij}-E\big[Z_{ij}\mid Z^{[b]}\big]\Big)e_j .
\end{gathered}
\]

\emph{Step 3: the conditional mean.} On $E_B$, $\norm h_{H^{\bar\gamma}}\le2B$, and Condition 2C together with the upper bound in \eqref{m:eq:infobound} at $\vartheta=\theta_0$ gives $\norm g_{L^2(\mathcal X)}\le\norm{\I_{\theta_0}h}_{L^2(\mathcal X)}+c\norm h^2_{L^2}\le C_B\norm h_{H^{-1}}+c\norm h^2_{L^2}$, with $c$ the constant of Condition 2C for the radius $2B$. With Step 1 and the first elementary inequality, on $E_B$
\[
\begin{aligned}
\norm{p_Nh+\beta(g)}_{H^s}&\le(1+\lambda_J)^{(s+1)/2}\Big[\big(1+C_B^2c_B^{-2}\big)\norm h_{H^{-1}}\\
&\qquad+cC_Bc_B^{-2}\norm h^2_{L^2}\Big].
\end{aligned}
\]
By \eqref{m:eq:pilotrate}, $\norm h_{H^{-1}}=O_P(\delta_N)$, and by the second elementary inequality with $\xi=0$, $\norm h^2_{L^2}\le\norm h^{2\bar\gamma/(\bar\gamma+1)}_{H^{-1}}(2B)^{2/(\bar\gamma+1)}=O_P(\delta_N)$ on $E_B$, because $\bar\gamma\ge1$ in Condition 2. The conditional mean $p_Nh+\beta(g)$ is therefore $O_P(\rho_N(s))$ in $H^s$.

\emph{Step 4: the fluctuation.} On $E_B$, Condition 2A gives $\norm{u_\vartheta}_\infty\le U_B$ for $\vartheta\in\{\theta_0,\bar\theta\}$, with $U_B$ its constant for the radius $B$, so $\norm g_\infty\le2U_B$. Since $\varepsilon_i\sim N(0,I_W)$ is independent of $X_i$, $E[Z_{ij}^2\mid Z^{[b]}]=E|w^b_j(X_i)|^2_W+E\langle w^b_j(X_i),g(X_i)\rangle^2_W\le(1+4U_B^2)\,T^{-1}\norm{w^b_j}^2_{L^2(\mathcal X)}$, and Step 1 with $\phi=e_j$ gives $\norm{w^b_j}_{L^2(\mathcal X)}\le TC_Bc_B^{-2}(1+\lambda_j)^{1/2}$. The coordinates of $\mathcal E$ are averages of $n$ conditionally independent centered terms, so on $E_B$
\[
\begin{aligned}
E\big[\norm{\mathcal E}^2_{H^s}\mid Z^{[b]}\big]&=\sum_{j\le J}(1+\lambda_j)^sE\big[\langle e_j,\mathcal E\rangle^2\mid Z^{[b]}\big]\le\frac{K_B}n\sum_{j\le J}(1+\lambda_j)^{s+1}\\
&\le K_B\frac Jn(1+\lambda_J)^{s+1}\lesssim K_B\,\rho_N(s)^2,
\end{aligned}
\]
with $K_B=(1+4U_B^2)\,TC_B^2c_B^{-4}$. Chebyshev's inequality given $Z^{[b]}$ yields $P\big(E_B\cap\{\norm{\mathcal E}_{H^s}>M\rho_N(s)\}\big)\lesssim K_B/M^2$, and since $P(E_B^c)$ is small for $B$ large, uniformly in $N$, $\norm{\mathcal E}_{H^s}=O_P(\rho_N(s))$.

\emph{Step 5: conclusion.} Steps 3 and 4, the triangle inequality over the two folds and the bound on the truncation term give \eqref{m:eq:sumrate}. If $\bar\gamma<\gamma-1$, the bound at $s=0$ is $o_P(\tilde\delta_N(0))$, since $N^{(1-\gamma)/(2\gamma+d)}/\tilde\delta_N(0)=N^{-(\gamma-\bar\gamma-1)/((2\gamma+d)(\bar\gamma+1))}\to0$, and the bound at $s=\bar\gamma$ is $o_P(1)$, since $\bar\gamma+1-\gamma<0$. For a linear forward map with $T^{-1}\I^*\I e_j=\iota_je_j$ for every $j$, as for the heat and Stokes equations by \eqref{m:eq:gram} and \eqref{m:eq:modeinfo}, $g=-\I h$ exactly and $\bar e_j=\iota_j^{-1}e_j$, so that for $j\le J$
\[
\langle e_j,\beta(g)\rangle=-\iota_j^{-1}T^{-1}\langle\I e_j,\I h\rangle_{L^2(\mathcal X)}=-\iota_j^{-1}\langle T^{-1}\I^*\I e_j,h\rangle=-\langle e_j,h\rangle .
\]
The conditional mean then vanishes identically, Step 3 is not needed, and Step 4 uses only that $\bar\theta$ is bounded in $H^{\bar\gamma}$ in probability, which Definition~\ref{m:def:summary} requires.
\end{proof}

\begin{lemma}[Preliminary estimates with the rate in $H^{-1}$]\label{m:lem:halfmean}
Suppose that Conditions 1 and 2 of Nickl~\cite{nicklbvm2024} hold with $\gamma$ and $k$ as in Theorem~\ref{m:thm:nickl2} and that $\bar\gamma<\gamma-d/2$. Let $\bar\theta^{[b]}_N$ be computed from $Z^{[b]}$ alone, with values in the parameter space and with $\norm{\bar\theta^{[b]}_N}_{H^{\bar\gamma}}=O_P(1)$. If $\norm{u_{\bar\theta^{[b]}_N}-u_{\theta_0}}_{L^2(\mathcal X)}=O_P(\delta_N)$, then \eqref{m:eq:pilotrate} holds, and \eqref{m:eq:pilotrate} implies the rates of Definition~\ref{m:def:summary} for every $-1\le\xi\le\bar\gamma$. The posterior mean computed from $Z^{[b]}$ with the prior \eqref{m:eq:prior} at sample size $|\mathcal I_b|$ satisfies \eqref{m:eq:pilotrate} and these rates; for the two dimensional Navier Stokes equations this holds in the setting of Theorem~2.11 of Konen and Nickl~\cite{konennickl2026} when $\bar\gamma\le11$.
\end{lemma}

\begin{proof}
On $\{\norm{\bar\theta^{[b]}_N}_{H^{\bar\gamma}}\le B\}$, Condition 2E gives $\norm{\bar\theta^{[b]}_N-\theta_0}_{H^{-1}}\le c^{-1}\norm{u_{\bar\theta^{[b]}_N}-u_{\theta_0}}_{L^2(\mathcal X)}$, with $c$ its constant for the radius $B$, which proves the first claim. Given \eqref{m:eq:pilotrate} and $\norm{\bar\theta^{[b]}_N-\theta_0}_{H^{\bar\gamma}}=O_P(1)$, the second elementary inequality of the proof of Proposition~\ref{m:prop:sumrate} gives $\norm{\bar\theta^{[b]}_N-\theta_0}_{H^\xi}=O_P(\delta_N^{(\bar\gamma-\xi)/(\bar\gamma+1)})=O_P(\tilde\delta_N(\xi))$ for $-1\le\xi\le\bar\gamma$, which is the second.

For the posterior mean $\bar\theta$, write $n_b=|\mathcal I_b|\ge(N-1)/2$, so that $\delta_{n_b}\asymp\delta_N$ and $n_b\delta^2_{n_b}=n_b^{d/(2\gamma+d)}$. Choose $\bar\gamma'$ with $\max(3\zeta+d+2,\bar\gamma)<\bar\gamma'<\gamma-d/2$, a nonempty range because $\gamma>3\zeta+3d/2+2$. Condition 2 holds with $\bar\gamma'$ in place of $\bar\gamma$, since balls of $H^{\bar\gamma'}$ lie in those of $H^{\bar\gamma}$, so Theorem~3 of Nickl~\cite{nicklbvm2024} applies to the data $Z^{[b]}$ with $\bar\gamma'$ in its range (31): by (37) and (38) there, for $L$ large and with probability tending to one, the posterior gives mass at least $1-e^{-n_b\delta^2_{n_b}}$ to a set on which $\norm\theta_{H^{\bar\gamma'}}\le L$ and $\norm{u_\theta-u_{\theta_0}}_{L^2(\mathcal X)}\le L\delta_{n_b}$, and by (39) there $E^\Pi[\norm{\theta-\theta_0}^2_{L^2}\mid Z^{[b]}]^{1/2}=O_P(1)$ and $\norm{\bar\theta-\theta_0}_{H^{\bar\gamma'}}=O_P(1)$. On that set Condition 2E gives $\norm{\theta-\theta_0}_{H^{-1}}\le c_L^{-1}L\delta_{n_b}$, with $c_L$ its constant for the radius $L$. Jensen's inequality for the Bochner integral, and the Cauchy Schwarz inequality off the set, give
\[
\norm{\bar\theta-\theta_0}_{H^{-1}}\le c_L^{-1}L\delta_{n_b}+E^\Pi\big[\norm{\theta-\theta_0}^2_{L^2}\mid Z^{[b]}\big]^{1/2}e^{-n_b\delta^2_{n_b}/2}=O_P(\delta_N),
\]
since $e^{-n_b\delta^2_{n_b}/2}$ tends to zero faster than any power of $N$. With $\norm{\bar\theta}_{H^{\bar\gamma}}\le\norm{\bar\theta}_{H^{\bar\gamma'}}=O_P(1)$, the second claim gives the rates.

For the two dimensional Navier Stokes equations, whose parameters are divergence free fields, Theorem~3 of Nickl is replaced by Theorem~2.13 of Konen and Nickl~\cite{konennickl2026}, whose proof is that of Theorem~3 of Nickl and which holds in the setting of their Theorem~2.11: their (2.31), (2.32) and (2.33) take the place of (37), (38) and (39), with their index $\beta=11$ in place of $\bar\gamma'$, and part E of their Condition~2.12, which they verify, takes the place of Condition 2E on the set (2.31). The argument above then applies unchanged. Since (2.33) there bounds the posterior moments in $\Hds{\xi}$ only for $0\le\xi\le\beta$, it gives $\norm{\bar\theta}_{\Hds{\bar\gamma}}=O_P(1)$, and with it the last assertion, when $\bar\gamma\le11$.
\end{proof}

\begin{remark}[Where \eqref{m:eq:infobound} holds]\label{m:rem:infobound}
At $\vartheta=\theta_0$ both bounds follow from Condition 2. For the lower bound, take $h\in H^{\bar\gamma}_0$ with $\norm h_{H^{\bar\gamma}}=1$, which loses nothing since both sides are homogeneous in $h$, and $0<t\le1$. Condition 2E at $\theta_0+th$ and Condition 2C give $c\,t\norm h_{H^{-1}}\le\norm{\Gcal(\theta_0+th)-\Gcal(\theta_0)}_{L^2(\mathcal X)}\le t\norm{\I_{\theta_0}h}_{L^2(\mathcal X)}+c't^2\norm h^2_{L^2}$, with constants for the radius $\norm{\theta_0}_{H^{\bar\gamma}}+1$; dividing by $t$ and letting $t\to0$ gives the bound on $H^{\bar\gamma}_0$, and the continuity of $\I_{\theta_0}$ on $L^2$ extends it to $L^2_0(\Omega)$. For the upper bound, let $S$ be the information operator $\I^*_{\theta_0}\I_{\theta_0}$ compressed to $L^2_0(\Omega)$, so that $\Delta S=\Delta\I^*_{\theta_0}\I_{\theta_0}$, let $D=(\Id-\Delta)^{1/2}$, and fix $\eta\ge\max(\eta_0,0)$. Condition 2F bounds $\Delta S$ on $H^\eta_0$, hence $S$ from $H^\eta_0$ to $H^{\eta+2}_0$, so that $D^{\eta+2}SD^{-\eta}$ is bounded on $L^2_0(\Omega)$, with norm $C$ say. The matrix $X_J=p_NDSDp_N$ on the span of $e_1,\dots,e_J$ is symmetric and similar to $D^{\eta+1}X_JD^{-\eta-1}=p_ND^{\eta+2}SD^{-\eta}p_N$, so its norm, which is its spectral radius, is at most $C$ for every $J$. Hence $\norm{\I_{\theta_0}h}^2_{L^2(\mathcal X)}=\langle Sh,h\rangle=\langle X_JD^{-1}h,D^{-1}h\rangle\le C\norm h^2_{H^{-1}}$ for $h$ in that span, and by density for every $h\in L^2_0(\Omega)$. Condition \eqref{m:eq:infobound} asks for the two bounds uniformly on balls of $H^{\bar\gamma}$.

For the heat and Stokes equations $\I_\vartheta=\Lcal$ for every $\vartheta$, and $\norm{\Lcal h}^2_{L^2(\mathcal X)}=\sum_j\frac{1-e^{-2\nu\lambda_jT}}{2\nu\lambda_j}\langle e_j,h\rangle^2$ by \eqref{m:eq:modeinfo}. On $[\lambda_1,\infty)$ the function $\lambda\mapsto(1+\lambda)(1-e^{-2\nu\lambda T})/(2\nu\lambda)$ lies between $(1-e^{-2\nu\lambda_1T})/(2\nu)$ and $(1+1/\lambda_1)/(2\nu)$, so \eqref{m:eq:infobound} holds with constants that do not depend on $B$.

For the reaction diffusion class of Theorem~\ref{m:thm:nickl1}, for which Theorem~5 of Nickl~\cite{nicklbvm2024} verifies Condition 2 with $\bar\gamma>3+d/2$, the linearization at $\vartheta$, which is the operator that the adjoint solve of the summary uses, is $\I_\vartheta h=U_h$, the solution of $\partial_tU-\Delta U-VU=0$, $U(0)=h$, with potential $V=f'(u_\vartheta)$, which is equation (107) there with source zero. The argument following (102) there bounds the norm in $C^1([0,T],C^1)$ of $f'(\tilde u)$, for $\tilde u$ any convex combination of two solutions whose initial conditions have $H^{\bar\gamma}$ norms summing to at most $B'$, by a constant $M(B')$, using $\bar\gamma>3+d/2$; taking both solutions equal to $u_\vartheta$ bounds the norm of $V$ by $M(2B)$. Proposition~5 there, in its case (i) with $a=-1$ and source zero, then gives the upper bound in \eqref{m:eq:infobound}, and Lemma~4 there gives the lower bound for $h\in H^1$, which extends to $h\in L^2_0(\Omega)$ by density and the upper bound; both constants depend only on $M(2B)$, $f$, $T$ and $d$.

For the two dimensional Navier Stokes equations with forcing $f\in\Hds{1}$, Propositions~3.9 and~3.10 of Konen and Nickl~\cite{konennickl2026}, applied with $a=-1$, for which their exponent $a^*$ equals $2$, give for every $L>0$ a constant $C$, depending on $\nu$, $T$, $\norm f_{\Hds{1}}$ and $L$, such that $\int_0^T\norm{\I_\vartheta[h](t)}^2_{L^2}dt\le C\norm h^2_{\Hds{-1}}$ and $\norm h^2_{\Hds{-1}}\le C\int_0^T\norm{\I_\vartheta[h](t)}^2_{L^2}dt$ for all $\vartheta\in\Hds{2}$ with $\norm\vartheta_{\Hds{2}}\le L$ and all $h\in\Hds{-1}$, in their homogeneous norms, which on mean zero fields are equivalent to ours. Since $\bar\gamma>2$ in the range of Section~\ref{m:sec:main}, a ball of $\Hds{\bar\gamma}$ lies in such a ball, and \eqref{m:eq:infobound} holds.
\end{remark}

\begin{corollary}[The summary under the hypotheses of Proposition~\ref{m:prop:suff}]\label{m:cor:sumsuff}
Under the hypotheses of Proposition~\ref{m:prop:suff}, for every $-1\le s\le\gamma$,
\[
\norm{\hat S^{J_N}_N-\theta_0}_{H^s}=O_P\big(N^{(s+1-\gamma)/(2\gamma+d)}\big)+o_P\big(N^{(s+k+2)/(2\gamma+d)-1/2}\big).
\]
\end{corollary}

\begin{proof}
Since $\hat S^{J_N}_N$ lies in the span of $e_1,\dots,e_{J_N}$, with $\hat\theta_N$ the centering \eqref{m:eq:thetahat},
\[
\hat S^{J_N}_N-\theta_0=p_N\big(\hat S^{J_N}_N-\hat\theta_N\big)+p_N\big(\hat\theta_N-\theta_0\big)-\big(\theta_0-p_N\theta_0\big).
\]
For $f$ in that span, $\norm f_{H^s}\le(1+\lambda_{J_N})^{(s+k+2)/2}\norm f_{H^{-k-2}}$, and $p_N$ does not increase the norm of $H^{-k-2}$, so Proposition~\ref{m:prop:suff} bounds the first term by $o_P\big((1+\lambda_{J_N})^{(s+k+2)/2}N^{-1/2}\big)$, where $(1+\lambda_{J_N})^{(s+k+2)/2}\asymp N^{(s+k+2)/(2\gamma+d)}$. In the second term the coordinate of $e_j$, $j\le J_N$, is the average of the $N$ independent centered variables $\langle w_j(X_i),\varepsilon_i\rangle_W$ with $w_j=\I_{\theta_0}p_N(\bar e_j)$, whose second moment is $T^{-1}\norm{w_j}^2_{L^2(\mathcal X)}$. By Remark~\ref{m:rem:infobound}, \eqref{m:eq:infobound} holds at $\vartheta=\theta_0$, with constants $c_0,C_0$ say, and Step 1 of the proof of Proposition~\ref{m:prop:sumrate}, with $\theta_0$ in place of $\bar\theta$, gives $\norm{w_j}_{L^2(\mathcal X)}\le TC_0c_0^{-2}(1+\lambda_j)^{1/2}$. Hence $E\norm{p_N(\hat\theta_N-\theta_0)}^2_{H^s}\le TC_0^2c_0^{-4}N^{-1}\sum_{j\le J_N}(1+\lambda_j)^{s+1}$, which is of order $N^{2(s+1-\gamma)/(2\gamma+d)}$, and Markov's inequality bounds the second term. The third term is bounded in the proof of Proposition~\ref{m:prop:sumrate}.
\end{proof}

\begin{remark}[What the rate in $H^{-1}$ is for]\label{m:rem:sumopen}
The proof of Proposition~\ref{m:prop:sumrate}, whose notation we keep, bounds the error $p_Nh$, $h=\bar\theta_N-\theta_0$, of the preliminary estimate and the conditional mean $\beta(g)$ of the correction separately, and \eqref{m:eq:pilotrate} is what makes each of them small, although the correction is built to cancel the error. Let $R=u_{\bar\theta_N}-u_{\theta_0}-\I_{\bar\theta_N}h$ be the remainder of the linearization at $\bar\theta_N$, so that $g=-\I_{\bar\theta_N}h-R$. With the information operator taken on the mean zero space, Step 1 of that proof gives $\beta(v)=p_N(\I^*_{\bar\theta_N}\I_{\bar\theta_N})^{-1}p_N\I^*_{\bar\theta_N}v$ for $v\in L^2(\mathcal X)$, and $p_Nh=p_N(\I^*_{\bar\theta_N}\I_{\bar\theta_N})^{-1}\I^*_{\bar\theta_N}\I_{\bar\theta_N}p_Nh$, so that
\[
p_Nh+\beta(g)=p_N\big(\I^*_{\bar\theta_N}\I_{\bar\theta_N}\big)^{-1}\big[\I^*_{\bar\theta_N}\I_{\bar\theta_N},p_N\big]h-\beta(R),
\]
where $[\,\cdot\,,\cdot\,]$ is the commutator. With $\theta_0$ in place of $\bar\theta_N$, the first term is, up to sign, the vector of the projection errors of Step 1 of the proof of Proposition~\ref{m:prop:suff}, which that proof makes small in $H^{-k-2}$ through the rate $\delta_N$ in $H^{-1}$ or, under the rates of Definition~\ref{m:def:summary} alone, through the condition $\gamma<(\bar\gamma+1)(\kappa-d/2)$. If $\norm R_{L^2(\mathcal X)}$ is at most a multiple of $\norm h^2_{L^2}$ uniformly on balls of $H^{\bar\gamma}$, a linearization at points other than $\theta_0$ that Condition 2C does not state, the second term is $O_P\big((1+\lambda_{J_N})^{(s+1)/2}\tilde\delta_N(0)^2\big)$ in $H^s$, which is $o_P(\tilde\delta_N(0))$ for $s=0$ and $o_P(1)$ for $s=\bar\gamma$ when $1\le\bar\gamma<\gamma-1$. The first term vanishes when the information operator at $\bar\theta_N$ is diagonal in $(e_j)$. In general it couples the retained modes with the discarded ones, and \eqref{m:eq:infobound} bounds it in $H^s$ only by a multiple of $(1+\lambda_{J_N})^{(s+1)/2}\norm h_{H^{-1}}$. Under \eqref{m:eq:pilotrate} this bound is $O_P(N^{(s+1-\gamma)/(2\gamma+d)})$, but for $h=\tilde\delta_N(0)e_1$, an error that meets all the rates of Definition~\ref{m:def:summary}, it exceeds $\tilde\delta_N(0)$ in $L^2$. The first term itself need not be small. For the linear equation $\partial_tu=\partial_x(a\,\partial_xu)$ on the circle with $a(x)=2+\cos(2\pi x)$ and $e$ the retained cosine mode of highest frequency, the first term at $h=e$ has norm in $L^2$ converging to $(2-\sqrt3)/(2\sqrt3)\approx0.077$ as $N\to\infty$, because the information operator and its inverse couple $e$ with the first discarded cosine mode through the first Fourier coefficients of $1/a$ and of $a$. Since $h=\tilde\delta_N(0)e$ meets all the rates of Definition~\ref{m:def:summary} when $\bar\gamma<\gamma-1$, those rates alone do not give the hypothesis of Theorem~\ref{m:thm:qnnA} in general, even for a linear forward map. For a nonlinear forward map the hypothesis of Theorem~\ref{m:thm:qnnA} on the summary is therefore open in general when the preliminary estimate has only those rates, and also when \eqref{m:eq:infobound} is not known uniformly on balls of $H^{\bar\gamma}$, since Step 1 of the proof of Proposition~\ref{m:prop:sumrate} uses it at $\bar\theta_N$. A proof along the lines of Proposition~\ref{m:prop:sumrate} needs the linearization just described and a bound on the commutator of $\I^*_\vartheta\I_\vartheta$ with $p_N$, uniformly on balls of $H^{\bar\gamma}$, that is sharper than the one \eqref{m:eq:infobound} gives. For the reaction diffusion class Lemma~\ref{m:lem:rdinfo} below supplies both, and Proposition~\ref{m:prop:sumrd} proves the hypothesis for every preliminary estimate.
\end{remark}

\begin{lemma}[The information operator of the reaction diffusion class]\label{m:lem:rdinfo}
In the setting of Theorem~\ref{m:thm:nickl1}, with $\bar\gamma>3+d/2$ as in Theorem~5 of Nickl~\cite{nicklbvm2024}, let $\Lcal$ be the solution operator of $\partial_tu=\Delta u$, as in \eqref{m:eq:modeinfo} with $\nu=1$, and $\Kcal_\vartheta=\I^*_\vartheta\I_\vartheta-\Lcal^*\Lcal$, both compressed to $L^2_0(\Omega)$. For every $B>0$ there is a constant $C$, depending on $B$, $f$, $T$ and $d$, such that for all $\vartheta$ and $\vartheta+h$ with $H^{\bar\gamma}$ norms at most $B$:
\begin{enumerate}[label=(\roman*),nosep]
\item $|\langle\Kcal_\vartheta\phi,\psi\rangle|\le C\norm\phi_{H^{-1}}\norm\psi_{H^{-3}}$ for $\phi,\psi\in L^2_0(\Omega)$, so that the symmetric operator $\Kcal_\vartheta$ extends to maps from $H^{-1}_0$ to $H^3_0$ and from $H^{-3}_0$ to $H^1_0$ with norms at most $C$;
\item $\norm{\Gcal(\vartheta+h)-\Gcal(\vartheta)-\I_\vartheta h}_{L^2(\mathcal X)}\le C\norm h^2_{L^2}$;
\item the weights of \eqref{m:eq:summary} satisfy \eqref{m:eq:lipweight} for every $\kappa\ge-1$.
\end{enumerate}
\end{lemma}

\begin{proof}
Write $V=f'(u_\vartheta)$, $K_f=\sup|f'|$, $L^2H^b=L^2([0,T],H^b(\Omega))$ with the norms (13) of Nickl~\cite{nicklbvm2024}, which include the constant mode, and $\mathcal Z$ for the Duhamel operator of the proof of Lemma~\ref{m:lem:rdlin}, so that $U_\phi=\I_\vartheta\phi$ solves $U_\phi=\Lcal\phi+\mathcal Z(VU_\phi)$; constants depend only on $B$, $f$, $T$ and $d$. By (90) of Nickl with $\gamma=\bar\gamma$, the embedding of $H^{\bar\gamma}$ into $C^3$ and the chain rule, $\sup_{t\le T}\norm{V(t)}_{C^3}$ is bounded, so multiplication by $V(t)$ is bounded on $H^b$ for every integer $|b|\le3$, by the Leibniz rule and duality. On each mode $\mathcal Z$ is convolution with $e^{-\lambda_jr}1_{r>0}$, whose integral over $[0,T]$ is at most $\min(T,\lambda_j^{-1})\le(1+T)(1+\lambda_j)^{-1}$, so by Young's inequality $\mathcal Z$ maps $L^2H^b$ to $L^2H^{b+2}$ with norm at most $1+T$ for every real $b$; \eqref{m:eq:modeinfo} gives $\norm{\Lcal\phi}_{L^2H^{b+1}}\lesssim\norm\phi_{H^b}$, and Remark~\ref{m:rem:infobound} gives $\norm{U_\phi}_{L^2H^0}\le C_B\norm\phi_{H^{-1}}$.

(i) For $\psi\in H^{-3}$ the mild equation has, by Picard iteration, a unique solution $U_\psi$ in $C([0,T],H^{-3})$, equal to $\I_\vartheta\psi$ when $\psi\in L^2_0(\Omega)$. Gr\"onwall's inequality bounds $\sup_t\norm{U_\psi(t)}_{H^{-3}}$ by a multiple of $\norm\psi_{H^{-3}}$, and the mild equation then gives $\norm{U_\psi}_{L^2H^{-2}}\lesssim\norm\psi_{H^{-3}}$. With $R_\phi=\mathcal Z(VU_\phi)$, substituting $U_\psi=\Lcal\psi+R_\psi$ and then $U_\phi=\Lcal\phi+R_\phi$ gives $\langle\Kcal_\vartheta\phi,\psi\rangle=\langle R_\phi,\Lcal\psi\rangle_{L^2(\mathcal X)}+\langle U_\phi,R_\psi\rangle_{L^2(\mathcal X)}$, where $\norm{R_\phi}_{L^2H^2}\le(1+T)K_fC_B\norm\phi_{H^{-1}}$, $\norm{\Lcal\psi}_{L^2H^{-2}}\lesssim\norm\psi_{H^{-3}}$ and $\norm{R_\psi}_{L^2H^0}\lesssim\norm{U_\psi}_{L^2H^{-2}}$.

(ii) The remainder $r=u_{\vartheta+h}-u_\vartheta-\I_\vartheta h$ solves (105) of Nickl with $f'(u_\vartheta)$ in place of $f'(u_{\theta_0})$, a source $m$ and $r(0)=0$. Testing with $r$ and Gr\"onwall's inequality give $\sup_t\norm{r(t)}^2_{L^2}\le e^{(2K_f+1)T}\int_0^T\norm{m(t)}^2_{H^{-1}}dt$, which replaces (109) with $a=0$ in the proof of Theorem~7 of Nickl; that proof then bounds the integral by a multiple of $\norm h^4_{L^2}$ using only a Taylor expansion of $f$, (18) there, the embedding $H^1\subset L^4$, and (99) and (100) there, whose constants depend only on $f$ and $T$.

(iii) Let $A_\vartheta=\I^*_\vartheta\I_\vartheta$, inverted as in Step 1 of the proof of Proposition~\ref{m:prop:sumrate}, so that by \eqref{m:eq:infobound} for the radius $\max(B,\norm{\theta_0}_{H^{\bar\gamma}})$, $A_\vartheta^{-1}$ maps $H^1$ to $H^{-1}$ with norm $O(1)$. Condition (I2) of Lemma~\ref{m:lem:rdlin}, with the equivalence of norms in its proof, bounds $\I_\vartheta-\I_{\theta_0}$ from $H^{-1}$ to $L^2(\mathcal X)$ by a multiple of $\norm{\vartheta-\theta_0}_{H^\zeta}$, hence also $A_\vartheta-A_{\theta_0}$ from $H^{-1}$ to $H^1$, and $A_\vartheta^{-1}-A_{\theta_0}^{-1}=A_\vartheta^{-1}(A_{\theta_0}-A_\vartheta)A_{\theta_0}^{-1}$ gives $\norm{\bar e_j^{\,\vartheta}-\bar e_j^{\,\theta_0}}_{H^{-1}}\lesssim\norm{\vartheta-\theta_0}_{H^\zeta}\norm{e_j}_{H^1}$. Splitting the difference in \eqref{m:eq:lipweight} as $(\I_\vartheta-\I_{\theta_0})p_N\bar e^{\,\vartheta}_j+\I_{\theta_0}p_N(\bar e^{\,\vartheta}_j-\bar e^{\,\theta_0}_j)$ gives \eqref{m:eq:lipweight}, since $\norm{e_j}_{H^1}\le\norm{e_j}_{H^{\kappa+2}}$ for $\kappa\ge-1$.
\end{proof}

\begin{proposition}[The summary in Sobolev norms: reaction diffusion]\label{m:prop:sumrd}
In the setting of Lemma~\ref{m:lem:rdinfo} let $\bar\gamma<\gamma-1$. Then for every choice of the preliminary estimates that Definition~\ref{m:def:summary} allows, $\norm{\hat S^{J_N}_N-\theta_0}_{L^2}=o_P(\tilde\delta_N(0))$ and $\norm{\hat S^{J_N}_N-\theta_0}_{H^{\bar\gamma}}=o_P(1)$, which is the hypothesis of Theorem~\ref{m:thm:qnnA} on the summary.
\end{proposition}

\begin{proof}
Steps 1, 2 and 4 of the proof of Proposition~\ref{m:prop:sumrate} and its bound on the truncation term do not use \eqref{m:eq:pilotrate}, and \eqref{m:eq:infobound} holds by Remark~\ref{m:rem:infobound}; we keep their notation and replace Step 3. Let $A=\I^*_{\bar\theta}\I_{\bar\theta}$ and $Q=\Id-p_N$. Since $\Lcal^*\Lcal$ is diagonal in $(e_j)$, $[A,p_N]=Q\Kcal_{\bar\theta}p_N-p_N\Kcal_{\bar\theta}Q$, so on $E_B$ Lemma~\ref{m:lem:rdinfo}(i), with $\norm{Q\phi}_{H^1}\le(1+\lambda_J)^{-1}\norm\phi_{H^3}$ and $\norm{Qh}_{H^{-3}}\le(1+\lambda_J)^{-1}\norm h_{H^{-1}}$, gives $\norm{[A,p_N]h}_{H^1}\le2C(1+\lambda_J)^{-1}\norm h_{H^{-1}}$, and Lemma~\ref{m:lem:rdinfo}(ii) at $\bar\theta$ with increment $-h$ gives $\norm R_{L^2(\mathcal X)}\le C\norm h^2_{L^2}$. The identity of Remark~\ref{m:rem:sumopen}, Step 1 and the first elementary inequality bound $p_Nh+\beta(g)$ in $H^s$ on $E_B$ by $c_B^{-2}C(1+\lambda_J)^{(s+1)/2}(2(1+\lambda_J)^{-1}\norm h_{H^{-1}}+C_B\norm h^2_{L^2})$, where $\norm h_{H^{-1}}\le\norm h_{L^2}=O_P(\tilde\delta_N(0))$. With Steps 4 and 5 there, $\norm{\hat S^{J_N}_N-\theta_0}_{H^s}$ is therefore $O_P$ of $N^{(s+1-\gamma)/(2\gamma+d)}+(1+\lambda_J)^{(s-1)/2}\tilde\delta_N(0)+(1+\lambda_J)^{(s+1)/2}\tilde\delta_N(0)^2$ for $-1\le s\le\gamma$. Since $1+\lambda_J\asymp N^{2/(2\gamma+d)}$ and $\tilde\delta_N(0)=N^{-\gamma\bar\gamma/((\bar\gamma+1)(2\gamma+d))}$, the last two terms are $o(\tilde\delta_N(0))$ at $s=0$ and $o(1)$ at $s=\bar\gamma$, because $\gamma>\bar\gamma+1$ and $\bar\gamma\ge1$, and Step 5 there treats the first term.
\end{proof}

\begin{proposition}[The regularized summary on $J_N$ modes: heat and Stokes]\label{m:prop:regheat}
Let $\Gcal$ be the solution map of the heat equation $\partial_tu=\nu\Delta u$ on $\T^d$, with $W=\R$, or of the Stokes equations on $\T^2$, with $W=\R^2$ and $d=2$, and let $(e_j)$ be a real orthonormal eigenbasis of the mean zero functions, respectively of the divergence free mean zero fields, with $\sup_x|e_j(x)|_W\le\sqrt2$, as for the trigonometric basis; thus $s_j(t)=e^{-\nu\lambda_jt}$ and $\lambda_j\asymp j^{2/d}$. Let the data follow \eqref{m:eq:model} at $\theta_0\in H^\gamma$ in the closed span of $(e_j)$, where $\gamma>d+1$ is the exponent of the prior \eqref{m:eq:prior}, and let $k>d/2$. Retain $D=D_N$ modes, give their coefficients the truncation $N(0,\Sigma_0)$, $\Sigma_0=\diag(\sigma^2_{0,j})_{j\le D}$, $\sigma^2_{0,j}=\lambda_j^{-\gamma}/(N\delta_N^2)$, of \eqref{m:eq:prior}, and write $\Pi^D(\cdot\mid Z^{(N)})$ for the resulting posterior. Suppose that
\[
D_N\to\infty,\qquad\delta_N^2D_N^{2+2/d}\to0,\qquad ND_N^{-2(k+2+\gamma)/d}\to0 ,
\]
which holds for $D_N\asymp J_N$. Let $\hat S^{\rm reg}_N$ be the regularized summary \eqref{m:eq:regsummary} computed on the $D$ retained modes from a pilot $\bar\theta_N$ in their span, which may depend on the data in any way. Then, with $\hat\theta_N$ the centering \eqref{m:eq:thetahat} and $\Wone_1$ the Wasserstein distance of order one on $H^{-k-2}$:
\begin{enumerate}[label=(\roman*),nosep]
\item $\sqrt N\norm{\hat S^{\rm reg}_N-\hat\theta_N}_{H^{-k-2}}\to^P0$;
\item $\Wone_1\big(\Law(\sqrt N(\theta-\hat S^{\rm reg}_N)\mid Z^{(N)}),\Ncal_{\theta_0}\big)\to^P0$ for $\theta\sim\Pi^D(\cdot\mid Z^{(N)})$;
\item $\sqrt N(\hat S^{\rm reg}_N-\theta_0)\to^d\Ncal_{\theta_0}$ in $H^{-k-2}$;
\item for the heat equation with integer $\gamma>3d+2$ and $k>3d$, Theorem~\ref{m:thm:nickl2} applies, (ii) holds also for $\theta$ drawn from the posterior $\PiN$ under the untruncated prior \eqref{m:eq:prior}, and the mean $\tilde\theta_N$ of $\PiN$ satisfies $\sqrt N\norm{\hat S^{\rm reg}_N-\tilde\theta_N}_{H^{-k-2}}\to^P0$;
\item for the heat equation, whenever the hypotheses of Proposition~\ref{m:prop:suff} hold, $\sqrt N\norm{\hat S^{\rm reg}_N-\hat S^{J_N}_N}_{H^{-k-2}}\to^P0$.
\end{enumerate}
\end{proposition}

\begin{proof}
For $D_N\asymp J_N\asymp N\delta_N^2=N^{d/(2\gamma+d)}$ the growth conditions hold: $\delta_N^2D_N^{2+2/d}\asymp N^{2(d+1-\gamma)/(2\gamma+d)}\to0$ because $\gamma>d+1$, and $ND_N^{-2(k+2+\gamma)/d}\asymp N^{(d-2k-4)/(2\gamma+d)}\to0$ because $k>d/2$. Write $\Lambda=\diag(\lambda_j)_{j\le D}$, $\mathcal D=\diag(1/\iota_j)_{j\le D}$ as in Section~\ref{m:sec:qnn}, $\theta_0^D=(\theta_{0,j})_{j\le D}$ and $\theta_0^{>D}=\sum_{j>D}\theta_{0,j}e_j$, and identify $a\in\R^D$ with $\sum_{j\le D}a_je_j$. For $i\le N$ let $\Phi_i:\R^D\to W$, $\Phi_ia=\sum_{j\le D}a_js_j(t_i)e_j(x_i)$, be the $i$th row of the design matrix \eqref{m:eq:design} (a pair of rows for Stokes), so that $\Phi^\top\Phi=\sum_i\Phi_i^*\Phi_i$ and $\Phi^\top v=\sum_i\Phi_i^*v_i$. By \eqref{m:eq:iota} with $T_0=0$, $c_1/\lambda_j\le\iota_j\le\min\{1,(2\nu T\lambda_j)^{-1}\}$ with $c_1=(1-e^{-2\nu\lambda_1T})/(2\nu T)$, so that $\kappa_j=N\iota_j\sigma^2_{0,j}=\iota_j\lambda_j^{-\gamma}\delta_N^{-2}\ge c_1\lambda_j^{-\gamma-1}\delta_N^{-2}$. Since $(1+\lambda_j)^{-k-2}\le\lambda_j^{-k-2}$ and $\lambda_j^{-k-2}\iota_j^{-1}\le c_1^{-1}\lambda_j^{-k-1}$, for every $b\in\R^D$
\[
\begin{gathered}
\norm{\mathcal D^{1/2}b}^2_{H^{-k-2}}\le\big\|\Lambda^{-(k+2)/2}\mathcal D^{1/2}b\big\|^2\le c_1^{-1}\lambda_1^{-k-1}\norm b^2,\\
\big\|\Lambda^{-(k+2)/2}\mathcal D^{1/2}\big\|^2_{HS}\le c_1^{-1}\sum_{j\ge1}\lambda_j^{-k-1}=:C_k<\infty,
\end{gathered}
\]
the last because $k+1>d/2$. The constants $C$ below depend on $d$, $\nu$, $T$, $\gamma$, $k$ and $\theta_0$ only.

\emph{Step 0: the summary is the truncated posterior mean.} For the heat and Stokes equations $\I_{\bar\theta_N}$ does not depend on $\bar\theta_N$, so on the retained modes $F_N=\Phi^\top\Phi$ and $\I^*_{\bar\theta_N}[\sum_ir_i\delta_{X_i}]=\Phi^\top(Y-\Phi\bar\theta_N)$, and \eqref{m:eq:regsummary} is the statistic $S(Y)$ of Lemma~\ref{m:lem:exact} with $\varsigma=1$ and $\Sigma_N=V_N=(\Phi^\top\Phi+\Sigma_0^{-1})^{-1}$. The cancellation in the proof of that lemma uses only $V_N(\Phi^\top\Phi+\Sigma_0^{-1})=\Id$ and not the form of $Y$, so $\hat S^{\rm reg}_N=V_N\Phi^\top Y$ for every data vector, although here $Y=\Phi\theta_0^D+u^{>D}+\varepsilon$ with $u^{>D}=(u_{\theta_0^{>D}}(X_i))_{i\le N}$ is not generated by the $D$ mode model. By \eqref{m:eq:linpost}, $\Pi^D(\cdot\mid Z^{(N)})=N(\hat S^{\rm reg}_N,V_N)$. No cross fitting enters.

\emph{Step 1: whitening and the Gram matrix.} Let $\Gamma=N^{-1}\sum_{i\le N}\big(\mathcal D^{1/2}\Phi_i^*\Phi_i\mathcal D^{1/2}-\Id\big)$ and $M=\Id+\diag(\kappa_j^{-1})$. Since $N\iota_j/\kappa_j=\sigma_{0,j}^{-2}$,
\[
\begin{gathered}
\Phi^\top\Phi+\Sigma_0^{-1}=N\mathcal D^{-1/2}(M+\Gamma)\mathcal D^{-1/2},\\
V_N=N^{-1}\mathcal D^{1/2}(M+\Gamma)^{-1}\mathcal D^{1/2}.
\end{gathered}
\]
The matrices $\mathcal D^{1/2}\Phi_i^*\Phi_i\mathcal D^{1/2}$ are independent and identically distributed with mean $\Id$, by the orthonormality of $(e_j)$ and the definition \eqref{m:eq:gram} of $\iota_j$, so $\Gamma$ is an average of $N$ independent centered matrices. Let $\varrho(X_i)=\norm{\Phi_i\mathcal D^{1/2}}^2_{HS}=\sum_{j\le D}\iota_j^{-1}s_j(t_i)^2|e_j(x_i)|^2_W$. Then $E\varrho(X_i)=D$, $\varrho(X_i)\le2\sum_{j\le D}\iota_j^{-1}\le2c_1^{-1}D\lambda_D$ and $\norm{\mathcal D^{1/2}\Phi_i^*\Phi_i\mathcal D^{1/2}}_{HS}\le\varrho(X_i)$, hence
\[
E\norm\Gamma^2_{HS}\le\frac1NE\varrho(X_1)^2\le\frac{2D^2\lambda_D}{c_1N}\le C\,\frac{D^{2+2/d}}N .
\]
The squared length $\varrho$ of a whitened row has mean $D$ but is of order $D\lambda_D$ for observations near $t=0$, before any mode has decayed; this is why the bound is of order $D^{2+2/d}/N$ and not $D^2/N$. It tends to zero, since $D^{2+2/d}/N=\delta_N^2D^{2+2/d}/(N\delta_N^2)$ and $N\delta_N^2\to\infty$. On the event $E_N=\{\norm\Gamma_{op}\le\tfrac12\}$, whose probability tends to one by Chebyshev's inequality, $M+\Gamma\succeq\tfrac12\Id$ because $M\succeq\Id$, so $\norm{(M+\Gamma)^{-1}}_{op}\le2$. It suffices to prove each convergence in probability on $E_N$.

\emph{Step 2: decomposition.} For every $j\ge1$ let $\xi_j=(N\iota_j)^{-1}\sum_{i\le N}s_j(t_i)\langle e_j(x_i),\varepsilon_i\rangle_W$, and let $\xi^D=(\xi_j)_{j\le D}$. For the heat and Stokes equations $\bar e_j=e_j/\iota_j$, and $\I_{\theta_0}p_N(\bar e_j)(t,x)=\iota_j^{-1}s_j(t)e_j(x)$ for $j\le J_N$ and $0$ otherwise, so \eqref{m:eq:thetahat} reads $\hat\theta_N=\theta_0+\sum_{j\le J_N}\xi_je_j$. Since $V_N\Phi^\top\Phi=\Id-V_N\Sigma_0^{-1}$, Step 0 gives
\[
\hat S^{\rm reg}_N-\hat\theta_N=A_1+A_2+A_3+A_4-A_5-A_6,
\]
with $A_1=-V_N\Sigma_0^{-1}\theta_0^D$ (the prior shrinkage), $A_2=V_N\Phi^\top u^{>D}$ (the aliasing of the discarded modes), $A_3=V_N\Phi^\top\varepsilon-\xi^D$ (the noise), $A_4=\sum_{J_N<j\le D}\xi_je_j$, $A_5=\sum_{D<j\le J_N}\xi_je_j$ and $A_6=\theta_0^{>D}$; at most one of $A_4$ and $A_5$ is nonzero. We show that $\sqrt N\norm{A_m}_{H^{-k-2}}\to^P0$ for each $m$.

\emph{Step 3: the prior shrinkage.} By Step 1, $V_N\Sigma_0^{-1}=\mathcal D^{1/2}(M+\Gamma)^{-1}(M-\Id)\mathcal D^{-1/2}$. The resolvent identity $(M+\Gamma)^{-1}=M^{-1}-(M+\Gamma)^{-1}\Gamma M^{-1}$ and $M^{-1}(M-\Id)=\diag((1+\kappa_j)^{-1})$ split $A_1=A_1'+A_1''$, with
\[
\begin{gathered}
A_1'=-\big(\theta_{0,j}/(1+\kappa_j)\big)_{j\le D},\qquad A_1''=\mathcal D^{1/2}(M+\Gamma)^{-1}\Gamma w,\\
w=\big(\iota_j^{1/2}\theta_{0,j}/(1+\kappa_j)\big)_{j\le D}.
\end{gathered}
\]
The diagonal part $A_1'$ is of order $\theta_{0,j}$ on the modes with $\kappa_j\le1$, as Remark~\ref{m:rem:reg} anticipates, and the weights of $H^{-k-2}$ make it negligible. Let $s=\min\{2,(k+2+\gamma)/(\gamma+1)\}$. Since $(1+\kappa)^{-2}\le\kappa^{-s}$ for $\kappa>0$ and $0\le s\le2$, and since $s(\gamma+1)-k-2\le\gamma$,
\[
\begin{aligned}
N\norm{A_1'}^2_{H^{-k-2}}&\le N\sum_{j\le D}\lambda_j^{-k-2}\kappa_j^{-s}\theta_{0,j}^2\\
&\le c_1^{-s}N\delta_N^{2s}\sum_j\lambda_j^{s(\gamma+1)-k-2}\theta^2_{0,j}\le C\,N\delta_N^{2s}\norm{\theta_0}^2_{H^\gamma},
\end{aligned}
\]
and $N\delta_N^{2s}=N^{1-2s\gamma/(2\gamma+d)}\to0$ because $2s\gamma>2\gamma+d$: for $s=2$ since $2\gamma>d$, and otherwise since $2\gamma(k+1)>\gamma d+2\gamma>d(\gamma+1)$, by $k>d/2$ and $2\gamma>d$. For the cross part, $(1+\kappa_j)^2\ge4\kappa_j$ gives $\norm w^2\le\sum_j\iota_j\theta^2_{0,j}/(4\kappa_j)=\tfrac14\delta_N^2\sum_j\lambda_j^\gamma\theta_{0,j}^2\le\tfrac14\delta_N^2\norm{\theta_0}^2_{H^\gamma}$, so on $E_N$
\[
\begin{aligned}
N\norm{A_1''}^2_{H^{-k-2}}&\le c_1^{-1}\lambda_1^{-k-1}N\norm{(M+\Gamma)^{-1}\Gamma w}^2\\
&\le c_1^{-1}\lambda_1^{-k-1}\norm{\theta_0}^2_{H^\gamma}\,N\delta_N^2\norm\Gamma^2_{HS},
\end{aligned}
\]
whose expectation is at most $C\delta_N^2D^{2+2/d}\to0$ by Step 1. This is the only place where the full condition $\delta_N^2D_N^{2+2/d}\to0$ is used, through the concentration of the whitened realized information on the population one; elsewhere its consequence $D^{2+2/d}/N\to0$ suffices. Proposition~\ref{m:prop:regfixed} below shows that neither is needed.

\emph{Step 4: aliasing.} $A_2=N^{-1}\mathcal D^{1/2}(M+\Gamma)^{-1}g$ with $g=\sum_i\mathcal D^{1/2}\Phi_i^*u_{\theta_0^{>D}}(X_i)$, a sum of independent centered vectors: for $j\le D$ the $j$th coordinate of a summand has mean $\iota_j^{-1/2}\sum_{l>D}\theta_{0,l}E\big[s_j(t_1)s_l(t_1)\langle e_j(x_1),e_l(x_1)\rangle_W\big]=0$, the series converging absolutely because $\sum_l|\theta_{0,l}|\le(\sum_l\lambda_l^{-\gamma})^{1/2}\norm{\theta_0}_{H^\gamma}<\infty$. Hence, by Step 1 and $\iota_l\le(2\nu T\lambda_l)^{-1}$,
\[
\begin{aligned}
E\norm g^2&=N\,E\norm{\mathcal D^{1/2}\Phi_1^*u_{\theta_0^{>D}}(X_1)}^2\le N\sup\varrho\;E|u_{\theta_0^{>D}}(X_1)|_W^2\\
&\le\frac{2ND\lambda_D}{c_1}\sum_{l>D}\iota_l\theta_{0,l}^2\le C\,ND\lambda_D^{-\gamma}\norm{\theta_0}^2_{H^\gamma},
\end{aligned}
\]
and on $E_N$, $N\norm{A_2}^2_{H^{-k-2}}\le4c_1^{-1}\lambda_1^{-k-1}N^{-1}\norm g^2$, whose expectation is at most $CD\lambda_D^{-\gamma}\to0$, since $\lambda_D\asymp D^{2/d}$ and $\gamma>d/2$.

\emph{Step 5: the noise.} Let $\chi=N^{-1/2}\mathcal D^{1/2}\Phi^\top\varepsilon$; given the design, $\chi\sim N(0,\Id+\Gamma)$. By Step 1, $V_N\Phi^\top\varepsilon=N^{-1/2}\mathcal D^{1/2}(M+\Gamma)^{-1}\chi$ and $\xi^D=N^{-1}\mathcal D\Phi^\top\varepsilon=N^{-1/2}\mathcal D^{1/2}\chi$, so $\sqrt NA_3=\mathcal D^{1/2}Q\chi$ with
\[
Q=(M+\Gamma)^{-1}-\Id=-\diag\big((1+\kappa_j)^{-1}\big)-(M+\Gamma)^{-1}\Gamma M^{-1}.
\]
Given the design, on $E_N$,
\[
\begin{aligned}
E\big[N\norm{A_3}^2_{H^{-k-2}}\,\big|\,(X_i)\big]&\le\norm{\Id+\Gamma}_{op}\big\|\Lambda^{-(k+2)/2}\mathcal D^{1/2}Q\big\|^2_{HS}\\
&\le3\big(a_N+4c_1^{-1}\lambda_1^{-k-1}\norm\Gamma^2_{HS}\big),\\
a_N&=\sum_{j\le D}\frac{\lambda_j^{-k-2}\iota_j^{-1}}{(1+\kappa_j)^2}.
\end{aligned}
\]
For fixed $j$ the summand of $a_N$ tends to zero, since $\kappa_j\to\infty$ as $\delta_N\to0$, and it is at most $c_1^{-1}\lambda_j^{-k-1}$, which is summable; so $a_N\to0$ by dominated convergence, and $N\norm{A_3}^2_{H^{-k-2}}\to^P0$ by Markov's inequality given the design and Step 1.

\emph{Step 6: the tails.} $E[N\xi_j^2]=\iota_j^{-1}$, since the summands that define $\xi_j$ are independent and centered with second moment $E[s_j(t_1)^2|e_j(x_1)|_W^2]=\iota_j$. Hence $E\,N\norm{A_4}^2_{H^{-k-2}}\le c_1^{-1}\sum_{j>J_N}\lambda_j^{-k-1}\to0$ and $E\,N\norm{A_5}^2_{H^{-k-2}}\le c_1^{-1}\sum_{j>D}\lambda_j^{-k-1}\to0$, while $N\norm{A_6}^2_{H^{-k-2}}\le N\lambda_D^{-k-2-\gamma}\norm{\theta_0}^2_{H^\gamma}\le CND^{-2(k+2+\gamma)/d}\to0$; this is where the growth of $D$ is limited from below. With Steps 3 to 5 this proves (i).

\emph{Step 7: the truncated posterior.} By Steps 0 and 1, $\sqrt N(\theta-\hat S^{\rm reg}_N)\sim N(0,\mathcal D^{1/2}(M+\Gamma)^{-1}\mathcal D^{1/2})$ under $\Pi^D(\cdot\mid Z^{(N)})$. Let $(Z_j)_{j\ge1}$ be independent standard normal variables, independent of the data, and $Z^D=(Z_j)_{j\le D}$. Then $\mathcal D^{1/2}(M+\Gamma)^{-1/2}Z^D$ has that law, $\sum_j\iota_j^{-1/2}Z_je_j$ has the law $\Ncal_{\theta_0}=N(0,\diag(1/\iota_j))$, and this coupling gives
\[
\begin{aligned}
&\Wone_1\big(\Law(\sqrt N(\theta-\hat S^{\rm reg}_N)\mid Z^{(N)}),\Ncal_{\theta_0}\big)\\
&\quad\le\big\|\Lambda^{-(k+2)/2}\mathcal D^{1/2}(M^{-1/2}-\Id)\big\|_{HS}\\
&\qquad+C_k^{1/2}\norm{(M+\Gamma)^{-1/2}-M^{-1/2}}_{op}+\Big(\sum_{j>D}\frac{\lambda_j^{-k-2}}{\iota_j}\Big)^{1/2}.
\end{aligned}
\]
The first term is at most $a_N^{1/2}$, because $1-(\kappa_j/(1+\kappa_j))^{1/2}\le1-\kappa_j/(1+\kappa_j)=(1+\kappa_j)^{-1}$, and the third tends to zero. For the second we use that symmetric matrices $A,B\succeq a\Id$, $a>0$, satisfy $\norm{A^{-1/2}-B^{-1/2}}_{op}\le\norm{A-B}_{op}/(2a^{3/2})$. Indeed $X=A^{1/2}-B^{1/2}$ solves $A^{1/2}X+XB^{1/2}=A-B$; at a unit eigenvector $x$ of $X$ with eigenvalue $\mu$ this gives $|\mu|\,\langle x,(A^{1/2}+B^{1/2})x\rangle\le\norm{A-B}_{op}$ with $\langle x,(A^{1/2}+B^{1/2})x\rangle\ge2a^{1/2}$, so $\norm X_{op}\le\norm{A-B}_{op}/(2a^{1/2})$, and $A^{-1/2}-B^{-1/2}=A^{-1/2}(B^{1/2}-A^{1/2})B^{-1/2}$. With $A=M+\Gamma$, $B=M$ and $a=\tfrac12$ on $E_N$, the second term is at most $(2C_k)^{1/2}\norm\Gamma_{op}\to^P0$. This proves (ii).

\emph{Step 8: the limit law of the summary.} Let $\mathcal D_J$, $\Lambda_J$ and $\Gamma_J$ be $\mathcal D$, $\Lambda$ and $\Gamma$ formed with $J_N$ in place of $D$. Given the design, $\sqrt N\mathcal D_J^{-1/2}(\xi_j)_{j\le J_N}\sim N(0,\Id+\Gamma_J)$, so $\sqrt N(\hat\theta_N-\theta_0)=\sum_{j\le J_N}\sqrt N\xi_je_j$ has the law of $\mathcal D_J^{1/2}(\Id+\Gamma_J)^{1/2}(Z_j)_{j\le J_N}$, and the coupling of Step 7, with $|\sqrt{1+\mu}-1|\le|\mu|$ for $\mu\ge-1$, gives
\[
\begin{aligned}
&E\Big[\Big\|\mathcal D_J^{1/2}(\Id+\Gamma_J)^{1/2}(Z_j)_{j\le J_N}-\sum_{j\ge1}\iota_j^{-1/2}Z_je_j\Big\|^2_{H^{-k-2}}\,\Big|\,(X_i)\Big]\\
&\qquad\le C_k\norm{\Gamma_J}^2_{op}+\sum_{j>J_N}\frac{\lambda_j^{-k-2}}{\iota_j}.
\end{aligned}
\]
This tends to zero in probability, because $E\norm{\Gamma_J}^2_{HS}\le CJ_N^{2+2/d}/N\asymp N^{(d+2-2\gamma)/(2\gamma+d)}\to0$ by Step 1. Hence $\sqrt N(\hat\theta_N-\theta_0)\to^d\Ncal_{\theta_0}$ in $H^{-k-2}$, and (iii) follows from (i).

\emph{Step 9: the untruncated posterior.} For the heat equation Conditions 1 and 2 of Nickl~\cite{nicklbvm2024} hold with $\eta_0=0$ and every $\zeta\in(d/2,\gamma-d/2)$. Condition 1 is \eqref{m:eq:prior}, whose RKHS is $H^\gamma_0$, with $\gamma>1+d/2$. In Condition 2, $\Gcal$ is linear, so $\I_{\theta_0}=\Gcal$ and the remainder in 2C vanishes; 2A, 2B and 2D follow from $|u_h(t,x)|\le\sqrt2\sum_j|h_j|\le\sqrt2(\sum_j\lambda_j^{-r})^{1/2}\norm h_{H^r}$ for $r>d/2$, used with $r=\zeta$ in 2D and with $r=\bar\gamma$ in 2A, where $\bar\gamma$ is the exponent of Condition 2, which the proof of Nickl's theorem fixes in (31) there with $\bar\gamma>3\zeta+d+2>d/2$, and from $\norm{u_h}^2_{L^2(\mathcal X)}=T\sum_j\iota_jh_j^2\le T\norm h^2_{L^2}$; 2E follows from $T\sum_j\iota_jh_j^2\ge c_1T\norm h^2_{H^{-1}}$; and 2F holds because $\Delta\I^*_{\theta_0}\I_{\theta_0}$ multiplies $e_j$ by $-T\lambda_j\iota_j=-(1-e^{-2\nu\lambda_jT})/(2\nu)$, which lies between $-1/(2\nu)$ and $-c_1T$, so that it is a homeomorphism of every $H^\eta_0$. If $\gamma>3d+2$ is an integer and $k>3d$, some $\zeta>d/2$ satisfies $\gamma>3\zeta+3d/2+2$ and $k>\zeta+5d/2$, and Theorem~\ref{m:thm:nickl2} applies. Nickl forms $\hat\theta_N$ with the projection onto the span of the first $J_N$ trigonometric polynomials, (32) and (36) there; we choose $J_N\asymp N\delta_N^2$ so that the first $J_N$ modes $(e_j)$ fill every eigenspace of $\Delta$ that they meet, as (35) there allows, and then the two spans agree and his $\hat\theta_N$ is \eqref{m:eq:thetahat}. The proof of Theorem~\ref{m:thm:nickl2} gives, in (80) of Nickl~\cite{nicklbvm2024}, $\Wone_1\big(\Law(\sqrt N(\theta-\hat\theta_N)\mid Z^{(N)}),\Ncal_{\theta_0}\big)\to^P0$ for $\theta\sim\PiN$, and, in the display after (80), $\sqrt N\norm{\tilde\theta_N-\hat\theta_N}_{H^{-k-2}}\to^P0$. Since a shift by $c$ moves a law by at most $\norm c_{H^{-k-2}}$ in $\Wone_1$, (i) gives (iv).

\emph{Step 10: the one step summary.} Under the hypotheses of Proposition~\ref{m:prop:suff}, $\sqrt N\norm{\hat S^{J_N}_N-\hat\theta_N}_{H^{-k-2}}\to^P0$, and (v) follows from (i) by the triangle inequality. For the heat equation with integer $\gamma>3d+2$ and $k>3d$ these hypotheses reduce, by Step 9 and since \eqref{m:eq:lipweight} holds with $L=0$ and $\I^*_{\theta_0}\I_{\theta_0}$ is diagonal in $(e_j)$, to those on the pilot estimates in Definition~\ref{m:def:summary} and Proposition~\ref{m:prop:suff}. For the Stokes equations the prior and the forward map live on the divergence free fields, while Conditions 1 and 2 are stated on all of $H^\gamma_0(\Omega,W)$; they, and with them the proof of Proposition~\ref{m:prop:suff}, would have to be read on the divergence free subspace, which Remark~1 of Nickl~\cite{nicklbvm2024} suggests without a stated theorem. This is the step that (v) and the first claim of (iv) would need; Proposition~\ref{m:prop:regfixed} below gives the other two without it.
\end{proof}

Proposition~\ref{m:prop:regheat} compares the regularized summary with the centering \eqref{m:eq:thetahat}, and not directly with the one step statistic or with the posterior mean of Theorem~\ref{m:thm:nickl2}, because neither shortcut is available. Even for the heat equation the projected one step statistic $\hat S^{J_N}_N$ depends on its pilot estimates, through the deviation of the realized Gram matrix on each half sample from its mean, so it is neither $\hat\theta_N$ nor $p_N\hat\theta_N$; this is why (v) needs the hypotheses of Proposition~\ref{m:prop:suff}. And $\hat S^{\rm reg}_N$ is the posterior mean under the truncated prior, not the mean $\tilde\theta_N$ under the untruncated prior to which Theorem~\ref{m:thm:nickl2} refers; (iv) connects the two through $\hat\theta_N$.

\begin{proposition}[The regularized summary under a fixed prior, and the untruncated posteriors]\label{m:prop:regfixed}
Let $\Gcal$, $W$, $(e_j)$ and the data be as in Proposition~\ref{m:prop:regheat}, but with $\theta_0\in H^\beta$, $\beta>d/2$, and $k>d/2-1$, and give the $D=D_N$ retained coefficients the fixed prior $\sigma^2_{0,j}=\rho^2\lambda_j^{-\alpha}$ of Section~\ref{m:sec:heat}, $\rho>0$, $\alpha>d/2$. Suppose that $\beta\ge\alpha-k-1$, that $D_N\to\infty$ with $\sup_NND_N^{-2(k+2+\beta)/d}<\infty$, a condition void when $\theta_0$ has finitely many nonzero coefficients, and that $\alpha>d$ or $D_N^{2+2/d}/N\to0$, and form \eqref{m:eq:thetahat} with the first $J_N$ modes, for any $J_N\to\infty$. Then (i) to (iii) of Proposition~\ref{m:prop:regheat} hold, and (v) for the heat equation. If $\alpha>d$, they hold also for $D=\infty$, that is, under the untruncated prior. The conclusions of Proposition~\ref{m:prop:regheat} hold without its condition $\delta_N^2D_N^{2+2/d}\to0$, and (i) to (iii) there also for $D=\infty$, that is, for $\PiN$ and its mean $\tilde\theta_N$. Hence the last two claims of (iv) there hold for the Stokes equations as well, and for the heat equation without its restrictions on $\gamma$ and $k$; and when $\theta_0$ meets both sets of hypotheses, the regularized summaries under the two priors, and if $\alpha>d$ the untruncated posterior means, are $\sqrt N$ equivalent in $H^{-k-2}$.
\end{proposition}

\begin{proof}
We keep the notation of the proof of Proposition~\ref{m:prop:regheat}, with $\kappa_j=N\iota_j\sigma^2_{0,j}$ and with constants $C$ that depend on $\alpha$ and $\beta$ in place of $\gamma$, and not on $N$, $D$ or $\rho$. Steps 0, 2, 6, 8 and 10 stand, with $N\norm{A_6}^2_{H^{-k-2}}\le N\lambda_D^{-k-2-\beta}\sum_{j>D}\lambda_j^\beta\theta^2_{0,j}\to0$ in Step 6, and with Step 8 applied to some $J'_N\to\infty$ with $J'^{2+2/d}_N/N\to0$, which by (i) gives a $\sqrt N$ equivalent centering. We whiten instead by the mean $N\mathcal D^{-1/2}M\mathcal D^{-1/2}=N\diag(\iota_j)+\Sigma_0^{-1}$ of the posterior precision: let $\tilde\Gamma=M^{-1/2}\Gamma M^{-1/2}$ and $B_i=\Phi_i\mathcal D^{1/2}M^{-1/2}$, so that $V_N=N^{-1}\mathcal D^{1/2}M^{-1/2}(\Id+\tilde\Gamma)^{-1}M^{-1/2}\mathcal D^{1/2}$ and $\tilde\Gamma=N^{-1}\sum_i(B_i^*B_i-M^{-1})$, an average of independent centered matrices. Since $\iota_j^{-1}\kappa_j=N\sigma^2_{0,j}$, $\norm{B_i}^2_{HS}$ has mean $m_N=\sum_{j\le D}\kappa_j/(1+\kappa_j)$ and is at most $2\bar m_N$, $\bar m_N=\sum_{j\le D}N\sigma^2_{0,j}/(1+\kappa_j)$, so as in Step 1, and with $EB_i^*B_i=M^{-1}$ for the second bound, for $v\in\R^D$
\[
E\norm{\tilde\Gamma}^2_{HS}\le2\bar m_Nm_N/N,\qquad E\norm{\tilde\Gamma v}^2\le2\bar m_N\norm{M^{-1/2}v}^2/N .
\]
Let $L=(N\rho^2)^{d/(2\alpha+2)}$, the order of the effective dimension of Theorem~\ref{m:thm:jeff}. As $\lambda_j\asymp j^{2/d}$ and $c_1\le\lambda_j\iota_j\le(2\nu T)^{-1}$, $\kappa_j\asymp(j/L)^{-(2\alpha+2)/d}$, and bounding the summands of $m_N$ and $\bar m_N$ by $\min\{1,\kappa_j\}$ and $\min\{N\sigma^2_{0,j},\iota_j^{-1}\}$ gives $m_N\le C\min\{D,L\}$ and, as $2\alpha/d>1$, $\bar m_N\le C\min\{D,L\}^{1+2/d}$. So $E\norm{\tilde\Gamma}^2_{HS}\le C\min\{D,L\}^{2+2/d}/N\to0$, if $\alpha>d$ whatever $D$ is, since $L^{2+2/d}/N\asymp N^{(d-\alpha)/(\alpha+1)}$, and otherwise by hypothesis. The event $\tilde E_N=\{\norm{\tilde\Gamma}_{op}\le\frac12\}$ replaces $E_N$, and $\Gamma$ itself need not be small.

In Step 3, by $(1+\kappa_j)^2\ge4\kappa_j$ and $N/\kappa_j\le\lambda_j^{\alpha+1}/(c_1\rho^2)$, the $j$th summand of $N\norm{A_1'}^2_{H^{-k-2}}$ is at most $\lambda_j^{\alpha-k-1}\theta^2_{0,j}/(4c_1\rho^2)$, summable as $\beta\ge\alpha-k-1$, and tends to zero, so $N\norm{A_1'}^2_{H^{-k-2}}\to0$ by dominated convergence. As $(M+\Gamma)^{-1}\Gamma=M^{-1/2}(\Id+\tilde\Gamma)^{-1}\tilde\Gamma M^{1/2}$, the second bound above with $v=M^{1/2}w$ gives $E[N\norm{A_1''}^2_{H^{-k-2}}1_{\tilde E_N}]\le8c_1^{-1}\lambda_1^{-k-1}\bar m_N\norm w^2$, and with $\beta'=\min\{\beta,2\alpha\}$ the bounds $\iota_j\le Cj^{-2/d}$ and $(1+\kappa_j)^{-2}\le C\min\{1,(j/L)^{(4\alpha+4)/d}\}$ give $\norm w^2\le CL^{-(2+2\beta')/d}\norm{\theta_0}^2_{H^{\beta'}}$, since $2+2\beta'<4\alpha+4$, so this is at most $CL^{1-2\beta'/d}\to0$, as $\beta'>d/2$. In Step 4, $A_2=N^{-1}\mathcal D^{1/2}M^{-1/2}(\Id+\tilde\Gamma)^{-1}M^{-1/2}g$, and the bound becomes $CD\lambda_D^{-\beta}$. In Step 5, $Q=(M^{-1}-\Id)-M^{-1/2}(\Id+\tilde\Gamma)^{-1}\tilde\Gamma M^{-1/2}$, the first part contributes at most $a_N$ to $E\,N\norm{A_3}^2_{H^{-k-2}}$ because $E[\Id+\Gamma]=\Id$, and $M^{-1/2}\chi\sim N(0,M^{-1}+\tilde\Gamma)$ given the design, so $E[N\norm{A_3}^2_{H^{-k-2}}1_{\tilde E_N}]\le2a_N+12c_1^{-1}\lambda_1^{-k-1}E\norm{\tilde\Gamma}^2_{HS}$. In Step 7 we couple through $\mathcal D^{1/2}M^{-1/2}(\Id+\tilde\Gamma)^{-1/2}Z^D$, and the middle term is at most $(2C_k)^{1/2}\norm{\tilde\Gamma}_{op}$ on $\tilde E_N$. This proves (i) to (iii) and (v).

If $\alpha>d$, these bounds do not depend on $D$, except those on $A_2$, $A_5$, $A_6$ and the tail in Step 7, which vanish as $D\to\infty$ for fixed $N$. Given the design, $\theta$ and $Y$ are jointly Gaussian under the untruncated prior, since $|s_je_j|_W\le\sqrt2$ and $\sum_j\sigma^2_{0,j}<\infty$. So for $\theta'\sim N(0,\Sigma_0)$ and $\varepsilon'\sim N(0,\Id)$, independent of each other and of the data, $\theta'+K(Y-\Phi\theta'-\varepsilon')$ with $K=\Sigma_0\Phi^*(\Phi\Sigma_0\Phi^*+\Id)^{-1}$ has the untruncated posterior law, and by the Woodbury identity the same holds for $\Pi^D(\cdot\mid Z^{(N)})$ with $\theta'$, $\Sigma_0$ and $\Phi$ restricted to the first $D$ modes. As $D\to\infty$ for fixed data these draws converge in $L^2$, by dominated convergence, so the truncated posteriors and their means converge to the untruncated ones in $\Wone_1$ on $H^{-k-2}$ and in $\ell^2$, and Fatou's lemma, applied as $D\to\infty$ for fixed $N$ to the indicators of the open events in (i) and (ii), gives (i) to (iii) for $D=\infty$. Under \eqref{m:eq:prior}, $\kappa_j=\iota_j\lambda_j^{-\gamma}\delta_N^{-2}\to\infty$, $A_1'$ is bounded in Step 3 of that proof, and the rest holds with $\delta_N^{-2}$ for $N\rho^2$ and $\gamma$ for $\alpha$ and $\beta$, since then $L=\delta_N^{-d/(\gamma+1)}$ and $L^{2+2/d}/N=N^{-(2\gamma^2-\gamma d+d)/((2\gamma+d)(\gamma+1))}\to0$ whatever $D$ is; Steps 9 and 10 then give (iv) and (v) from (i). The last claims follow from (i) and (ii) by the triangle inequality, since a shift by $c$ moves a law by at most $\norm c_{H^{-k-2}}$ in $\Wone_1$.
\end{proof}

\begin{remark}[What a nonlinear version of Proposition~\ref{m:prop:regheat} needs]\label{m:rem:regnonlin}
For a nonlinear forward map let $P_D$ be the projection onto the $D$ retained modes, let the pilot $\bar\theta_N$ lie in their span, let $\Phi$ and $F_N=\Phi^\top\Phi$ be formed with $\I_{\bar\theta_N}$ on those modes, and write $h=\bar\theta_N-\theta_0$ and $\theta_0^D$, $\theta_0^{>D}$, $\Lambda$ as in the proof of Proposition~\ref{m:prop:regheat}. Substituting $Y_i=\Gcal(\theta_0)(X_i)+\varepsilon_i$ in \eqref{m:eq:regsummary} gives, exactly,
\[
\begin{gathered}
\hat S^{\rm reg}_N-\theta_0^D=\big(F_N+\Sigma_0^{-1}\big)^{-1}\Big(\Phi^\top\varepsilon+\Phi^\top q-\Sigma_0^{-1}\theta_0^D\Big),\\
q_i=\big(\I_{\bar\theta_N}\theta_0^{>D}+(\I_{\bar\theta_N}-\I_{\theta_0})h-R_0\big)(X_i),
\end{gathered}
\]
with $R_0=\Gcal(\theta_0+h)-\Gcal(\theta_0)-\I_{\theta_0}h$ the remainder of Condition 2C, so that the pilot enters only through $F_N$, $\Phi$ and $q$. For the heat and Stokes equations $q_i=u_{\theta_0^{>D}}(X_i)$, and the proof of Proposition~\ref{m:prop:regheat} bounds the three terms. For a nonlinear map the same proof needs five things that neither Conditions 1 and 2 nor \eqref{m:eq:lipweight} supply. (a) Since $\Phi$ depends on the data through $\bar\theta_N$, the noise term is not conditionally Gaussian; the step must be cross fitted as in Definition~\ref{m:def:summary}, or the noise term and the Gram matrix bounded uniformly over the localization set. (b) The whitening of Step 1, now by the compression $H_D(\vartheta)=T^{-1}P_D\I^*_\vartheta\I_\vartheta P_D$ of the information to the retained modes, needs $\norm{\I_\vartheta f}_{L^2(\mathcal X)}\ge c\norm f_{H^{-1}}$ uniformly for $\vartheta$ near $\theta_0$, where Conditions 2C and 2E give it at $\vartheta=\theta_0$ only, and $\sup_j\norm{\I_\vartheta e_j}_\infty<\infty$; with both, the squared length of a whitened row is at most a constant times $D\lambda_D$, against a mean of $D$, the excess again coming from observations near $t=0$, and the second moment bound of Step 1 holds. (c) The step inverts the compression $H_D(\vartheta)$ of the information, while the centering \eqref{m:eq:thetahat} uses the compression $p_NT(\I^*_{\theta_0}\I_{\theta_0})^{-1}p_N$ of its inverse. The two agree when the information is diagonal in $(e_j)$; in general the difference of the two noise terms is governed by the difference of these operators, which has to vanish in the weighted norm of $H^{-k-2}$, and this needs a Galerkin form of Condition 2F, for instance $\sup_D\norm{H_D(\vartheta)^{-1}P_D\psi}_{H^\eta}\lesssim\norm\psi_{H^{\eta+2}}$ uniformly for $\vartheta$ near $\theta_0$. With $H_D(\vartheta)\succeq c\Lambda^{-1}$ from (b), the same bound would also control the prior shrinkage $(F_N+\Sigma_0^{-1})^{-1}\Sigma_0^{-1}\theta_0^D$, which is no longer diagonal, by splitting the test function into low and high modes. (d) The term $(\I_{\bar\theta_N}-\I_{\theta_0})h$ needs $\vartheta\mapsto\I_\vartheta$ to be Lipschitz into the bounded operators on $L^2$, which is more than \eqref{m:eq:lipweight}, a condition on the weights of Definition~\ref{m:def:summary} only. (e) The aliasing term $\I_{\bar\theta_N}\theta_0^{>D}$ is no longer orthogonal to the retained modes, and its projection is governed by the coupling of retained and discarded modes in $\I^*_{\bar\theta_N}\I_{\bar\theta_N}$, which for Navier Stokes is the compact part $\Kcal$ of \eqref{m:eq:fredholm}. For the reaction diffusion class of Theorem~\ref{m:thm:nickl1} the bound $\sup_j\norm{\I_\vartheta e_j}_\infty\le\sqrt2\,e^{T\sup|f'|}$ in (b) follows from the maximum principle, and Remark~\ref{m:rem:infobound} gives the first bound in (b) uniformly on balls of $H^{\bar\gamma}$ for that class and for the two dimensional Navier Stokes equations. For that class Proposition~\ref{m:prop:regrd} meets (a) by cross fitting, avoids (d) by linearizing at the pilot, and meets (c) and (e) by Lemma~\ref{m:lem:rdinfo}; for the Navier Stokes equations the rest is open.
\end{remark}

\begin{proposition}[The regularized summary on $J_N$ modes: reaction diffusion]\label{m:prop:regrd}
In the setting of Lemma~\ref{m:lem:rdinfo} let $k>3d$, so that Theorem~5 of Nickl~\cite{nicklbvm2024} gives Conditions 1 and 2 with $\gamma$ and $k$ as in Theorem~\ref{m:thm:nickl2}. Retain $D=J_N$ modes with $\Sigma_0$ as in Proposition~\ref{m:prop:regheat}, let the preliminary estimates $\bar\theta^{[b]}_N$ of Definition~\ref{m:def:summary} lie in the span of $e_1,\dots,e_{J_N}$, as their projections $p_N\bar\theta^{[b]}_N$ do with the same rates, and let $\hat S^{\rm reg}_N=\frac12(\hat S^{{\rm reg},1}_N+\hat S^{{\rm reg},2}_N)$, where $\hat S^{{\rm reg},b}_N$ is \eqref{m:eq:regsummary} computed with the pilot $\bar\theta^{[b]}_N$ and with $F_N$ and the sum over $i$ restricted to $\mathcal I_{3-b}$. Then:
\begin{enumerate}[label=(\roman*),nosep]
\item $\sqrt N\norm{\hat S^{\rm reg}_N-\hat\theta_N}_{H^{-k-2}}\to^P0$;
\item $\Wone_1\big(\Law(\sqrt N(\theta-\hat S^{\rm reg}_N)\mid Z^{(N)}),\Ncal_{\theta_0}\big)\to^P0$ for $\theta\sim\PiN$, the mean $\tilde\theta_N$ of $\PiN$ satisfies $\sqrt N\norm{\hat S^{\rm reg}_N-\tilde\theta_N}_{H^{-k-2}}\to^P0$, and $\sqrt N(\hat S^{\rm reg}_N-\theta_0)\to^d\Ncal_{\theta_0}$ in $H^{-k-2}$;
\item under the hypotheses of Proposition~\ref{m:prop:suff} other than \eqref{m:eq:lipweight}, which Lemma~\ref{m:lem:rdinfo} supplies, $\sqrt N\norm{\hat S^{\rm reg}_N-\hat S^{J_N}_N}_{H^{-k-2}}\to^P0$.
\end{enumerate}
\end{proposition}

\begin{proof}
Since $\norm{\theta_0-p_N\theta_0}_{H^\xi}\le\lambda_{J_N}^{-(\gamma-\xi)/2}\norm{\theta_0}_{H^\gamma}$ and $(\gamma-\xi)(\bar\gamma+1)\ge\gamma(\bar\gamma-\xi)$ for $0\le\xi\le\bar\gamma<\gamma$, the projections keep the rates. Fix $b$; write $\bar\theta$, $\mathcal I$, $n$, $J$, $h$, $E_B$, $c_B$ and $C_B$ as in the proof of Proposition~\ref{m:prop:sumrate}, $C$ for the constant of Lemma~\ref{m:lem:rdinfo}, and $\Lambda$, $\mathcal D$, $\theta_0^D$, $\theta_0^{>D}$, $c_1$, $w$ and $a_N$ as in the proof of Proposition~\ref{m:prop:regheat}, with $\nu=1$, $D=J$ and $\kappa_j=n\iota_j\sigma^2_{0,j}\ge\frac13c_1\lambda_j^{-\gamma-1}\delta_N^{-2}$. For $b'\in\R^J$ let $\norm{b'}_{-k-2}=\norm{\sum_jb'_je_j}_{H^{-k-2}}$, so that $\norm{\mathcal D^{1/2}b'}_{-k-2}\le C_0\norm{b'}$ by the first display of that proof. Let $\Phi_ia=\sum_{j\le J}a_j(\I_{\bar\theta}e_j)(X_i)$ for $i\in\mathcal I$, $R=u_{\bar\theta}-u_{\theta_0}-\I_{\bar\theta}h$ and $V_N=(F_N+\Sigma_0^{-1})^{-1}$. Since $Y_i-u_{\bar\theta}(X_i)=\varepsilon_i-\Phi_i(\bar\theta-\theta_0^D)+(\I_{\bar\theta}\theta_0^{>D}-R)(X_i)$, the rearrangement of Remark~\ref{m:rem:regnonlin}, made at $\bar\theta$, gives exactly $\hat S^{{\rm reg},b}_N-\theta_0^D=V_N(\Phi^\top\varepsilon+\Phi^\top((\I_{\bar\theta}\theta_0^{>D}-R)(X_i))_{i\in\mathcal I}-\Sigma_0^{-1}\theta_0^D)$.

\emph{Whitening.} Let $H=(T^{-1}\langle\I_{\bar\theta}e_j,\I_{\bar\theta}e_l\rangle_{L^2(\mathcal X)})_{j,l\le J}$ and $G=\mathcal D^{1/2}H\mathcal D^{1/2}$; by \eqref{m:eq:modeinfo}, $G=\Id+\tilde\Kcal$ with $\tilde\Kcal=(T^{-1}\iota_j^{-1/2}\iota_l^{-1/2}\langle\Kcal_{\bar\theta}e_j,e_l\rangle)_{j,l\le J}$. Since $(1+\lambda_j)\iota_j$ lies between two positive constants, on $E_B$ \eqref{m:eq:infobound} gives $c_*\Id\preceq G\preceq C_*\Id$ and Lemma~\ref{m:lem:rdinfo}(i) gives $\norm{(\Id+\Lambda)\tilde\Kcal}_{op}\le C_2$, with constants depending on $B$ only. Let $\Gamma=n^{-1}\sum_{i\in\mathcal I}(\mathcal D^{1/2}\Phi_i^*\Phi_i\mathcal D^{1/2}-G)$, $M=\Id+\diag(\kappa_j^{-1})$ and $W=M+\tilde\Kcal+\Gamma$, so that $V_N=n^{-1}\mathcal D^{1/2}W^{-1}\mathcal D^{1/2}$. By the maximum principle in the proof of Lemma~\ref{m:lem:rdlin}, $|\I_{\bar\theta}e_j|\le\sqrt2e^{K_fT}$, so a whitened row has squared length at most a multiple of $J\lambda_J$ and conditional mean $\tr G\le C_*J$, and as in Step 1 of the proof of Proposition~\ref{m:prop:regheat}, $E[\norm\Gamma^2_{HS}\mid Z^{[b]}]\lesssim J^2\lambda_J/n\to0$. On $E_N=\{\norm\Gamma_{op}\le c_*/2\}$, $W\succeq(c_*/2)\Id$, and we use the resolvent identity $W^{-1}=M^{-1}-W^{-1}(\tilde\Kcal+\Gamma)M^{-1}$, where $M^{-1}\diag(\kappa_j^{-1})=\diag((1+\kappa_j)^{-1})$.

\emph{Decomposition.} Let $w_j$, $w^b_j$ and $\hat\Psi^b_{N,j}$ be as in the proof of Proposition~\ref{m:prop:suff}, $A=\I^*_{\bar\theta}\I_{\bar\theta}$ and $(A^{-1})_{JJ}=(\langle e_l,A^{-1}e_j\rangle)_{l,j\le J}$. Since $w^b_j=T\sum_{l\le J}\langle e_l,A^{-1}e_j\rangle\I_{\bar\theta}e_l$ and $n^{-1}\mathcal D^{1/2}G^{-1}\mathcal D^{1/2}=n^{-1}H^{-1}$, $\hat S^{{\rm reg},b}_N-(\hat\Psi^b_{N,j})_{j\le J}=T_1+\dots+T_6$ with
\[
\begin{aligned}
T_1&=-V_N\Sigma_0^{-1}\theta_0^D, & T_2&=V_N\Phi^\top\big((\I_{\bar\theta}\theta_0^{>D})(X_i)\big)_i,\\
T_3&=-V_N\Phi^\top\big(R(X_i)\big)_i, & T_4&=\big(V_N-n^{-1}H^{-1}\big)\Phi^\top\varepsilon,\\
T_5&=\big(H^{-1}-T(A^{-1})_{JJ}\big)n^{-1}\Phi^\top\varepsilon, & T_6&=\Big(\frac1n\sum_{i\in\mathcal I}(w^b_j-w_j)(X_i)\,\varepsilon_i\Big)_{j\le J}.
\end{aligned}
\]
The first paragraph and Step 3 of the proof of Proposition~\ref{m:prop:suff} treat the half sample weights and the coordinates $j>J$, so (i) follows once $\sqrt N\norm{T_m}_{-k-2}\to^P0$ for each $m$, which we show on $E_B\cap E_N$.

\emph{$T_1$.} By the resolvent identity, $T_1=-\mathcal D^{1/2}W^{-1}\diag(\kappa_j^{-1})\mathcal D^{-1/2}\theta_0^D=A_1'+\mathcal D^{1/2}W^{-1}(\Gamma+\tilde\Kcal)w$, with $A_1'$ as in Step 3 of the proof of Proposition~\ref{m:prop:regheat}, which bounds $A_1'$ and $\mathcal D^{1/2}W^{-1}\Gamma w$ up to constants, with $2/c_*$ in place of $2$. For the rest, $\norm{\tilde\Kcal w}\le C_2\norm{(\Id+\Lambda)^{-1}w}\lesssim\norm{\sum_j\theta_{0,j}(1+\kappa_j)^{-1}e_j}_{H^{-3}}$, whose square is at most a multiple of $\delta_N^{2s}\norm{\theta_0}^2_{H^\gamma}$ for $s=(\gamma+3)/(\gamma+1)$, since $(1+\kappa_j)^{-2}\le\kappa_j^{-s}$; and $N\delta_N^{2s}\to0$ because $(4-d)\gamma>d$.

\emph{$T_2$.} The summands $\mathcal D^{1/2}\Phi_i^*(\I_{\bar\theta}\theta_0^{>D})(X_i)$ have conditional mean $T^{-1}\mathcal D^{1/2}(\langle e_j,\Kcal_{\bar\theta}\theta_0^{>D}\rangle)_{j\le J}$, since $\Lcal^*\Lcal e_j$ is orthogonal to $\theta_0^{>D}$, whose norm is at most a multiple of $\norm{\Kcal_{\bar\theta}\theta_0^{>D}}_{H^1}\le C\lambda_J^{-(\gamma+3)/2}\norm{\theta_0}_{H^\gamma}$ by Lemma~\ref{m:lem:rdinfo}(i); it contributes $O(N\lambda_J^{-\gamma-3})=O(N^{(d-6)/(2\gamma+d)})$ to $N\norm{T_2}^2_{-k-2}$. The centered part is bounded as $A_2$ in Step 4 of the proof of Proposition~\ref{m:prop:regheat}, with $\norm{\I_{\bar\theta}\theta_0^{>D}}^2_{L^2(\mathcal X)}\le C_B^2\lambda_J^{-\gamma-1}\norm{\theta_0}^2_{H^\gamma}$, and contributes $O_P(J\lambda_J^{-\gamma})$.

\emph{$T_3$.} Since $n^{-1}\mathcal D^{1/2}\Phi^\top\Phi\mathcal D^{1/2}=G+\Gamma$, the Cauchy Schwarz inequality bounds $\norm{n^{-1}\mathcal D^{1/2}\Phi^\top(R(X_i))_i}^2$ by $\norm{G+\Gamma}_{op}n^{-1}\sum_{i\in\mathcal I}R(X_i)^2$, whose conditional mean is at most a multiple of $\norm R^2_{L^2(\mathcal X)}\lesssim\norm h^4_{L^2}$ by Lemma~\ref{m:lem:rdinfo}(ii). Hence $N\norm{T_3}^2_{-k-2}=O_P(N\tilde\delta_N(0)^4)$, which tends to zero because $2\gamma(\bar\gamma-1)>d(\bar\gamma+1)$.

\emph{$T_4$.} Let $\chi=n^{-1/2}\mathcal D^{1/2}\Phi^\top\varepsilon$, which given $Z^{[b]}$ and the design is centered with covariance $G+\Gamma$, and $\Lambda_k=\diag((1+\lambda_j)^{-(k+2)/2})$. Then $\sqrt nT_4=-\mathcal D^{1/2}W^{-1}(\Gamma+\diag(\kappa_j^{-1}))G^{-1}\chi$, and by the resolvent identity, on $E_N$, the Hilbert Schmidt norm of $\Lambda_k\mathcal D^{1/2}W^{-1}(\Gamma+\diag(\kappa_j^{-1}))G^{-1}$ is at most $c_*^{-1}a_N^{1/2}+2C_0c_*^{-2}(2\norm\Gamma_{HS}+\norm{\tilde\Kcal\diag((1+\kappa_j)^{-1})}_{HS})$. Here $a_N\to0$ as in Step 5 of the proof of Proposition~\ref{m:prop:regheat}, and the last norm is at most $C_2(\sum_{j\le J}(1+\kappa_j)^{-2}(1+\lambda_j)^{-2})^{1/2}$, which tends to zero by dominated convergence, since $\kappa_j\to\infty$ for each $j$ and $\sum_j(1+\lambda_j)^{-2}<\infty$ for $d\le3$. Markov's inequality given the design concludes.

\emph{$T_5$.} Write $A$ in blocks on the span of $e_1,\dots,e_J$ and on $QH^{-1}_0$, $Q=\Id-p_N$: $A_{11}=TH$, $A_{12}=p_N\Kcal_{\bar\theta}Q$, $A_{21}=Q\Kcal_{\bar\theta}p_N$ and $A_{22}=QAQ$, since $\Lcal^*\Lcal$ is diagonal. As in Step 1 of the proof of Proposition~\ref{m:prop:sumrate}, $A_{11}^{-1}$, $A_{22}^{-1}$ and $(A^{-1})_{JJ}$ map $H^1$ to $H^{-1}$ with norms at most $c_B^{-2}$, and, as in the proof of Proposition~\ref{m:prop:sumrd}, $A_{12}$ and $A_{21}$ map $H^{-1}$ to $H^1$ with norms at most $C(1+\lambda_J)^{-1}$. Since $(A^{-1})_{JJ}-A_{11}^{-1}=(A^{-1})_{JJ}A_{12}A_{22}^{-1}A_{21}A_{11}^{-1}$ and $E[\norm{n^{-1}\Phi^\top\varepsilon}^2_{H^1}\mid Z^{[b]}]\le C_B^2J/(nT)$, $N\norm{T_5}^2_{-k-2}\le N\norm{T_5}^2_{H^{-1}}=O_P(J\lambda_J^{-4})$, which tends to zero.

\emph{$T_6$.} As in Step 2 of the proof of Proposition~\ref{m:prop:suff}, with \eqref{m:eq:lipweight} for $\kappa=-1$ from Lemma~\ref{m:lem:rdinfo}(iii), $N\,E[\norm{T_6}^2_{-k-2}\mid Z^{[b]}]\lesssim\norm h^2_{H^\zeta}\sum_j(1+\lambda_j)^{-k-1}\to^P0$.

Part (ii) follows from (i) as in Step 9 of the proof of Proposition~\ref{m:prop:regheat}, with Theorem~\ref{m:thm:nickl2} for the limit law, and (iii) from (i), Proposition~\ref{m:prop:suff} and the triangle inequality. No condition on the exponents enters (i): the bias $T_1+T_2+T_3$ has no projection error.
\end{proof}

\begin{remark}[What the generator has to learn]\label{m:rem:affine}
Propositions~\ref{m:prop:linsuff} and~\ref{m:prop:suff} together say that the map from summary to posterior is affine in the linear case and asymptotically affine in the nonlinear case, provided $s$ is, in the linear case, the preconditioned summary $\tilde S_N$ of Proposition~\ref{m:prop:linsuff}, which is the regularized summary \eqref{m:eq:regsummary} by Lemma~\ref{m:lem:exact}, and in the nonlinear case the one step summary of Definition~\ref{m:def:summary}, and not the raw adjoint statistic, whose slope is $V_N$:
\[
Q_j(\tau\mid s)=s_j+\frac{z_\tau\sigma_j}{\sqrt N}+o(N^{-1/2}),\qquad\sigma_j^2=T\big\langle e_j,(\I^*_{\theta_0}\I_{\theta_0})^{-1}e_j\big\rangle .
\]
In the nonlinear case the expansion is to be read in the Wasserstein sense, that is with the remainder $o(N^{-1/2})$ taken in $L^1$ in $\tau$, which is what Proposition~\ref{m:prop:suff} gives for the posterior given all the data; for the posterior given $s$ alone it needs Assumption~\ref{m:ass:U}, as noted after that proposition. In the nonlinear case it is also local. It describes the posterior for data generated at a fixed $\theta_0$, which concentrates on a shrinking neighborhood of $\theta_0$, the localization set of Section~\ref{m:sec:main}, and the summary \eqref{m:eq:onestep} is a one step efficient estimator only there; nothing is claimed for a draw of $\theta$ from the prior far from $\theta_0$, and Section~\ref{m:sec:loc} turns this caveat into a rule for training the generator. The generator is asked to learn an identity map plus a scale. The scale is the diagonal of the inverse information operator, which by \eqref{m:eq:iota} is $1/\iota_j=2\nu\lambda_jT/(1-e^{-2\nu\lambda_jT})$ up to a compact perturbation in the dissipative case. The correct architecture is a skip connection carrying $s$ to the output plus a learned, low rank plus diagonal scale.
\end{remark}

\subsection{Information as expected utility, and the role of quadratic variation}\label{m:sec:pr}

\subsubsection{Information as a change in expected utility} Polson and Roberts~\cite{polsonroberts1993} take a Bayesian decision problem in which the experimenter observes $X$ on $[0,T]$ solving $dX_t=S(\vartheta,t,X_t)dt+\sigma(t,X_t)dB_t$ and must report beliefs about $\vartheta$, scored by a proper utility. The information in the experiment is defined as the change in expected utility from prior to posterior; for the logarithmic score this is the mutual information $I(\vartheta;X)=E\log\frac{d\Pi(\cdot\mid X)}{d\Pi}$. Their main result, for locally asymptotically mixed normal families, is an asymptotic relationship between the Shannon information of the posterior and the Fisher information of the process. In our setting the corresponding statement is an exact identity in the linear Gaussian models of the heat and Stokes equations, and it gives the decision theoretic content of the Bernstein von Mises limit.

\begin{proposition}[Expected utility form of the Bernstein von Mises limit]\label{m:prop:pr}
For the heat or the Stokes equation, consider the first $J$ modes with prior $\Pi^J=N(0,\Sigma_0)$, $\Sigma_0=\diag(\sigma^2_{0,j})$, in the Gaussian experiment that observes $\theta_j+(N\iota_j)^{-1/2}\xi_j$, $j\le J$, with independent standard normal $\xi_j$ and $\iota_j$ as in \eqref{m:eq:gram}. Its information matrix is the population value $\diag(N\iota_j)$ of the Gram matrix in \eqref{m:eq:gram}, and it is the $J$ mode experiment of the Bernstein von Mises limit, whose covariance is $T(\I^*\I)^{-1}/N=\diag(1/(N\iota_j))$. Write
\[
\kappa_j=N\iota_j\sigma^2_{0,j}
\]
for the per mode ratio of information to prior precision, and let $I_N=E^{\Pi}\big[\log\tfrac{d\Pi^J(\cdot\mid Z^{(N)})}{d\Pi^J}\big|Z^{(N)}\big]$ be the realized information gain, with $\Pi^J(\cdot\mid Z^{(N)})$ the posterior. Then, for every $N$ and $J$,
\[
I_N=\tfrac12\sum_{j\le J}\Big[\log(1+\kappa_j)-\frac{\kappa_j}{1+\kappa_j}\Big]+\tfrac12\,\tilde\theta_N^{J\top}\Sigma_0^{-1}\tilde\theta^J_N ,
\]
with $\tilde\theta^J_N$ the posterior mean, and averaging over the prior predictive law gives the Lindley expected information gain
\[
\mathbf{E}\,I_N=\tfrac12\log\det\big(\Id+\tfrac NT\,\I^*\I\,\Sigma_0\big)=\tfrac12\sum_{j\le J}\log\big(1+\kappa_j\big) .
\]
In the linear model \eqref{m:eq:design} with $D=J$, that is, with the realized information $\Phi^\top\Phi$ in place of $\diag(N\iota_j)$, both identities hold conditionally on the design, with the $\kappa_j$ replaced by the eigenvalues of $\Sigma_0^{1/2}\Phi^\top\Phi\,\Sigma_0^{1/2}$.
Every mode contributes the strictly positive amount $\tfrac12\log(1+\kappa_j)$, and every mode has posterior variance $\sigma^2_{0,j}/(1+\kappa_j)$ strictly below its prior variance; the modes with $\kappa_j>1$ are those whose posterior variance falls below \emph{half} the prior variance. Under the prior \eqref{m:eq:prior} the number of these is
\[
J^{\rm eff}_N=\#\{j\le J_N:\kappa_j>1\}\asymp N^{\gamma d/((\gamma+1)(2\gamma+d))}=J_N^{\gamma/(\gamma+1)}=o(J_N) ,
\]
so the experiment is informative in every retained mode, but strongly informative in only a vanishing fraction of them.
\end{proposition}

\begin{proof}
The posterior is $N(m,V)$ with $m=\tilde\theta^J_N$, and the Kullback Leibler divergence between $N(m,V)$ and $N(0,\Sigma_0)$ is $\tfrac12[\log\det(\Sigma_0V^{-1})-J+\mathrm{tr}(\Sigma_0^{-1}V)+m^\top\Sigma_0^{-1}m]$. Here $V^{-1}=\tfrac NT\I^*\I+\Sigma_0^{-1}$, whose $j$th diagonal entry is $N\iota_j+\sigma_{0,j}^{-2}$ by \eqref{m:eq:modeinfo} and which is diagonal because $\I^*\I$ is diagonal in $(e_j)$ for the heat and Stokes equations. So $\Sigma_0V^{-1}=\Id+\tfrac NT\Sigma_0\I^*\I$ has eigenvalues $1+\kappa_j$ and $\mathrm{tr}(\Sigma_0^{-1}V)=\sum_j(1+\kappa_j)^{-1}$; collecting terms gives the first display.

Under the prior predictive law the posterior mean satisfies $m\sim N(0,\Sigma_0-V)$, so $\mathbf{E}[m^\top\Sigma_0^{-1}m]=\mathrm{tr}(\Sigma_0^{-1}(\Sigma_0-V))=\sum_j\kappa_j(1+\kappa_j)^{-1}$, which cancels the second summand of the first display and leaves $\tfrac12\sum_j\log(1+\kappa_j)$. This is the classical identity that the prior predictive average of the realized information gain is the mutual information $I(\theta;Z^{(N)})$. With the realized information the posterior is $N(m,V)$ with $V^{-1}=\Phi^\top\Phi+\Sigma_0^{-1}$ by \eqref{m:eq:linpost}, $\Sigma_0V^{-1}$ is similar to $\Id+\Sigma_0^{1/2}\Phi^\top\Phi\,\Sigma_0^{1/2}$, and the same two computations, conditionally on the design, give the last assertion.

For the count, \eqref{m:eq:prior} gives $\sigma^2_{0,j}=\lambda_j^{-\gamma}/(N\delta_N^2)$, hence $\kappa_j=\iota_j\lambda_j^{-\gamma}\delta_N^{-2}$, and by \eqref{m:eq:iota} $\iota_j=c_j/(2\nu T\lambda_j)$ with $c_j=1-e^{-2\nu\lambda_jT}\in[1-e^{-2\nu\lambda_1T},1)$. So $\kappa_j>1$ if and only if $\lambda_j^{\gamma+1}<c_j(2\nu T\delta_N^2)^{-1}$, and with $\lambda_j\asymp j^{2/d}$ and $\delta_N=N^{-\gamma/(2\gamma+d)}$ this gives $J^{\rm eff}_N\asymp N^{\gamma d/((\gamma+1)(2\gamma+d))}$. Comparing with $J_N\asymp N\delta_N^2=N^{d/(2\gamma+d)}$ gives the stated power of $J_N$.
\end{proof}

\begin{remark}[What the effective dimension is not]\label{m:rem:jeff}
The summand is $\log(1+\kappa_j)$, not $\log\kappa_j$: $\tfrac12\sum_{j\le J_N}\log\kappa_j$ is negative and divergent, since $\kappa_j\to0$ for the modes near the truncation, and so cannot be a Kullback Leibler divergence. With $\gamma=3$, $d=2$ and the prior \eqref{m:eq:prior} it takes the values $-104$, $-355$, $-1275$ and $-4579$ at $N=10^6,10^8,10^{10},10^{12}$, against the correct $0.68$, $6.72$, $24.98$ and $71.54$ (these values are computed for $\nu=0.01$ and $T=1$, the spectrum $\lambda_k=4\pi^2|k|^2$, $k\in\Z^2\setminus\{0\}$, of the torus $[0,1]^2$ with the modes ordered by $|k|$, and $J_N$ the integer nearest to $N^{1/4}$, which gives $J_N=32$, $100$, $316$ and $1000$ and $J^{\rm eff}_N=0$, $8$, $20$ and $36$). Which count governs a simulation budget depends on the prior and on the generator, and three counts are in play. $J_N$ and $J^{\rm eff}_N$ belong to the rescaled prior \eqref{m:eq:prior}, and it is $J_N$, not $J^{\rm eff}_N$, that governs the simulation budget of the identity skip generator of Definition~\ref{m:def:gen} in Theorem~\ref{m:thm:qnnA}, whose head acts on all $J_N$ retained modes. The count $J_{\rm eff}$ of Theorem~\ref{m:thm:jeff} applies the same criterion $N\iota_j\sigma^2_{0,j}\ge1$ to the fixed prior $\sigma^2_{0,j}=\rho^2(\lambda_j/\lambda_1)^{-\alpha}$ and has the order of $J_N$ when $d=2$ and $\alpha=\gamma$; with $\rho=1$, $\alpha=3$ and the settings above it is $100$, $340$, $1048$ and $3332$, about three times $J_N$. It sets the cut of the hybrid generator of Theorem~\ref{m:thm:jeff}(ii), which estimates the whitened second moment on the first $\min(J_{\rm eff},J_N)$ modes and uses the analytic scale beyond, and it bounds that generator's budget, not the budget of the generator of Definition~\ref{m:def:gen}. For $\alpha=\gamma$ the cut is $J_N$ itself, so in these settings the hybrid generator saves nothing against the generator of Definition~\ref{m:def:gen}; the comparison in which $J_{\rm eff}$ matters is the one with the raw data in Theorem~\ref{m:thm:raw}.
\end{remark}

\begin{remark}[Design]
For the heat and Stokes equations with observation points uniform in space and observation times drawn from a density $q$ on $[0,T]$, the information matrix stays diagonal, with per mode information $\iota_j(q)=\int_0^Te^{-2\nu\lambda_jt}q(t)\,dt$ in place of $\iota_j$, so Proposition~\ref{m:prop:pr} holds verbatim with $\kappa_j(q)=N\iota_j(q)\sigma^2_{0,j}$ and makes the observation time design problem a utility maximization: choosing $q$ to maximize $\sum_j\log(1+\kappa_j(q))$ is a Bayesian $D$ optimal design in the Polson Roberts sense. By \eqref{m:eq:iota} it front loads observations near $t=0$, because $\iota_j$ is smaller by a factor of order $e^{-2\nu\lambda_jT_0}$ when sampling starts at $T_0>0$, and every summand is smaller with it. The dissipation barrier of Table~\ref{m:tab:iota} is a statement about expected utility: an experiment that begins at $T_0>0$ forgoes $\tfrac12\sum_j\big[\log(1+\kappa_j)-\log(1+\kappa_j(T_0))\big]$ of the information $\tfrac12\sum_j\log(1+\kappa_j)$ of Proposition~\ref{m:prop:pr}, with $\kappa_j(T_0)=N\iota_j(T_0)\sigma^2_{0,j}$ and $\iota_j(T_0)$ as in \eqref{m:eq:iota}, the value for $q$ uniform on $[T_0,T]$. For a nonlinear forward map the proofs of Lemma~\ref{m:lem:prexact} and of the bound \eqref{m:eq:prvt} apply verbatim to any design independent of $\theta$, with the Fisher information computed under it, and they bound the expected information gain by the information averaged over the prior, which under \eqref{m:eq:prior}, for a fixed number of modes and under the regularity hypotheses of the second part of Proposition~\ref{m:prop:prvt}, is the information at the center of the prior up to a factor $1+o(1)$, and not the information at $\theta_0$ (Proposition~\ref{m:prop:prrd}); the criterion with the efficient information at $\theta_0$ (Corollary~\ref{m:cor:prreal}) describes the Gaussian limit experiment, and then only for a design under which the Bernstein von Mises theorem holds, which must be assumed, since Theorem~\ref{m:thm:nickl2} is proved for independent uniform times and points.
\end{remark}

\subsubsection{The posterior itself under a nonlinear forward map}

Proposition~\ref{m:prop:pr} is exact for linear Gaussian models. For a nonlinear forward map its two displays have to be separated, because they are of different kinds: the first is a statement about the posterior for data from the truth $\theta_0$, the second is an average over the prior predictive law and so cannot depend on $\theta_0$ at all. Throughout, $J$ is a fixed number of modes, $\theta^J=(\langle\theta,e_j\rangle)_{j\le J}$, $\Pi^J$ and $\Pi^J(\cdot\mid Z^{(N)})$ are its laws under the prior and the posterior, $\Sigma^J_0=\diag(\sigma^2_{0,j})_{j\le J}$, $m_N$ and $V_N$ are the posterior mean and covariance of $\theta^J$, and $I_N=\KL\big(\Pi^J(\cdot\mid Z^{(N)})\,\|\,\Pi^J\big)$ is the realized information gain of Proposition~\ref{m:prop:pr}. We write $\mathbf P$ for the prior predictive law, under which $\vartheta\sim\Pi$ and $Z^{(N)}\sim P^\N_\vartheta$ given $\vartheta$, and $\mathbf E$ for its expectation. By the computation in the proof of Proposition~\ref{m:prop:pr}, for $m\in\R^J$ and positive definite $C$,
\[
\begin{aligned}
\KL\big(N(m,C)\,\|\,\Pi^J\big)=\tfrac12\Big[&\log\det\big(\Sigma^J_0C^{-1}\big)-J\\
&{}+\tr\big((\Sigma^J_0)^{-1}C\big)+m^\top(\Sigma^J_0)^{-1}m\Big],
\end{aligned}
\]
and the first display of Proposition~\ref{m:prop:pr} is this formula at $m=m_N$, $C=V$.

\begin{lemma}[Exact identities for any forward map]\label{m:lem:prexact}
Let $\Gcal$ be a measurable forward map in \eqref{m:eq:model}, and let the prior be a centered Gaussian measure under which the coordinates $\langle\theta,e_j\rangle$ are independent with variances $\sigma^2_{0,j}>0$, as under \eqref{m:eq:prior}. Then for every $N$ and $J$ the following hold.
\begin{enumerate}[label=(\roman*),nosep]
\item Almost surely $I_N=\KL\big(N(m_N,V_N)\,\|\,\Pi^J\big)+\Gamma_N$, where $\Gamma_N=\KL\big(\Pi^J(\cdot\mid Z^{(N)})\,\|\,N(m_N,V_N)\big)\ge0$, and $\Gamma_N\le\KL\big(\Pi^J(\cdot\mid Z^{(N)})\,\|\,G\big)$ for every nondegenerate Gaussian law $G$ on $\R^J$.
\item If $\int\norm{\Gcal(\vartheta)}^2_{L^\infty(\mathcal X)}\,d\Pi(\vartheta)<\infty$, then $\mathbf E\,I_N=I(\theta^J;Z^{(N)})<\infty$, and
\begin{equation}\label{m:eq:prexact}
\mathbf E\,I_N=\tfrac12\,\mathbf E\log\det\big(\Sigma^J_0V_N^{-1}\big)+\mathbf E\,\Gamma_N .
\end{equation}
\end{enumerate}
In the linear Gaussian models of Proposition~\ref{m:prop:pr} the posterior is $N(m_N,V)$ with $V$ not depending on the data, so $\Gamma_N=0$, and (i) and \eqref{m:eq:prexact} are its two displays.
\end{lemma}

\begin{proof}
The posterior has density $e^{\ell_N}/\int e^{\ell_N}d\Pi$ with respect to the prior, where $\ell_N(\vartheta)=-\tfrac12\sum_{i\le N}|Y_i-\Gcal(\vartheta)(X_i)|_W^2\le0$ is the log likelihood up to a constant, and $\int e^{\ell_N}d\Pi>0$ because $\ell_N$ is finite. So $\Pi^J(\cdot\mid Z^{(N)})$ has a Lebesgue density bounded by a multiple of that of $\Pi^J$; it therefore has finite second moments and a nonsingular covariance $V_N$, and $I_N<\infty$. For a nondegenerate Gaussian law $G$ the logarithm of the density of $N(m_N,V_N)$ with respect to $G$ is a polynomial of degree two, whose integrals against $\Pi^J(\cdot\mid Z^{(N)})$ and against $N(m_N,V_N)$ agree, since the two laws have the same first and second moments. Hence
\[
\KL\big(\Pi^J(\cdot\mid Z^{(N)})\,\|\,G\big)=\Gamma_N+\KL\big(N(m_N,V_N)\,\|\,G\big).
\]
With $G=\Pi^J$ this is (i), and the inequality in (i) follows because the last term is nonnegative.

For (ii), the prior predictive average of the information gain is the mutual information, as recalled in Chapter~\ref{ch:entropy}, here applied to $\theta^J$. It is finite: given the parameter the observations are independent, so $I(\theta^J;Z^{(N)})\le I(\theta;Z^{(N)})\le\sum_{i\le N}I(\theta;(X_i,Y_i))$, and since $X_i$ is independent of $\theta$ and the Gaussian law maximizes entropy for a given covariance, each summand is at most the capacity of a Gaussian channel whose input has second moment at most $\int\norm{\Gcal(\vartheta)}^2_{L^\infty}d\Pi(\vartheta)$. Since $\mathbf E\,m_N=\mathbf E\,\theta^J=0$, the law of total covariance gives $\Sigma^J_0=\mathbf E\,V_N+\mathbf E\big[m_Nm_N^\top\big]$, hence
\[
\mathbf E\big[m_N^\top(\Sigma^J_0)^{-1}m_N\big]=J-\mathbf E\,\tr\big((\Sigma^J_0)^{-1}V_N\big).
\]
Both terms are integrable, and so is $\KL(N(m_N,V_N)\,\|\,\Pi^J)$, which lies between $0$ and $I_N$ by (i); hence so is $\log\det(\Sigma^J_0V_N^{-1})$, and taking expectations in the formula for $\KL(N(m_N,V_N)\,\|\,\Pi^J)$ gives $\mathbf E\,\KL(N(m_N,V_N)\,\|\,\Pi^J)=\tfrac12\mathbf E\log\det(\Sigma^J_0V_N^{-1})$. With (i) this is \eqref{m:eq:prexact}.
\end{proof}

\begin{theorem}[Kullback Leibler form of the Bernstein von Mises theorem]\label{m:thm:prkl}
Assume the hypotheses of Theorem~\ref{m:thm:nickl2}. Let $\Gcal_{\rm lin}(\theta)=\Gcal(\theta_0)+\I_{\theta_0}[\theta-\theta_0]$, and let $\Pi_{\rm lin}(\cdot\mid Z^{(N)})$ be the posterior in model \eqref{m:eq:model} with $\Gcal_{\rm lin}$ in place of $\Gcal$, the same prior \eqref{m:eq:prior} and the same data. It is a Gaussian measure, and
\begin{equation}\label{m:eq:prkl}
\KL\big(\Pi(\cdot\mid Z^{(N)})\,\big\|\,\Pi_{\rm lin}(\cdot\mid Z^{(N)})\big)\to0\qquad\text{in }P^\N_{\theta_0}\text{ probability}.
\end{equation}
\end{theorem}

\begin{proof}
The posterior $\Pi_{\rm lin}(\cdot\mid Z^{(N)})$ is Gaussian because the prior is and the likelihood of $\Gcal_{\rm lin}$ is the exponential of a quadratic function of finitely many linear functionals of $\theta$. We use four facts from Nickl's proof of Theorem~\ref{m:thm:nickl2}~\cite{nicklbvm2024}, stated for the sets $\bar\Theta_{N,L}$ of (37) there, which consist of the $\theta$ with $\norm\theta_{H^{\bar\gamma}}\le L$, $\norm{\Gcal(\theta)-\Gcal(\theta_0)}_{L^2(\mathcal X)}\le L\delta_N$, $\norm{\theta-\theta_0}_{H^\xi}\le L\tilde\delta_N(\xi)$ for $0\le\xi<\bar\gamma$, and one more condition on the reproducing kernel Hilbert space norm, with $\tilde\delta_N(\xi)=\delta_N^{(\bar\gamma-\xi)/(\bar\gamma+1)}$ and $\bar\gamma$, $\zeta$ as there. Write $g_\theta=\Gcal(\theta)-\Gcal(\theta_0)$, $\bar g_\theta=\I_{\theta_0}[\theta-\theta_0]$ and $r_\theta=g_\theta-\bar g_\theta$.
\begin{enumerate}[label=(\alph*),nosep]
\item Theorem~3 there: for every $b>0$ there is $L$ such that $\Pi(\bar\Theta_{N,L}\mid Z^{(N)})\ge1-e^{-bN\delta_N^2}$ with probability tending to one, and $E^\Pi[\norm{\theta-\theta_0}^q_{H^\xi}\mid Z^{(N)}]=O_P(\tilde\delta_N(\xi)^q)$ for $1\le q<\infty$ and $0\le\xi\le\bar\gamma$; these are (38) and (39) there.
\item Lemma~1 there, with $t=0$: for every $L$, the suprema over $\bar\Theta_{N,L}$ of $\big|\sum_i\langle\varepsilon_i,r_\theta(X_i)\rangle_W\big|$ and of $\big|\sum_i\big(|g_\theta(X_i)|^2_W-E|g_\theta(X_1)|^2_W\big)\big|$ tend to zero in $L^1$.
\item $N\tilde\delta_N(0)^3=N^{-a}$ for some $a>0$ under the hypotheses on $\gamma$, as shown after (58) there.
\item Conditions 2C and 2D: $\I_{\theta_0}$ is continuous from $L^2(\Omega)$ into $L^2(\mathcal X)$; on $\bar\Theta_{N,L}$, $\norm{r_\theta}_{L^2(\mathcal X)}\lesssim\norm{\theta-\theta_0}^2_{L^2}\lesssim\tilde\delta_N(0)^2$; and $\norm{\I_{\theta_0}h}_\infty\lesssim\norm h_{H^\zeta}$.
\end{enumerate}

\emph{Step 1: the linearized model satisfies the same hypotheses.} Conditions 2A to 2D hold for $\Gcal_{\rm lin}$ by Conditions 2A, 2C and 2D for $\Gcal$, the remainder in 2C being zero, and 2F holds because the information operator of $\Gcal_{\rm lin}$ is $\I^*_{\theta_0}\I_{\theta_0}$. For 2E, let $\norm\theta_{H^{\bar\gamma}}\le B$, $h=\theta-\theta_0$ and $0<s\le1$; Conditions 2E and 2C at $\theta_0+sh$, whose $H^{\bar\gamma}$ norm is at most $B+2\norm{\theta_0}_{H^{\bar\gamma}}$, give $s\norm{\I_{\theta_0}h}_{L^2(\mathcal X)}\ge cs\norm h_{H^{-1}}-c's^2\norm h^2_{L^2}$, and letting $s\to0$ gives $\norm{\I_{\theta_0}h}_{L^2(\mathcal X)}\ge c\norm h_{H^{-1}}$. So (a) and (b) hold for $\Gcal_{\rm lin}$ with its own sets $\bar\Theta^{\rm lin}_{N,L}$, under the same law of the data, since $\Gcal_{\rm lin}(\theta_0)=\Gcal(\theta_0)$. Since $\tilde\delta_N(0)^2=o(\delta_N)$, (d) gives $\bar\Theta_{N,L}\subset\bar\Theta^{\rm lin}_{N,2L}$ and $\bar\Theta^{\rm lin}_{N,L}\subset\bar\Theta_{N,2L}$ for $N$ large.

\emph{Step 2: an elementary bound.} Let $\ell_N$ and $\ell^{\rm lin}_N$ be the log likelihoods of the two models, $R=\ell_N-\ell^{\rm lin}_N$, $P=\Pi(\cdot\mid Z^{(N)})$ and $Q=\Pi_{\rm lin}(\cdot\mid Z^{(N)})$. Then $dP/dQ=e^R/E_Qe^R$. For every measurable $\Theta$, $E_Qe^R\ge Q(\Theta)\exp(-\sup_\Theta|R|)$, and $R\le-\ell^{\rm lin}_N$ because $\ell_N\le0$, so
\[
\begin{aligned}
\KL(P\,\|\,Q)&=E_PR-\log E_Qe^R\\
&\le2\sup_\Theta|R|+E_P\big[(-\ell^{\rm lin}_N)\,1_{\Theta^c}\big]-\log Q(\Theta).
\end{aligned}
\]

\emph{Step 3: the three terms.} Fix $L$ as in (a) with $b=1$, for both models, and take $\Theta=\bar\Theta_{N,2L}$. Under $P^\N_{\theta_0}$, $Y_i=\Gcal(\theta_0)(X_i)+\varepsilon_i$, and expanding the squares,
\[
R(\theta)=\sum_{i\le N}\langle\varepsilon_i,r_\theta(X_i)\rangle_W-\tfrac12\sum_{i\le N}\big(|g_\theta(X_i)|^2_W-|\bar g_\theta(X_i)|^2_W\big).
\]
The mean of the second sum is $\tfrac NT\big(\norm{g_\theta}^2_{L^2(\mathcal X)}-\norm{\bar g_\theta}^2_{L^2(\mathcal X)}\big)=\tfrac NT\langle r_\theta,r_\theta+2\bar g_\theta\rangle_{L^2(\mathcal X)}$, which on $\Theta$ is $O(N\tilde\delta_N(0)^3)=O(N^{-a})$ by (c) and (d). The centered sums and the first sum are controlled by (b), for $\Gcal$ on $\Theta$ and for $\Gcal_{\rm lin}$ on $\bar\Theta^{\rm lin}_{N,4L}\supset\Theta$, so $\sup_\Theta|R|\to0$ in probability. Next, $-\ell^{\rm lin}_N(\theta)\le\sum_i|\varepsilon_i|_W^2+N\norm{\bar g_\theta}^2_\infty\le\sum_i|\varepsilon_i|_W^2+CN\norm{\theta-\theta_0}^2_{H^\zeta}$ by (d), so by the Cauchy Schwarz inequality and (a), with $q=4$ and $\xi=\zeta$, $E_P[(-\ell^{\rm lin}_N)1_{\Theta^c}]=O_P(N)e^{-N\delta_N^2/2}$, which tends to zero because $N\delta_N^2=N^{d/(2\gamma+d)}$. Finally $Q(\Theta)\ge Q(\bar\Theta^{\rm lin}_{N,L})\ge1-e^{-N\delta_N^2}$ with probability tending to one, by (a) for $\Gcal_{\rm lin}$, so $-\log Q(\Theta)\to0$. Step 2 gives \eqref{m:eq:prkl}.
\end{proof}

\begin{corollary}[The first display for the posterior]\label{m:cor:prreal}
Under the hypotheses of Theorem~\ref{m:thm:prkl}, $\Gamma_N\le\KL\big(\Pi(\cdot\mid Z^{(N)})\,\|\,\Pi_{\rm lin}(\cdot\mid Z^{(N)})\big)$ for every $J$, so $\Gamma_N\to0$ in $P^\N_{\theta_0}$ probability uniformly in $J$. For fixed $J$, in $P^\N_{\theta_0}$ probability, $NV_N\to\Sigma_J(\theta_0)=\big(T\langle e_j,(\I^*_{\theta_0}\I_{\theta_0})^{-1}e_{j'}\rangle\big)_{j,j'\le J}$, the covariance of the first $J$ coordinates under $\Ncal_{\theta_0}$, which is positive definite, and
\begin{equation}\label{m:eq:prreal}
I_N=\tfrac12\sum_{j\le J}\Big[\log(1+\kappa_j)-\frac{\kappa_j}{1+\kappa_j}\Big]+\tfrac12\,m_N^\top(\Sigma^J_0)^{-1}m_N+o_P(1),
\end{equation}
where the $\kappa_j$ are the eigenvalues of $N(\Sigma^J_0)^{1/2}\Sigma_J(\theta_0)^{-1}(\Sigma^J_0)^{1/2}$. For the heat and Stokes equations $\Sigma_J(\theta_0)=\diag(1/\iota_j)$ and $\kappa_j=N\iota_j\sigma^2_{0,j}$, as in Proposition~\ref{m:prop:pr}.
\end{corollary}

\begin{proof}
The first assertion is Lemma~\ref{m:lem:prexact}(i) with $G$ the law of $\theta^J$ under $\Pi_{\rm lin}(\cdot\mid Z^{(N)})$, together with the fact that the Kullback Leibler divergence does not increase under the map $\theta\mapsto\theta^J$. Write $N(\mu_N,C_N)$ for that law and $\epsilon_N$ for the Kullback Leibler divergence of $\Pi^J(\cdot\mid Z^{(N)})$ from it, so that $\epsilon_N\to0$. For a unit vector $w\in\R^J$ let $X=\langle C_N^{-1/2}w,\theta^J-\mu_N\rangle$, which is standard normal under $N(\mu_N,C_N)$. The variational inequality $E_P\psi\le\KL(P\,\|\,Q)+\log E_Qe^\psi$, which is the nonnegativity of $\KL(P\,\|\,Q_\psi)$ for the tilted law $dQ_\psi=e^\psi\,dQ/E_Qe^\psi$, applied with $\psi=\lambda X$ for all $\lambda\in\R$ and with $\psi=\lambda(X^2-1)$ for $|\lambda|\le1/4$, where $\log E\,e^{\lambda(X^2-1)}\le2\lambda^2$, gives $|E^\Pi[X\mid Z^{(N)}]|\le\sqrt{2\epsilon_N}$ and $|E^\Pi[X^2\mid Z^{(N)}]-1|\le2\sqrt{2\epsilon_N}$ once $\epsilon_N\le1/8$. Hence $\norm{C_N^{-1/2}(m_N-\mu_N)}\le\sqrt{2\epsilon_N}$ and $\norm{C_N^{-1/2}V_NC_N^{-1/2}-\Id}\le4\sqrt{\epsilon_N}$. By Step 1 of the proof of Theorem~\ref{m:thm:prkl}, Theorem~\ref{m:thm:nickl2} applies to $\Gcal_{\rm lin}$ and gives the convergence in $\Wone_1$ of the law of $\sqrt N(\theta-E^{\Pi_{\rm lin}}[\theta\mid Z^{(N)}])$ under $\Pi_{\rm lin}(\cdot\mid Z^{(N)})$ to $\Ncal_{\theta_0}$; since $\theta\mapsto\theta^J$ is Lipschitz from $H^{-k-2}$ to $\R^J$, $N(0,NC_N)$ converges weakly to $N(0,\Sigma_J(\theta_0))$, so $NC_N\to\Sigma_J(\theta_0)$, and with the previous bound $NV_N\to\Sigma_J(\theta_0)$. The matrix $\Sigma_J(\theta_0)$ is positive definite: for $w\in\R^J$, $w\ne0$, the function $g=(\I^*_{\theta_0}\I_{\theta_0})^{-1}\sum_{j\le J}w_je_j$ is nonzero and lies in $H^\eta$ for every $\eta\ge\eta_0$ by Condition 2F, so $w^\top\Sigma_J(\theta_0)w=T\langle\I^*_{\theta_0}\I_{\theta_0}g,g\rangle=T\norm{\I_{\theta_0}g}^2_{L^2(\mathcal X)}>0$ by the inequality $\norm{\I_{\theta_0}h}_{L^2(\mathcal X)}\ge c\norm h_{H^{-1}}$ of Step 1 of the proof of Theorem~\ref{m:thm:prkl}, applied with $h=g$. Now insert the convergence of $NV_N$ into Lemma~\ref{m:lem:prexact}(i). The smallest $\kappa_j$ is at least $\lambda_J^{-\gamma}\delta_N^{-2}/\norm{\Sigma_J(\theta_0)}\to\infty$, so $\log\det(\Sigma^J_0V_N^{-1})=\sum_j\log\kappa_j+o_P(1)=\sum_j\log(1+\kappa_j)+o_P(1)$ and $\sum_j\kappa_j/(1+\kappa_j)=J+o(1)$, while $\tr((\Sigma^J_0)^{-1}V_N)\le\norm{NV_N}\,\delta_N^2\sum_{j\le J}\lambda_j^\gamma\to0$. With $\Gamma_N\to0$ this is \eqref{m:eq:prreal}. For the heat and Stokes equations $\I^*\I$ is diagonal with entries $T\iota_j$ by \eqref{m:eq:gram}.
\end{proof}

\begin{proposition}[The Navier Stokes equations]\label{m:prop:prklns}
Let $\Gcal(\theta)=u_\theta$ solve \eqref{m:eq:nsproj} on $\Omega=[0,1]^2$ with $\nu>0$, $f,\theta_0\in\Hds{\infty}$ and $W=\R^2$, with the prior \eqref{m:eq:prior} on $\Hds{}$, $d=2$ and $\gamma>12$, which is that of Theorem~2.11 of Konen and Nickl~\cite{konennickl2026} with $\alpha=\gamma$, and with $(e_j)$ a real orthonormal eigenbasis of $A$ in $\Hds{}$ with nondecreasing eigenvalues. Then the conclusions of Theorem~\ref{m:thm:prkl} and Corollary~\ref{m:cor:prreal} hold.
\end{proposition}

\begin{proof}
Norms are the homogeneous ones of Konen and Nickl, equivalent to ours on mean zero fields, $\norm\cdot_\infty$ is the supremum over $\mathcal X$, $C$ does not depend on $N$, and the rest of the notation is that of the proof of Theorem~\ref{m:thm:prkl}, with $\bar\gamma=11$ and $\zeta=2$, their $\beta$ and $\zeta$ in (2.22), so that $\tilde\delta_N(2)=\delta_N^{3/4}$. We replace facts (a) to (d) of that proof by three others.
\begin{enumerate}[label=(\alph*$'$),nosep]
\item Theorem~2.13 of Konen and Nickl: for every $b>0$ there is $L$ such that, with probability tending to one, $\Pi(\bar\Theta_{N,L}\mid Z^{(N)})\ge1-e^{-bN\delta_N^2}$ for the set $\bar\Theta_{N,L}$ of (2.31) there, whose elements satisfy $\norm\theta_{\Hds{11}}\le L$ and $\norm{\theta-\theta_0}_{\Hds2}\le L\tilde\delta_N(2)$; and $E^\Pi[\norm{\theta-\theta_0}^4_{\Hds2}\mid Z^{(N)}]=O_P(\tilde\delta_N(2)^4)$ by (2.33) there.
\item The same holds for $\Gcal_{\rm lin}$, $Q$ and the sets (2.31) defined with $\Gcal_{\rm lin}$: Theorem~2.13 of Konen and Nickl is Theorem~3 of Nickl~\cite{nicklbvm2024} on divergence free fields (Remark~1 of Nickl and Section~2.5 of Konen and Nickl), whose proof uses only parts A, B and E of Condition~2, and for $\Gcal_{\rm lin}$ these follow from (c$'$) with part A for $\Gcal$ and from the two bounds of Remark~\ref{m:rem:infobound} at $\vartheta=\theta_0$.
\item Propositions~3.7, 3.9 and~3.12 of Konen and Nickl with $a=a^*=2$ and the embedding of $\Hds2$ into $C(\Omega)^2$: $\norm{\bar g_\theta}_\infty\le C\norm{\theta-\theta_0}_{\Hds2}$, and for every $L$ there is $C$ with $\norm{r_\theta-r_{\theta'}}_\infty\le C\norm{\theta-\theta'}_{\Hds2}$ and $\norm{r_\theta}_\infty\le C\norm{\theta-\theta_0}^2_{\Hds2}$ when $\norm\theta_{\Hds2},\norm{\theta'}_{\Hds2}\le L$.
\end{enumerate}
Let $L$ be as in (a$'$) and (b$'$) with $b=1$, and $\Theta=\{\theta:\norm\theta_{\Hds{11}}\le L,\ \norm{\theta-\theta_0}_{\Hds2}\le L\tilde\delta_N(2)\}$. With probability tending to one, $P(\Theta^c)\le e^{-N\delta_N^2}$ and $Q(\Theta)\ge1-e^{-N\delta_N^2}$; so $-\log Q(\Theta)\to0$, and Step~3 of the proof of Theorem~\ref{m:thm:prkl}, with (a$'$) and (c$'$) for (a) and (d) and with $N\delta_N^2=N^{1/(\gamma+1)}$, gives $E_P[(-\ell^{\rm lin}_N)1_{\Theta^c}]\to0$ in probability. By Step~2 there it remains to bound $\sup_\Theta|R|$ by a random variable that tends to zero in probability, which we do without Lemma~1 of Nickl. By the display in Step~3, $R=S-\tfrac12D$ with $S(\theta)=\sum_{i\le N}\langle\varepsilon_i,r_\theta(X_i)\rangle_W$ and $D(\theta)=\sum_{i\le N}\langle r_\theta,r_\theta+2\bar g_\theta\rangle_W(X_i)$. On $\Theta$, (c$'$) gives $|D(\theta)|\le N\norm{r_\theta}_\infty(\norm{r_\theta}_\infty+2\norm{\bar g_\theta}_\infty)\le CN\tilde\delta_N(2)^3=CN^{(8-\gamma)/(8\gamma+8)}$ for every design, which tends to zero. For $S$, let $\eta_N=N^{-3/2}$. As $\lambda_j\asymp j$, deleting the coordinates $j>K$ moves a field of the ball of radius $L$ of $\Hds{11}$ by at most $\lambda_{K+1}^{-9/2}L\le\eta_N/2$ in $\Hds2$ for some $K\le C\eta_N^{-2/9}$; covering the ball of radius $L$ of this $K$ dimensional space by $(1+4L/\eta_N)^K$ balls of radius $\eta_N/2$, and choosing a point of $\Theta$ in each concentric ball of radius $\eta_N$ that meets $\Theta$, gives $\theta_1,\dots,\theta_M\in\Theta$ with $\log M\le CN^{1/3}\log N$ such that each $\theta\in\Theta$ is within $2\eta_N$ of some $\theta_k$ in $\Hds2$, and then $|S(\theta)-S(\theta_k)|\le2C\eta_N\sum_i|\varepsilon_i|_W=O_P(N^{-1/2})$ by (c$'$). Given the design, $S(\theta_k)$ is centered Gaussian with variance at most $N\norm{r_{\theta_k}}^2_\infty\le CN\delta_N^3$, and centered Gaussian $Z_1,\dots,Z_M$ with variances at most $\sigma^2$ satisfy $E\max_k|Z_k|\le\sigma(2\log(2M))^{1/2}$; so $\max_k|S(\theta_k)|=O_P(N^{(8-\gamma)/(12\gamma+12)}(\log N)^{1/2})$, and Step~2 gives \eqref{m:eq:prkl}.

In the proof of Corollary~\ref{m:cor:prreal}, Theorem~\ref{m:thm:nickl2} serves only to show $NC_N\to\Sigma_J(\theta_0)$, with $C_N$ the covariance of $\theta^J$ under $Q$, and Condition~2F only to show that $\Sigma_J(\theta_0)$ is positive definite. For the first, Theorem~2.11 of Konen and Nickl with $k=8$, applied to $P$, gives $E_Pe^{iw\cdot\sqrt N(\theta^J-m_N)}\to e^{-w^\top\Sigma_J(\theta_0)w/2}$ in probability for every $w\in\R^J$, as $\theta\mapsto\theta^J$ is Lipschitz from $\Hds{-10}$ to $\R^J$, $|e^{iw\cdot x}-e^{iw\cdot y}|\le|w||x-y|$, and the limit $\Ncal_{\theta_0}=N(0,T(\I^*_{\theta_0}\I_{\theta_0})^{-1})$ of Proposition~2.7 there gives $\theta^J$ the law $N(0,\Sigma_J(\theta_0))$. Under $Q$ this characteristic function has modulus $e^{-w^\top NC_Nw/2}$, and the two differ by at most twice the total variation distance of $P$ and $Q$, so by at most $(2\KL(P\,\|\,Q))^{1/2}\to0$ by Pinsker's inequality; hence $w^\top NC_Nw\to w^\top\Sigma_J(\theta_0)w$ for every $w$. For the second, if $w\ne0$ then $g=(\I^*_{\theta_0}\I_{\theta_0})^{-1}\sum_{j\le J}w_je_j$ is a nonzero element of $\Hds{-1}$ by Theorem~2.6 of Konen and Nickl with $a=-1$, and $w^\top\Sigma_J(\theta_0)w=T\norm{\I_{\theta_0}g}^2_{L^2(\mathcal X)}>0$ by the lower bound of Remark~\ref{m:rem:infobound}.
\end{proof}

Write
\[
\mathcal J^J_N(\vartheta)=\tfrac NT\big(\langle\I_\vartheta e_j,\I_\vartheta e_{j'}\rangle_{L^2(\mathcal X)}\big)_{j,j'\le J}
\]
for the $J$ block of the Fisher information of $N$ observations at $\vartheta$.

\begin{proposition}[Entropic van Trees bound]\label{m:prop:prvt}
Assume the hypotheses of Lemma~\ref{m:lem:prexact}(ii), and that for $\Pi$ almost every $\eta$ in the closed span of $(e_j)_{j>J}$ the map $x\mapsto\Gcal(\eta+\sum_{j\le J}x_je_j)$ from $\R^J$ into $C(\mathcal X)$ is continuously differentiable, with partial derivatives $\I_\vartheta e_j$ at $\vartheta=\eta+\sum_{j\le J}x_je_j$. If $\bar{\mathcal J}^J_N=\int\mathcal J^J_N(\vartheta)\,d\Pi(\vartheta)$ is finite, then for every $N$
\begin{equation}\label{m:eq:prvt}
\mathbf E\,I_N=I(\theta^J;Z^{(N)})\le\tfrac12\log\det\big(\Id+(\Sigma^J_0)^{1/2}\,\bar{\mathcal J}^J_N\,(\Sigma^J_0)^{1/2}\big).
\end{equation}
If moreover the prior is \eqref{m:eq:prior} with $\gamma>d$, and the matrix $\mathcal J^J_N(\vartheta)/N$, which does not depend on $N$, is continuous at $\vartheta=0$ for the supremum norm, bounded by a polynomial in $\norm\vartheta_\infty$, and invertible at $\vartheta=0$, then for fixed $J$
\begin{equation}\label{m:eq:prvtzero}
\mathbf E\,I_N\le\tfrac12\log\det\big(\Id+(\Sigma^J_0)^{1/2}\,\mathcal J^J_N(0)\,(\Sigma^J_0)^{1/2}\big)+o(1).
\end{equation}
\end{proposition}

\begin{proof}
Write $\theta=\theta^J+\eta$, with $\theta^J$ identified with $\sum_{j\le J}\langle\theta,e_j\rangle e_j$ and $\eta$ the component in the span of $(e_j)_{j>J}$. Under the prior $\theta^J$ and $\eta$ are independent, so the chain rule for mutual information gives $I(\theta^J;Z^{(N)})\le I(\theta^J;(Z^{(N)},\eta))=I(\theta^J;Z^{(N)}\mid\eta)$, and it suffices to bound the conditional mutual information for fixed $\eta$ and to average over $\eta$. Fix $\eta$ and let $p(x\mid z)$ be the conditional density of $\theta^J$ given $Z^{(N)}=z$ and $\eta$, which is proportional to $\phi(x)e^{\ell_N(\eta+\sum_jx_je_j)}$, with $\phi$ the density of $\Pi^J$. Its score is $\nabla_x\log p(x\mid z)=-(\Sigma^J_0)^{-1}x+s$, where $s_j=\sum_i\langle Y_i-\Gcal(\vartheta)(X_i),(\I_\vartheta e_j)(X_i)\rangle_W$ at $\vartheta=\eta+\sum_jx_je_j$. Under the joint law of $(\theta^J,Z^{(N)})$ given $\eta$, $Y_i-\Gcal(\vartheta)(X_i)=\varepsilon_i$ with $\vartheta=\eta+\theta^J$, so $s$ has conditional mean zero and conditional covariance $\mathcal J^J_N(\vartheta)$ given $\theta^J$, because $X_i$ is uniform on $\mathcal X$, of measure $T$. The cross term therefore vanishes, and the Fisher information matrix $F_z$ of $p(\cdot\mid z)$ has mean
\[
M_\eta=(\Sigma^J_0)^{-1}+\int\mathcal J^J_N\big(\eta+\textstyle\sum_jx_je_j\big)\,\phi(x)\,dx .
\]
For an invertible matrix $A$ the law of $A\theta^J$ given $(z,\eta)$ has entropy $h(p(\cdot\mid z))+\log|\det A|$ and Fisher information $\tr(A^{-\top}F_zA^{-1})$, so Theorem~\ref{thm:lsi} gives, trivially when that Fisher information is infinite,
\[
h\big(p(\cdot\mid z)\big)\ge-\log|\det A|+\frac J2\log\frac{2\pi e\,J}{\tr(A^{-\top}F_zA^{-1})} .
\]
Averaging over $z$, using the concavity of the logarithm and then taking $A=M_\eta^{1/2}$, the conditional entropy of $\theta^J$ given $(Z^{(N)},\eta)$ is at least $\tfrac12\log\det(2\pi e\,M_\eta^{-1})$. Since $\theta^J$ has entropy $\tfrac12\log\det(2\pi e\,\Sigma^J_0)$ and is independent of $\eta$, $I(\theta^J;Z^{(N)}\mid\eta)\le\tfrac12\log\det(\Sigma^J_0M_\eta)$. Averaging over $\eta$ and applying the concavity of $\log\det$ on positive definite matrices to $M_\eta$, whose average is $(\Sigma^J_0)^{-1}+\bar{\mathcal J}^J_N$ because $\eta$ and $\theta^J$ are independent, gives \eqref{m:eq:prvt}, since $\det\big(\Sigma^J_0((\Sigma^J_0)^{-1}+\bar{\mathcal J}^J_N)\big)$ is the determinant of the matrix in \eqref{m:eq:prvt}.

For \eqref{m:eq:prvtzero}, consider $\big(\bar{\mathcal J}^J_N-\mathcal J^J_N(0)\big)/N=\int\big(\mathcal J^J_N(\vartheta)-\mathcal J^J_N(0)\big)/N\,d\Pi(\vartheta)$, in which the integrand does not depend on $N$. Under \eqref{m:eq:prior}, $\theta=\theta'/(\sqrt N\delta_N)$ with $\theta'=\sum_j\lambda_j^{-\gamma/2}g_je_j$. Since $\gamma>d$ we may fix $s$ with $d/2<s<\gamma-d/2$; then, as $|e_j|\le\sqrt2$, the Cauchy Schwarz inequality gives $\sum_j|\langle\theta',e_j\rangle|\,\norm{e_j}_\infty\le C_s\norm{\theta'}_{H^s}$ with $C_s^2=2\sum_j(1+\lambda_j)^{-s}<\infty$, and $\norm{\theta'}^2_{H^s}=\sum_jw_jg_j^2$ with $w_j=(1+\lambda_j)^s\lambda_j^{-\gamma}$ summable, so that $E\,e^{t\norm{\theta'}^2_{H^s}}=\prod_j(1-2tw_j)^{-1/2}<\infty$ for $0<t<1/(2\max_jw_j)$. Hence the series for $\theta'$ converges uniformly, $\theta'$ is continuous, and $\Pi'(\norm{\theta'}_\infty>r)\le C\,e^{-cr^2}$ for some $C,c>0$ and all $r>0$, where $\Pi'$ is the law of $\theta'$. As $\sqrt N\delta_N\to\infty$, $\Pi(\norm\theta_\infty>\epsilon)\to0$ for every $\epsilon>0$, and splitting the integral at $\norm\vartheta_\infty=\epsilon$, with the continuity at $\vartheta=0$ on one side and the polynomial bound and the Gaussian tails on the other, shows that it tends to zero. Since $\mathcal J^J_N(0)/N$ is invertible and does not depend on $N$, there are $\epsilon_N\to0$ with $\bar{\mathcal J}^J_N\preceq(1+\epsilon_N)\mathcal J^J_N(0)$, so the matrix in \eqref{m:eq:prvt} is at most $\Id+(1+\epsilon_N)L_N$ in the Loewner order, $L_N=(\Sigma^J_0)^{1/2}\mathcal J^J_N(0)(\Sigma^J_0)^{1/2}$, and $\log\det(\Id+(1+\epsilon_N)L_N)\le\log\det(\Id+L_N)+J\log(1+\epsilon_N)$.
\end{proof}

For the reaction diffusion class of Theorem~\ref{m:thm:nickl1} the information at the center of the prior is explicit. Let $v$ solve $v'=f(v)$, $v(0)=0$, let $\chi(t)=\int_0^tf'(v(s))\,ds$, and set
\begin{equation}\label{m:eq:iotazero}
\iota^0_j=\frac1T\int_0^Te^{2\chi(t)-2\lambda_jt}\,dt,\qquad\kappa^0_j=N\iota^0_j\sigma^2_{0,j} .
\end{equation}

\begin{proposition}[Reaction diffusion: the expected gain is set at the center of the prior]\label{m:prop:prrd}
Let $\partial_tu=\Delta u+f(u)$ with $f\in C_c^\infty(\R)$ and $d\le3$, let $(e_j)$ be the real trigonometric basis, of the functions $\sqrt2\cos(2\pi k\cdot x)$ and $\sqrt2\sin(2\pi k\cdot x)$, let the prior be \eqref{m:eq:prior} with $\gamma>d$, and let $J$ be fixed.
\begin{enumerate}[label=(\roman*),nosep]
\item $\mathbf E\,I_N\le\tfrac12\sum_{j\le J}\log(1+\kappa^0_j)+o(1)$.
\item If $f$ vanishes on $[-a,a]$ for some $a>0$, then $\iota^0_j=\iota_j=(1-e^{-2\lambda_jT})/(2\lambda_jT)$, the value of \eqref{m:eq:gram} for the heat equation with $\nu=1$, and $\mathbf E\,I_N=\tfrac12\sum_{j\le J}\log(1+\kappa^0_j)+o(1)$.
\item Let $d\ge2$ and let $f$ be as in (ii), with moreover $f\ge0$, $f$ nondecreasing on $(-\infty,b]$ and strictly increasing on $(a,b)$ for some $b>a$. The first mode $e_1$, with $\lambda_1=4\pi^2$, is $\sqrt2\cos(2\pi x_p)$ or $\sqrt2\sin(2\pi x_p)$ for some coordinate $p$; let $\theta_0(x)=A\cos(2\pi x_q)$ with $q\ne p$, $a<A<b$, and $T$ so small that the solution of $w'=f(w)$, $w(0)=A$, stays at most $b$ on $[0,T]$. Then $\norm{\I_{\theta_0}e_1}^2_{L^2(\mathcal X)}>T\iota_1$, and for $J=1$
\[
\tfrac12\log\big(1+\sigma^2_{0,1}\mathcal J^1_N(\theta_0)\big)-\mathbf E\,I_N\ \to\ \tfrac12\log\frac{\norm{\I_{\theta_0}e_1}^2_{L^2(\mathcal X)}}{T\iota_1}>0 .
\]
\item Under the hypotheses of (iii), let $\psi(r)=(1-e^{-2\lambda_1r})/(2\lambda_1)$, let $\psi(T-t)f'(u_{\theta_0}(t,x))<2$ for $(t,x)\in\mathcal X$, as when $T\sup f'<2$, and let $\lambda_j=\lambda_1$ for $j\le J$. Then $\Sigma_J(0)-\Sigma_J(\theta_0)$ is positive definite, where $\Sigma_J(0)=\Id/\iota_1$ and $\Sigma_J$ is as in Corollary~\ref{m:cor:prreal}, and
\[
\begin{aligned}
&\tfrac12\log\det\big(\Id+N\sigma^2_{0,1}\Sigma_J(\theta_0)^{-1}\big)-\mathbf E\,I_N\\
&\qquad\to\ \tfrac12\log\det\big(\Sigma_J(0)\Sigma_J(\theta_0)^{-1}\big)>0 .
\end{aligned}
\]
\end{enumerate}
\end{proposition}

\begin{proof}
(i) We check the hypotheses of Proposition~\ref{m:prop:prvt}. Comparison with the spatially constant solutions shows $\norm{u_\vartheta}_\infty\le\norm\vartheta_\infty+T\sup|f|$, so the moment hypothesis of Lemma~\ref{m:lem:prexact}(ii) holds, and the mild formulation $u_\vartheta(t)=S_t\vartheta+\int_0^tS_{t-s}f(u_\vartheta(s))\,ds$, with $S_t$ the heat semigroup, together with the implicit function theorem makes $\vartheta\mapsto u_\vartheta$ continuously differentiable from $C(\Omega)$ into $C(\mathcal X)$, with derivative $\I_\vartheta h$ the solution of $\partial_tU=\Delta U+f'(u_\vartheta)U$, $U(0)=h$. The energy estimate $\norm{U(t)}_{L^2}\le e^{t\sup|f'|}\norm h_{L^2}$ bounds $\mathcal J^J_N/N$ uniformly, and the same estimate for $\I_\vartheta h-\I_0h$, whose equation has the forcing term $(f'(u_\vartheta)-f'(u_0))\I_0h$, together with $\norm{u_\vartheta-u_0}_\infty\le e^{T\sup|f'|}\norm\vartheta_\infty$, gives continuity at $\vartheta=0$. At $\vartheta=0$ the solution $u_0=v$ is constant in space, so $\I_0e_j=e^{\chi(t)-\lambda_jt}e_j$ and $\mathcal J^J_N(0)=\diag(N\iota^0_j)$, which is invertible. Now \eqref{m:eq:prvtzero} is (i).

(ii) If $f=0$ on $[-a,a]$ then $v=0$, $\chi=0$ and $\iota^0_j=\iota_j$. On the event $\mathcal E=\{\norm\theta_\infty<a\}$ comparison with the constant solutions $\pm\norm\theta_\infty$ gives $|u_\theta|<a$, hence $f(u_\theta)=0$ and $u_\theta=S_t\theta$ solves the heat equation. Let $Z_h$ be the data generated by the heat equation from the same $\theta$, design and noise, so that $Z^{(N)}=Z_h$ on $\mathcal E$. For each of $Z'=Z^{(N)}$ and $Z'=Z_h$ the chain rule gives $|I(\theta^J;Z')-I(\theta^J;Z'\mid1_{\mathcal E})|\le H(1_{\mathcal E})$, and $I(\theta^J;Z'\mid1_{\mathcal E})$ is the average of the mutual informations under the conditional priors given $\mathcal E$ and given $\mathcal E^c$, of which the first is the same for $Z^{(N)}$ and $Z_h$. By the channel bound in the proof of Lemma~\ref{m:lem:prexact}, the mutual informations given $\mathcal E^c$ are at most $\tfrac N2\log(1+M/q_N)$, where $q_N=\Pi(\mathcal E^c)$ and $M$ bounds $\int\norm{\Gcal(\vartheta)}^2_{L^\infty(\mathcal X)}\,d\Pi(\vartheta)$ for both forward maps, since the prior conditioned on $\mathcal E^c$ has density at most $1/q_N$ with respect to $\Pi$. By the Gaussian tail bound in the proof of Proposition~\ref{m:prop:prvt}, $q_N=\Pi'(\norm{\theta'}_\infty\ge a\sqrt N\delta_N)\le C\,e^{-ca^2N\delta_N^2}$, and $N\delta_N^2$ is a positive power of $N$; so $H(1_{\mathcal E})\to0$ and $q_NN\log(1+M/q_N)\to0$, and $I(\theta^J;Z^{(N)})=I(\theta^J;Z_h)+o(1)$. For the heat equation $\mathcal J^J_N(\vartheta)=\diag(N\iota_j)$ for every $\vartheta$, so \eqref{m:eq:prvt} gives $I(\theta^J;Z_h)\le\tfrac12\sum_j\log(1+\kappa_j)$ for every $N$. For the lower bound we condition on the design, which is independent of $\theta$. Given it, $Y=\Phi_J\theta^J+\xi$ with $\Phi_J=(e^{-\lambda_jt_i}e_j(x_i))_{i\le N,j\le J}$ and $\xi=\Phi_>\theta^>+\varepsilon$ independent of $\theta^J$, where $\Phi_>$ and $\theta^>$ collect the modes $j>J$, with prior covariance $\Sigma_>$. So the conditional mutual information is $\tfrac12\log\det\big(\Id+(\Sigma^J_0)^{1/2}\Phi_J^\top C_\xi^{-1}\Phi_J(\Sigma^J_0)^{1/2}\big)$ with $C_\xi=\Id+K$ the covariance of $\xi$, $K=\Phi_>\Sigma_>\Phi_>^\top$, and $C_\xi^{-1}\succeq\Id-K$ because $(\Id+K)^{-1}-(\Id-K)=(\Id+K)^{-1}K^2\succeq0$, so that $\Phi_J^\top C_\xi^{-1}\Phi_J\succeq\Phi_J^\top\Phi_J-B\Sigma_>B^\top$ with $B=\Phi_J^\top\Phi_>$. Since $|e_j|\le\sqrt2$ and the $e_j$ are orthonormal, $E(B_{jk}^2)\le2N\iota_k$, so $E\,\tr(B\Sigma_>B^\top)\le2JN\sum_{k>J}\iota_k\sigma^2_{0,k}\le2J\delta_N^{-2}\sum_k\lambda_k^{-\gamma}$, while $\Phi_J^\top\Phi_J/N$ concentrates exponentially fast at $\diag(\iota_j)$, its entries being averages of independent bounded variables. Given $\epsilon\in(0,1/2)$, outside an event of probability $O(1/(\epsilon N\delta_N^2))$ we have $\Phi_J^\top C_\xi^{-1}\Phi_J\succeq(1-2\epsilon)N\diag(\iota_j)$, and there the conditional mutual information is at least $\tfrac12\sum_j\log(1+\kappa_j)+\tfrac J2\log(1-2\epsilon)$; it is nonnegative everywhere, and $\tfrac12\sum_j\log(1+\kappa_j)=O(\log N)$. Averaging over the design and letting $\epsilon\to0$ gives $I(\theta^J;Z_h)\ge\tfrac12\sum_j\log(1+\kappa_j)-o(1)$.

(iii) The function $\theta_0$ is smooth with mean zero. Since it depends on $x_q$ only, so do $u_{\theta_0}$ and $V=f'(u_{\theta_0})$. As $f\ge0$, the constant $-A$ is a subsolution, and the solution $w$ of $w'=f(w)$, $w(0)=A$, is a solution constant in space; so $-A\le u_{\theta_0}\le w(t)\le b$, whence $V\ge0$. Moreover $V(0,\cdot)=f'(A\cos(2\pi x_q))$ is not identically zero, because $A\cos(2\pi x_q)$ takes every value in $(a,A)$ and $f$ is strictly increasing there, and by continuity neither is $V(s,\cdot)$ for small $s$. Separating variables, $\I_{\theta_0}e_1=e_1(x_p)e^{-\lambda_1t}W(t,x_q)$, where $W$ solves $\partial_tW=\partial^2_{x_q}W+VW$, $W(0)=1$, on the circle. Its Duhamel series $W=\sum_{n\ge0}W_n$, with $W_0=1$ and $W_{n+1}(t)=\int_0^tS^{(1)}_{t-s}(VW_n)(s)\,ds$, $S^{(1)}$ the heat semigroup of the circle, converges uniformly and has nonnegative terms, and $W_1(t,\cdot)>0$ for $t>0$ because the heat kernel of the circle is positive. So $W>1$ for $t>0$ and
\[
\norm{\I_{\theta_0}e_1}^2_{L^2(\mathcal X)}=\int_0^Te^{-2\lambda_1t}\int_0^1W(t,y)^2\,dy\,dt>\int_0^Te^{-2\lambda_1t}\,dt=T\iota_1 .
\]
Finally $\sigma^2_{0,1}\mathcal J^1_N(\theta_0)=\kappa_1\norm{\I_{\theta_0}e_1}^2_{L^2(\mathcal X)}/(T\iota_1)$ and $\kappa_1\to\infty$, so (ii) gives the limit.

(iv) Let $w\in\R^J$ be a unit vector and $e=\sum_{j\le J}w_je_j$, so that $\Delta e=-\lambda_1e$, $\norm e_{L^2}=1$ and $e\ne0$ almost everywhere. By Condition 2F, which Theorem~5 of Nickl~\cite{nicklbvm2024} verifies for this class, $g=(\I^*_{\theta_0}\I_{\theta_0})^{-1}e$ is smooth, and $w^\top\Sigma_J(\theta_0)w=T\langle e,g\rangle=T\norm U^2_{L^2(\mathcal X)}$, where $U=\I_{\theta_0}g$ solves $\partial_tU=\Delta U+VU$, $U(0)=g$. Let $c=1/(T\iota_1)$ and $\Psi(t,x)=\phi(t)e(x)$ with $\phi(t)=\sinh(\lambda_1(T-t))/\sinh(\lambda_1T)$. Then $\Psi(0)=e$, $\Psi(T)=0$, $\phi(t)=ce^{-\lambda_1t}\psi(T-t)$ and $\phi'-\lambda_1\phi=-ce^{-\lambda_1t}$, so $\partial_t\Psi+\Delta\Psi+V\Psi=-ce^{-\lambda_1t}(1-\psi(T-t)V)e$, and differentiating $\langle\Psi(t),U(t)\rangle_{L^2(\Omega)}$ and integrating over $[0,T]$ gives $\langle e,g\rangle=c\int_{\mathcal X}e^{-\lambda_1t}(1-\psi(T-t)V)eU$. As $\norm U^2_{L^2(\mathcal X)}=\langle e,g\rangle$, the Cauchy Schwarz inequality gives $\langle e,g\rangle\le c^2\int_{\mathcal X}e^{-2\lambda_1t}(1-\psi(T-t)V)^2e^2$. Since $0\le\psi(T-t)V<2$, with $\psi(T-t)V>0$ on a nonempty open subset of $\mathcal X$ by the proof of (iii), the integral is less than $\int_0^Te^{-2\lambda_1t}\,dt=T\iota_1$; so $\langle e,g\rangle<c^2T\iota_1=c$, that is, $w^\top\Sigma_J(\theta_0)w<1/\iota_1$. At $\vartheta=0$, $\I_0$ is the solution map of the heat equation, so $\Sigma_J(0)=\Id/\iota_1$ by \eqref{m:eq:gram}, and the limit follows from (ii), as $\sigma^2_{0,j}=\sigma^2_{0,1}$ for $j\le J$ and $N\sigma^2_{0,1}\to\infty$.
\end{proof}

The gap in Proposition~\ref{m:prop:prrd}(iii) is not small. For $d=2$, $a=1/2$, $b=5/2$, $A=1$, $T=1/80$ and $f(u)=34\,e^{2(u-1)}\tau\big(\tfrac{5u-5/2}{2}\big)\tau(6-2u)$, with $\tau(s)=e^{-1/s}/(e^{-1/s}+e^{-1/(1-s)})$ on $(0,1)$, $\tau=0$ on $(-\infty,0]$ and $\tau=1$ on $[1,\infty)$, the hypotheses of (iii) hold: $f$ vanishes outside $(1/2,3)$, is nondecreasing on $(-\infty,5/2]$ and strictly increasing on $(1/2,5/2)$, and equals $34\,e^{2(u-1)}$ on $[9/10,5/2]$, so that $w(t)=1-\tfrac12\log(1-68t)$ and $w(T)=1+\tfrac12\log(20/3)<5/2$. A numerical solution of the equations in the proof, by a Fourier method on the circle whose values are stable to six digits under refinement, gives $\norm{\I_{\theta_0}e_1}^2_{L^2(\mathcal X)}/(T\iota_1)=1.357$, a gap of $0.153$, and for the four modes with $\lambda_j=4\pi^2$, which are the first four whatever the order of ties, the same reading overstates $\mathbf E\,I_N$ by $0.611$ in the limit. Reading the second display with the efficient information $N\Sigma_J(\theta_0)^{-1}$ of Corollary~\ref{m:cor:prreal} instead does not help: in the same example $\Sigma_1(0)/\Sigma_1(\theta_0)=1.220$, a gap of $0.099$, and $0.451$ for the four modes. These efficient values are numerical: they come from the Gram matrix of $\I_{\theta_0}$ restricted to finitely many Fourier modes in $x_q$, which can only overstate the efficient information at $\theta_0$, and they are stable to five digits as the number of modes grows. That $\Sigma_1(0)>\Sigma_1(\theta_0)$, and that the gap for the four modes is positive, follow from (iv): $f'=0$ on $(-\infty,1/2]$; on $[1/2,9/10]$, $f'\le238e^{-1/5}<195$, as $\tau'(s)=\big(\delta^2+4+2(4+\delta^2)^{1/2}\big)/(2+2\cosh\delta)\le(8+3\delta^2/2)/(4+\delta^2)\le2$ with $\delta=1/s-1/(1-s)$, while $\psi\le\psi(T)<0.008$, so that $\psi(T-t)f'(u_{\theta_0})<1.6$; and where $u_{\theta_0}\ge9/10$, $f'(u_{\theta_0})\le68e^{2(w(t)-1)}=68/(1-68t)$ and $\psi(T-t)\le T-t$, so $\psi(T-t)f'(u_{\theta_0})\le68(T-t)/(3/20+68(T-t))<1$. Numerically, the bound in the proof of (iv) gives $\Sigma_1(0)/\Sigma_1(\theta_0)\ge1.199$.

\begin{remark}[Scope of Proposition~\ref{m:prop:pr}]\label{m:rem:prscope}
The proposition is stated for the heat and Stokes equations because there it is exact: the posterior is Gaussian and $\I^*\I$ is diagonal in $(e_j)$. For a nonlinear forward map in the setting of Theorem~\ref{m:thm:nickl2} and a fixed $J$ its two displays separate, as shown above: the first holds for the posterior itself with the efficient information at $\theta_0$ (Corollary~\ref{m:cor:prreal}; for the Navier Stokes equations, Proposition~\ref{m:prop:prklns}), through Theorem~\ref{m:thm:prkl}, a Kullback Leibler form of the Bernstein von Mises theorem that the Wasserstein statement of Theorem~\ref{m:thm:nickl2} does not give, the Kullback Leibler divergence not being continuous for that distance, and the second does not transfer with the information at $\theta_0$ (Propositions~\ref{m:prop:prvt} and~\ref{m:prop:prrd} and the example after the latter). Two things remain open. The first is $\mathbf E\,I_N$ up to $o(1)$ outside the class of Proposition~\ref{m:prop:prrd}(ii). The natural candidate is the efficient form at the center of the prior, $\tfrac12\log\det\big(\Id+N(\Sigma^J_0)^{1/2}\Sigma_J(0)^{-1}(\Sigma^J_0)^{1/2}\big)$ when $\I^*_0\I_0$ is invertible, which is the expected gain in the experiment of Proposition~\ref{m:prop:pr} with covariance $\Sigma_J(0)/N$. Since $N\Sigma_J(0)^{-1}\preceq\mathcal J^J_N(0)$, it is at most the right side of \eqref{m:eq:prvtzero}, with equality when $\I^*_0\I_0$ does not couple the first $J$ modes with the others, as in Proposition~\ref{m:prop:prrd}(ii), and below it by an amount that does not vanish as $N\to\infty$ otherwise; a proof would need efficient estimation uniformly over the parameters near the center of the prior. The second is the first display with $J=J_N\to\infty$, for which the proof of Corollary~\ref{m:cor:prreal} would need a rate in Theorem~\ref{m:thm:prkl} and a concentration of $NV_N$ uniform in $J$.
\end{remark}

\subsubsection{Quadratic variation, and why the diffusion coefficient is not in the statistic} Polson and Roberts~\cite{polsonroberts1994} construct Bayes factors for discretely observed diffusions with time varying drift and diffusion coefficient, through an approximation to the Bayes factor that uses only the discrete data. The fact we use in their setting is a standard one of stochastic calculus: the diffusion coefficient is determined by the quadratic variation and therefore does not enter the likelihood ratio between models with the same $\sigma$; on a fine time scale the data identify $\int_0^T\sigma^2(t,X_t)dt$ exactly in the continuous record limit, and the comparison of drifts is a Girsanov likelihood ratio conditional on it. This is the finite dimensional ancestor of two things in this part.

First, it is the same mechanism as the continuous record remark of Johannes, Polson and Stroud~\cite{jps2009} discussed in Section~\ref{m:sec:jps}: what is identified without inference is the quadratic variation, and what requires inference is the drift, that is, the part of the model that is absolutely continuous with respect to the reference measure. In our setting the analogous statement is that the noise level is identified by the residual sum of squares and the initial condition is identified by the adjoint statistic, and these are variation independent.

\begin{proposition}[Quadratic variation and the sufficient statistic]\label{m:prop:qv}
In the linear model \eqref{m:eq:design} with unknown noise variance $\varsigma^2$, the pair $\big(\Phi^\top Y,\ \norm{Y}^2\big)$ is sufficient for $(\theta,\varsigma^2)$, and the conditional law of $\theta$ given the pair and $\varsigma^2$ depends on the data only through $\Phi^\top Y$. If instead the observation noise is replaced by an observed diffusion path on $[0,T]$ with unknown diffusion coefficient, the coefficient is identified by the quadratic variation, the likelihood ratio for $\theta$ is the Girsanov density, and the sufficient statistic for $\theta$ is again the adjoint evolution driven by the data, now with the data measure replaced by $dY_t/\sigma^2(t)$.
\end{proposition}

\begin{proof}
The first claim is the factorization $\exp\{-\tfrac1{2\varsigma^2}(\norm Y^2-2\theta^\top\Phi^\top Y+\theta^\top\Phi^\top\Phi\theta)\}$. The second is Girsanov's theorem: for $dY_t=\Gcal(\theta)(t)dt+\sigma(t)dB_t$ the likelihood ratio against $\theta=0$ is $\exp\{\int_0^T\sigma^{-2}\Gcal(\theta)dY_t-\tfrac12\int_0^T\sigma^{-2}\Gcal(\theta)^2dt\}$, whose linear statistic is $\int_0^T\sigma^{-2}\Gcal(\theta)dY$, and Lemma~\ref{m:lem:adjoint} applies with $r(dt)=\sigma^{-2}(t)dY_t$.
\end{proof}

Second, the model selection problem of Polson and Roberts is the one GBC faces when the forward model itself is uncertain. Their approximate Bayes factor uses only the discrete data, and in the reading above it compares drifts at fixed quadratic variation; the GBC analog is to simulate under each candidate forward model in the step of Algorithm~\ref{m:alg:gbc} that simulates the training pairs, learn a generator for the model index alongside $\theta$, and read off the posterior model probabilities. The sufficiency analysis is unchanged within each model, and the summary for the model index is the profile of residual sums of squares, which is the finite sample version of the quadratic variation comparison.
